\documentclass[11pt,a4paper]{article}
\usepackage[margin=25mm]{geometry}
\usepackage{fontspec}
\setmainfont[BoldFont=lmroman10-bold.otf,ItalicFont=lmroman10-italic.otf,
 BoldItalicFont=lmroman10-bolditalic.otf]{lmroman10-regular.otf}
\setsansfont{lmsans10-regular.otf}
\setmonofont[Scale=0.87]{lmmono10-regular.otf}
\usepackage{amsmath,amssymb,amsthm,mathtools}
\usepackage{booktabs,longtable,array,enumitem,microtype}
\usepackage{xcolor}
\definecolor{inkblue}{HTML}{173B52}
\usepackage[colorlinks=true,linkcolor=inkblue,urlcolor=inkblue]{hyperref}
\usepackage{xurl}
\usepackage{bookmark}
\usepackage{fancyhdr}
\pagestyle{fancy}\fancyhf{}
\fancyhead[L]{\small GRH\_NIGHT: consequences and proof gates}
\fancyfoot[C]{\thepage}
\setlength{\headheight}{14pt}
\setlength{\emergencystretch}{3em}
\setlist{nosep}
\newcommand{\A}{A_D}
\newcommand{\M}{M_H}
\title{\textbf{GRH\_NIGHT}\\Consequences, methods, and the Grothendieck programme}
\author{GPT-6 Astra Pro research run\\Model effort setting not exposed}
\date{21 November 2026, 02:\allowbreak{}47 CET}
\begin{document}
\maketitle
\begin{abstract}
This finite catalogue records 147 consequences, refinements, counterexamples and explicit gates.\allowbreak{}The input $A_D$ is the Alp\"oge RH/Dirichlet-GRH theorem designated in the task. The input $M_H$ is Masullo\textquotesingle s claimed universal unitary-Hecke GRH and is always conditional. Classical proofs are located rather than reproduced; new formulations receive complete proofs. ``Proved here means established in this run, not a claim of publication novelty. The source formalization was not rebuilt. The full 940-page draft was not certified. Explicit source errors are distinguished from the conditional endpoint assumption. IN PROGRESS.
\end{abstract}
\tableofcontents
\newpage
\section*{How to read the evidence}
\addcontentsline{toc}{section}{How to read the evidence}
The six-column \path{CONSEQUENCES.tsv} is the authoritative statement and
dependency catalogue. \path{QUESTIONS.tsv} records the state of every
question treated. \path{CHECK_EVENTS.tsv} records actual check events,
including who performed them and whether a proof was read, independently
derived, checked by a finite computation, or carried from the supplied
record. A citation verified at the statement level is not a claim that the
whole classical proof was rechecked. All new proof texts in the final
checked part have a complete-reading event; computational examples are
not substituted for a general proof.

All critical-line assertions concern the stated nontrivial divisor, with
Tate shifts, trivial zeros, poles and omitted Euler factors handled as
specified. A signed Brauer factorization gives zero/pole locations; it
does not supply the nonnegative divisor orders required by Artin
holomorphy. The notation $A$, $A_D$ refers to the same adopted theorem,
and $M$, $M_H$ to the same conditional draft claim.

The original prompt asks for an open-ended exhaustive investigation.
This is a finite research record. Unresolved conjectures and unavailable
proofs are displayed as such rather than given invented resolutions.
Archive locators distinguish PDF page, repository footer and manuscript
numbering where they differ. The old papers absent from the supplied ZIP
are cited through their attached return records, without a claim of fresh
whole-proof verification.

\part{Checked findings and located consequences}

% BEGIN number_theory
\section{Number-theory consequences with the hypotheses separated}
\label{sec:number-theory}

Throughout this section $A$ denotes the user-designated Alp\"oge theorem
for the Riemann zeta function and Dirichlet L-functions.  It is an adopted
cited input, not a claim that this run rebuilt the unseen Lean dependencies.
$M$ denotes the conditional all-number-field unitary Hecke-GRH assertion
of Masullo, \S1.2.  Every use of $M$ below remains conditional.  The
accompanying catalogue is \path{work/number_theory/rows.json}; its
literature records include retrieval evidence and distinguish a checked
statement from a reread of an entire classical proof.

There are three different geometric translations.  Rational primes are
the closed points of $\operatorname{Spec}\mathbb Z$; prime ideals are the
closed points of $\operatorname{Spec}\mathcal O_K$; and a conjugacy condition
on Frobenius is a condition on a finite \'{e}tale cover of an open arithmetic
curve.  The last dictionary is the Galois-category equivalence in SGA~1,
Expos\'e~V, \S7.  The class group is
$\operatorname{Pic}(\operatorname{Spec}\mathcal O_K)$ and its class fields
are finite \'{e}tale covers.  These translations identify the objects to
which the estimates apply.  They do not turn number-field GRH into a
consequence of finite-field purity.  The inherited cote-1 trace programme
is a comparison of explicit formulas; no new archive scan is certified by
this section.  Exact archive work belongs to the geometry section.

\subsection{The classical arithmetic bounds being used}

\paragraph{NT-P01--NT-P04: rational primes and gaps.}
Schoenfeld, \emph{Sharper bounds for the Chebyshev functions
$\theta(x)$ and $\psi(x)$. II} (1976), Theorem~10, printed p.~337,
equations (6.2) and (6.3), gives, under RH,
\[
 |\psi(x)-x|<\frac{\sqrt{x}\log^2x}{8\pi}\quad(x\ge73.2),\qquad
 |\theta(x)-x|<\frac{\sqrt{x}\log^2x}{8\pi}\quad(x\ge599).
\]
Equation (6.1) replaces $\log^2x$ by $(\log x-2)\log x$ for
$x\ge23\cdot10^8$.  Corollary~1 gives
$|\pi(x)-\operatorname{li}(x)|<\sqrt{x}\log x/(8\pi)$ for
$x\ge2657$.  Here $\operatorname{li}$ has its usual normalization.
The constants in Theorem~10 were checked in the primary indexed text.
An independently retrieved primary proof locator for the same
prime-counting bound is Johnston (2022), Theorem~3.2, equation (3.6),
\emph{The Ramanujan Journal} \textbf{59}, 1307--1321.
It requires RH only to a height $T$ satisfying
$9.06\sqrt{x/\log x}/\log\log x\le T$; under $A$ one may choose
such a $T$ for each $x$.  Its \S2 explicitly says its main theorem applies
also to the ordinary, non-half-weight counting functions.  Equation (3.3)
there also improves the threshold for the displayed $\psi$ inequality
to $x\ge59$.  Thus the inaccessible original Corollary~1 page is no
longer a proof-source gate for NT-P03.

Dudek--Greni\'e--Molteni, \emph{Primes in explicit short intervals on RH}
(2016), Theorem~1.1, author PDF p.~1, gives a prime in
\[
 (x-c\sqrt{x}\log x,\ x+c\sqrt{x}\log x),
 \qquad c=\tfrac12+2/\log x,\qquad x\ge2.
\]
The same theorem gives at least $\sqrt{x}$ primes after replacing $c$ by
$c+1$.  These results imply the classical
$p_{n+1}-p_n=O(\sqrt{p_n}\log p_n)$ consequence of RH.  Neither this
bound nor the assertion of RH supplies a proof here of Cram\'er's
$O(\log^2p_n)$ conjecture, twin primes, or the binary Goldbach conjecture.
Those problems require additional information about prime correlations.

\paragraph{NT-P05--NT-P08: progressions and nonresidues.}
Greni\'e--Molteni--Perelli (2016), Theorem~1, author PDF p.~2, proves that
absolute $x_0,c_1,c_2>0$ work in
\[
 \pi(x+h;q,a)-\pi(x;q,a)\ge
 c_2\frac{h}{\varphi(q)\log x},\qquad
 c_1\varphi(q)\sqrt{x}\log x\le h\le x,\quad x\ge x_0,
\]
assuming the Dirichlet GRH modulo $q$ and $(a,q)=1$.  The following
sentence in that source gives the asymptotic if the ratio to the lower
scale tends to infinity.  We have not converted the unspecified absolute
constants into numerical constants.

Lamzouri--Li--Soundararajan (LLS), corrected
\path{arXiv:1309.3595v3}, is used for the sharp explicit conductor
statements.  Corollary~1.2, author p.~3, gives the least prime in a reduced
class modulo $q>3$ at most $(\varphi(q)\log q)^2$.
Corollary~1.1, p.~2, gives a quadratic nonresidue below $(\log q)^2$ for
prime $q\ge5$.  Theorem~1.1(1), p.~2, gives the proper-subgroup bound
recorded with its explicit functions $A(q),B(q)$ in NT-P08.
These are the already carried IA02 consequences, with the corrected
version and hypotheses rechecked; they are not new findings of this run.
Geometrically the progression is a Frobenius class in a cyclotomic cover,
and the nonresidue is detected by a finite rank-one character.

\paragraph{NT-P09: the least primitive root is a different quantifier.}
Shoup, \emph{Searching for primitive roots in finite fields} (1992),
Theorem~3, author PDF p.~4, gives under Dirichlet GRH
\[
 g(p)\ll r^4(\log r+1)^4(\log p)^2,
 \qquad r=\omega(p-1).
\]
The printed theorem was checked on the page image.  Here is the short
simplification to $g(p)\ll\log^6p$.  If $q_1<\cdots<q_r$ are the distinct
prime divisors of $p-1$, then $q_j\ge j+1$ and therefore
$p-1\ge\prod_{j=1}^r(j+1)$.  For $r\ge4$ the last at least $r/2$ factors
are at least $r/2$, so
$\log(p-1)\ge(r/2)\log(r/2)$.  Consequently
$r(\log r+1)\ll\log p$; enlarging the absolute constant handles $r<4$.
Substitution proves the assertion.  This finds a small generator for a
fixed $\mathbb F_p^\times$.  The density theorem for one fixed integer
as $p$ varies uses different, generally nonabelian, fields (NT-A01).

\subsection{A quantitative Frobenius trace transfer}

\paragraph{NT-C01 and NT-C02: the precise cited Chebotarev input.}
Let $L/K$ be a finite Galois extension, $G=\operatorname{Gal}(L/K)$,
$n_L=[L:\mathbb Q]$, and $\Delta_L=|\operatorname{disc}L|$.
For a conjugacy class $C$ put $\delta_C=|C|/|G|$.  For unramified
primes only, define
\[
 \psi_C(x)=\sum_{\substack{\mathfrak p\text{ unramified},\ m\ge1\\
                    (N\mathfrak p)^m\le x}}
 \mathbf1_C(\operatorname{Frob}_{\mathfrak p}^{m})\log N\mathfrak p,
 \qquad
 \pi_C(x)=\sum_{\substack{\mathfrak p\text{ unramified}\\N\mathfrak p\le x}}
             \mathbf1_C(\operatorname{Frob}_{\mathfrak p}).
\]
Set
\begin{align*}
 E_L(x)&=\sqrt{x}\left[
       \left(\frac{\log x}{2\pi}+2\right)\log\Delta_L+
       \left(\frac{\log^2x}{8\pi}+2\right)n_L\right],\qquad x\ge1,\\
 B_L(x)&=\sqrt{x}\left[
       \left(\frac1{2\pi}+\frac3{\log x}\right)\log\Delta_L+
       \left(\frac{\log x}{8\pi}+\frac1{4\pi}+\frac6{\log x}\right)n_L
                         \right],\qquad x\ge2,\\
 \operatorname{Li}_2(x)&=\int_2^x\frac{dt}{\log t}.
\end{align*}
Greni\'e--Molteni, \emph{An explicit Chebotarev density theorem under GRH}
(2019), Theorem~1.1 and Corollary~1.2, author PDF pp.~2--3, give
\[
 |\psi_C(x)-\delta_Cx|\le\delta_C E_L(x),\qquad
 |\pi_C(x)-\delta_C\operatorname{Li}_2(x)|\le\delta_C B_L(x),
\tag{NT-C}
\]
under RH for $\zeta_L$.  The theorem explicitly includes the unramified
version of $\psi_C$.  Its constants and ranges were checked on the page
images.  In particular Artin holomorphy is not an additional hypothesis.
$M$ supplies this input for every $L$; $A$ supplies it if $L/\mathbb Q$
is abelian by the factorization into primitive Dirichlet functions.

\paragraph{NT-F01. Statement and full proof.}
For a complex class function $f$ on $G$ put
\[
 \mu_f=\frac1{|G|}\sum_{g\in G}f(g),\qquad
 \|f\|_{1,G}=\frac1{|G|}\sum_{g\in G}|f(g)|,
\]
and define $\Psi_f$ and $\Pi_f$ by replacing $\mathbf1_C$ in the two
sums above by $f$.  Then
\[
 |\Psi_f(x)-\mu_fx|\le\|f\|_{1,G}E_L(x),\qquad
 |\Pi_f(x)-\mu_f\operatorname{Li}_2(x)|\le\|f\|_{1,G}B_L(x).
\tag{NT-F01}
\]
If $f=\chi_V$ is the character of a finite-dimensional complex
representation, then $\mu_f=\dim V^G$ and
$\|f\|_{1,G}\le\sqrt{\langle\chi_V,\chi_V\rangle_G}$.

To prove the estimates, choose one representative $g_C$ from each
conjugacy class.  The finite identities
\[
 f=\sum_C f(g_C)\mathbf1_C,
 \qquad \mu_f=\sum_C\delta_Cf(g_C)
\]
give $\Psi_f-\mu_fx=\sum_Cf(g_C)(\psi_C-\delta_Cx)$.
The triangle inequality and (NT-C) bound its modulus by
$E_L(x)\sum_C\delta_C|f(g_C)|=E_L(x)\|f\|_{1,G}$.
The same finite argument proves the assertion for $\Pi_f$.
The averaging operator $|G|^{-1}\sum_g\rho(g)$ is the projection to $V^G$,
so its trace is $\mu_{\chi_V}=\dim V^G$.  Cauchy--Schwarz for the
probability measure assigning mass $1/|G|$ to each group element yields
\[
 \left(\frac1{|G|}\sum_g|\chi_V(g)|\right)^2
 \le\frac1{|G|}\sum_g|\chi_V(g)|^2
 =\langle\chi_V,\chi_V\rangle_G.
\]
This proves every assertion.  This is a linear functional estimate on the
complexified representation ring; virtual differences need not be
nonnegative, so their error is controlled by the displayed norm.

\paragraph{NT-F02. Finite \texorpdfstring{\'{e}}{e}tale point counts, with the boundary.}
Let $U$ be a nonempty open subscheme of
$\operatorname{Spec}\mathcal O_K$, and let $X\to U$ be finite \'{e}tale
of degree $d$.  Let $S_X$ be the set of geometric points of its generic
fibre, let $L/K$ be its finite Galois splitting field, and let $r$ be the
number of $G$-orbits on $S_X$, equivalently the number of generic connected
components.  Let $S$ be the finite set of primes not in $U$ which are
unramified in $L$.  Under RH for $\zeta_L$,
\[
 \left|\sum_{\substack{\mathfrak p\in U,\ m\ge1\\(N\mathfrak p)^m\le x}}
       \#X(k_{\mathfrak p^m})\log N\mathfrak p-rx\right|
 \le rE_L(x)+d|S|\log x,\qquad x\ge1.
\tag{NT-F02}
\]
Here $k_{\mathfrak p^m}$ denotes the degree-$m$ extension of
$k_{\mathfrak p}$.  In particular, when no additional unramified prime is
removed, the error is $rE_L(x)$, independent of the degree $d$.

For the proof, finite \'{e}taleness implies that inertia at a prime of $U$
acts trivially on the generic geometric fibre.  The action of
$\operatorname{Gal}(L/K)$ on this fibre is faithful by the definition of
the splitting field, so $L/K$ is unramified at every prime of $U$.
Let $\chi(g)=\#S_X^g$ be the permutation character.  The elementary
Frobenius fixed-point formula for a finite \'{e}tale scheme over a finite
field is
$\#X(k_{\mathfrak p^m})=\chi(\operatorname{Frob}_{\mathfrak p}^{m})$:
its geometric points rational over that field are precisely the points
fixed by the $m$th Frobenius power.  Also $0\le\chi(g)\le d$.
Double counting pairs $(g,s)$ with $gs=s$ gives
\[
 \sum_{g\in G}\chi(g)=\sum_{s\in S_X}|G_s|=r|G|.
\]
Thus both $\mu_\chi$ and $\|\chi\|_{1,G}$ equal $r$.  NT-F01 bounds
the sum over every unramified prime by $rx$ with error $rE_L(x)$.
For each omitted prime, the nonnegative lost contribution is at most
\[
 d\left\lfloor\frac{\log x}{\log N\mathfrak p}\right\rfloor
          \log N\mathfrak p\le d\log x.
\]
If $N\mathfrak p>x$ the contribution is zero.  Summing this bound over
$S$ and applying the triangle inequality proves (NT-F02), including
$x=1$.  Notice that replacing the nonnegative character norm by the square
root of its second moment would give a weaker bound; the orbit count is
the relevant parameter here.

\paragraph{NT-F03. Splitting types.}
For a partition $\lambda$ of $d=|S_X|$, let $D_\lambda\subseteq G$ be
the elements with cycle partition $\lambda$ on $S_X$, and put
$\delta_\lambda=|D_\lambda|/|G|$.  This is a union of conjugacy classes.
Sum the second inequality (NT-C) over these classes and use the triangle
inequality.  Their densities sum to $\delta_\lambda$, proving
\[
 \left|\#\{\mathfrak p:\ N\mathfrak p\le x,\ \mathfrak p
 \text{ unramified},\ \operatorname{Frob}_{\mathfrak p}\in D_\lambda\}
 -\delta_\lambda\operatorname{Li}_2(x)\right|
 \le\delta_\lambda B_L(x).
\]
If $D_\lambda$ is empty both sides before the absolute value are zero.
For the canonical finite \'{e}tale model of the generic algebra over the
unramified locus, each Frobenius orbit is one closed point of its fibre,
and its orbit length is that point's residue degree.  Hence this cycle
partition is exactly its splitting type.  For the originally specified
model $X$ over $U$, restrict the counted primes to $U$; at most $|S|$
terms have been removed, so the same estimate holds with the additional
error $|S|$.  This proves both interpretations without assigning a fibre
of $X$ outside its specified base.

\paragraph{NT-F04. Avoiding a prescribed finite set.}
Let $D\subseteq G$ be a nonempty conjugacy-stable set of density
$\delta_D$, and let $S$ be any finite set of primes of $K$.  The same
finite-union argument gives
$\pi_D(x)\ge\delta_D(\operatorname{Li}_2(x)-B_L(x))$.
At most $|S|$ of the counted primes belong to $S$.  If
\[
 \delta_D\bigl(\operatorname{Li}_2(x)-B_L(x)\bigr)>|S|,
\]
the count outside $S$ is strictly positive.  It is an integer, so an
eligible prime exists.  This proves the avoidance assertion without adding
an unsupported arbitrary-exclusion clause to Bach--Sorenson's least-prime
theorem.  That theorem is NT-C03, Theorem~5.1, printed p.~1729, with
RH hypotheses specified in its \S2, p.~1720.

\subsection{Class numbers, Kummer towers, and Euclidean algorithms}

\paragraph{NT-H01--NT-H04: the regulator normalization.}
Write
\begin{align*}
 U(q)&=2e^\gamma\left(\log\log q-\log2+\tfrac12+
                                      \frac1{\log\log q}\right),\\
 V(q)&=\frac{12e^\gamma}{\pi^2}\left(\log\log q-\log2+\tfrac12+
                   \frac1{\log\log q}+\frac{14\log\log q}{\log q}\right).
\end{align*}
LLS Theorem~1.5, author pp.~3--4, gives
$V(q)^{-1}\le|L(1,\chi)|\le U(q)$ for primitive conductor
$q\ge10^{10}$.  Its Corollary~1.3 is the already carried imaginary
quadratic bound in NT-H02.  The inequality is not a two-sided estimate of
order $\log\log q$: the reciprocal lower bound must be retained.

Here is the full deduction of NT-H03, including its normalization.
For $D>0$ fundamental put $K=\mathbb Q(\sqrt D)$ and let $R_D$ be the
regulator of the full unit group modulo its torsion, so
$R_D=\log\varepsilon$ for a fundamental unit $\varepsilon>1$, regardless
of its norm's sign.  The analytic class-number formula is
\[
 \mathop{\rm Res}_{s=1}\zeta_K(s)=\frac{2h_DR_D}{\sqrt D}.
\]
An exact independent locator is Milne, \emph{Class Field Theory}, v4.03,
Chapter~VI, Corollary~2.12 and Example~2.13, printed p.~191.
Since $\zeta_K(s)=\zeta(s)L(s,\chi_D)$ and $\zeta$ has residue one,
$L(1,\chi_D)=2h_DR_D/\sqrt D>0$.
Multiplying the two inequalities in NT-H01 by $\sqrt D/2$ proves
\[
 \frac{\sqrt D}{2V(D)}\le h_DR_D\le\frac{\sqrt D}{2}U(D),
 \qquad D\ge10^{10}.
\]
This completes the short proof; no bound on $R_D$ was suppressed.

For completeness, the already inherited abelian-field formulation
NT-H04 is the following precise product calculation.  If $K/\mathbb Q$
is abelian with Dirichlet-character group $X_K$, then
\[
 h_KR_K=\frac{w_K\sqrt{\Delta_K}}{2^{r_1}(2\pi)^{r_2}}
              \prod_{\substack{\chi\in X_K\\\chi\ne1}}L(1,\chi).
\tag{NT-H04}
\]
Indeed Artin factorization gives the product for $\zeta_K/\zeta$,
and taking the residue at one gives this formula by the same class-number
formula.  Its left side and the prefactor are positive, so taking moduli
turns the product into $\prod|L(1,\chi)|$.  For any factors whose
conductors are at least $10^{10}$ insert $U$ and $1/V$; retain the finite
remaining product exactly or bound it separately.  This proves every
claimed product bound.  Dividing by an unestimated regulator does not
prove the corresponding bound on $h_K$.  This is the correction already
made in RETURN\_ZETA\_TWO IA02 and IA09, not a new result of this run.

\paragraph{NT-A01 and NT-A02: the nonabelian input in Artin's primitive-root problem.}
For a fixed integer $g\ne-1$ that is not a square, put
$K_n=\mathbb Q(\zeta_n,g^{1/n})$.  Hooley's classical theorem gives a
positive density
\[
 \delta(g)=\sum_{n\ge1}\frac{\mu(n)}{[K_n:\mathbb Q]}
\]
of primes where $g$ is a primitive root, assuming RH for these Kummer
fields.  We cite Hooley, \emph{On Artin's conjecture} (1967), the
unnumbered theorem on pp.~219--220.  The original internal page has not
been freshly obtained here; a fully inspected precise proof locator for
the density implication is Lenstra, \emph{On Artin's conjecture and
Euclid's algorithm in global fields} (1977), Theorem~(3.1), printed
p.~206, specialized to $K=F=\mathbb Q$, $W=\langle g\rangle$, $k=1$,
with density (2.15), p.~206, and positivity criterion (4.1), p.~207.
The series uses actual field degrees; replacing it by the universal
Artin constant for every $g$ loses the entanglement correction.

Lenstra's version also imposes a conjugacy-stable Frobenius condition
$C\subseteq\operatorname{Gal}(F/K)$ and asks that the residue image of
a finitely generated positive-rank $W\subseteq K^\times$ have index
dividing $k$.  Its auxiliary fields are formed from
$K(\zeta_{q(\ell)},W^{1/q(\ell)})$ with
$q(\ell)=\ell^{v_\ell(k)+1}$, taking the compositum for primes dividing
a squarefree $n$.  Theorem~(3.1) assumes RH for these fields' zeta
functions.  Theorem~(4.1) gives the finite obstruction criterion for
positive density.  $M$ supplies the hypotheses; $A$ alone does not
supply them because the Kummer fields need not be abelian over
$\mathbb Q$.  In geometric terms this is a tower of Kummer torsors of
$\mathbb G_m$, not merely the cyclotomic torsion cover.

\paragraph{NT-A03: a PID becomes Euclidean under the stated Kummer GRH.}
Lenstra's Theorem~(9.1), printed p.~218, states the result recorded in
NT-A03 for the ring of $S$-integers with $|S|\ge2$ and trivial ideal
class group.  Equation (9.2) gives its least Euclidean function.
The original page image confirms the weak inequality $|S|\ge2$;
the OCR's apparent strict inequality was not adopted.  The result is a
classical cited theorem with an exact proof locator, not a new proof of
the Euclidean algorithm.  The independent PID hypothesis and the
distinction from a norm-Euclidean algorithm remain in the catalogue.

\subsection{Growth, cancellation, and zero statistics}

\paragraph{NT-L02 and NT-L03. The fixed-character Lindel\"of deduction.}
The classical input is Simoni\v c, \emph{Estimates for L-functions in the
critical strip under GRH with effective applications} (2023), Theorem~1,
equations (4)--(6), proved in its \S3.  For a function $F$ in the Selberg
class with polynomial Euler product, zero-free to the right of the
critical line, it implies, for fixed $\delta>0$ and large $|t|$,
\[
 \log|F(\sigma+it)|,\qquad -\log|F(\sigma+it)|
 \ll_{F,\delta}(\log|t|)^{1-2\delta}\log\log|t|,
 \quad \tfrac12+\delta\le\sigma\le1.
\tag{NT-L-log}
\]
For $\delta\ge1/2$ one uses instead a smaller positive $\delta$ and
the absolutely convergent Euler product on $\sigma>1$.
The modulus bound for $\log F$ in the cited theorem gives both signs of
its real part.  Dividing the right side of (NT-L-log) by $\log|t|$
shows that it tends to zero for fixed $0<\delta<1/2$.  Hence, for every
$\epsilon>0$,
\[
 F(\sigma+it),\quad F(\sigma+it)^{-1}
 \ll_{F,\delta,\epsilon}(1+|t|)^\epsilon
 \qquad(\sigma\ge\tfrac12+\delta),
\tag{NT-L-right}
\]
with pole neighbourhoods omitted if $F$ has a pole.  The finite range of
$t$ is absorbed into the constant on any closed set avoiding poles and
zeros.  On $\sigma\ge1+\delta$ the Euler product gives the required
uniform bound directly; compact intermediate strips are already covered
by (NT-L-log) or the cited theorem's adjacent range.

Here is the passage to the critical line, often suppressed in invoking
Lindel\"of.  Fix $0<\delta<1/2$.  If $F$ has degree $d$, its functional
equation and Stirling's formula imply
\[
 |F(\tfrac12-\delta+it)|
 \ll_{F,\delta}(1+|t|)^{d\delta}
       |F^\vee(\tfrac12+\delta-it)|.
\]
The dual has the same zero-free hypothesis, so NT-L-right bounds this by
$O((1+|t|)^{d\delta+\eta})$ for every $\eta>0$.  The right boundary
$\sigma=1/2+\delta$ has exponent $\eta$.
The Phragm\'en--Lindel\"of principle on this strip gives at its midpoint
\[
 |F(\tfrac12+it)|\ll_{F,\delta,\eta}
                       (1+|t|)^{d\delta/2+\eta}.
\]
Its hypotheses hold because the standard functional equation and
finite-order analytic continuation give polynomial growth in each fixed
vertical strip; if there is a pole, apply the principle above and below
a fixed height containing it.  Choose $\delta$ and $\eta$ so that
$d\delta/2+\eta<\epsilon$.  The remaining bounded interval is compact,
and the critical line contains no pole for these L-functions.  This
proves $F(1/2+it)\ll_{F,\epsilon}(1+|t|)^\epsilon$.

For primitive Dirichlet functions the hypotheses are classical, and $A$
is the zero-free input; this proves NT-L02.  For a fixed non-pure-norm
unitary Hecke character, the Hecke L-function is entire of finite order,
its gamma shifts have nonnegative real parts, and its degree is
$[K:\mathbb Q]$.  On grouping its ideal Euler factors by rational primes,
one obtains a polynomial Euler product of that degree with local roots
of modulus at most one.  The Dirichlet coefficients are bounded by
$d_{[K:\mathbb Q]}(n)=O_\epsilon(n^\epsilon)$: each local coefficient
is bounded by the corresponding coefficient of $(1-p^{-s})^{-[K:\mathbb Q]}$.
Thus it meets the explicit axioms in Simoni\v c \S2.1, and $M$ supplies
zero-freeness.  If the character is a pure unitary norm power, its
L-function is a vertical translate of $\zeta_K$, to which the same
argument applies before translation.  Imprimitive Euler factors, when
included in the convention, form a fixed finite product; they and their
inverses are bounded on each $\sigma\ge1/2+\delta$, since their roots
have modulus at most one.  This proves NT-L03 in all the conventions just
specified.  Constants may depend on the field, character, infinity type,
and $\epsilon$; this proof makes no uniform conductor-aspect claim.

\paragraph{NT-L05 and NT-L06. Mertens, Liouville, and reciprocal tails.}
Titchmarsh, second edition, Theorem~14.25(C), printed p.~370, is the exact
classical equivalence
$\mathrm{RH}\Longleftrightarrow M(x)=O_\epsilon(x^{1/2+\epsilon})$
for every $\epsilon>0$.  We use that theorem rather than reproduce its
contour proof.  The following short deductions are included in full.
For a prime power $p^a$,
\[
 \sum_{d^2\mid p^a}\mu(p^a/d^2)=
 \begin{cases}1,&a\text{ even},\\-1,&a\text{ odd},\end{cases}
\]
because only the term with remaining exponent zero or one is nonzero.
Both sides are multiplicative, hence
$\lambda(n)=\sum_{d^2\mid n}\mu(n/d^2)$ for every $n$.
Consequently
\[
 \left|\sum_{n\le x}\lambda(n)\right|
 \le\sum_{d\le\sqrt x}|M(x/d^2)|
 \ll_\epsilon x^{1/2+\epsilon}
                \sum_{d\ge1}d^{-1-2\epsilon}
 \ll_\epsilon x^{1/2+\epsilon}.
\]
This proves Liouville cancellation, with the dependence on $\epsilon$
retained in the convergent sum.

Let $s=\sigma+it$ and choose $\eta>0$ with $1/2+\eta<\sigma$.
Partial summation for $y>x$ gives
\[
 \sum_{x<n\le y}\frac{\mu(n)}{n^s}
 =M(y)y^{-s}-M(x)x^{-s}
     +s\int_x^y M(u)u^{-s-1}\,du.
\]
The first term tends to zero as $y\to\infty$, and the integral converges
absolutely after using $|M(u)|\ll_\eta u^{1/2+\eta}$.  Its modulus is at
most a constant times
$|s|x^{1/2+\eta-\sigma}/(\sigma-1/2-\eta)$.
Together with the boundary term this proves the stated tail estimate.
It also proves local uniform convergence of the reciprocal Dirichlet
series on $\operatorname{Re}s>1/2$ by choosing $\eta$ below the distance
of a compact set from the boundary.  It agrees with $1/\zeta(s)$ on
$\operatorname{Re}s>1$ by the absolutely convergent Euler product and
therefore throughout that half-plane by analytic continuation.  At
$s=1$, $1/\zeta$ has value zero.  Hence $\sum\mu(n)/n=0$, as asserted.

\paragraph{NT-L04, NT-L07, NT-L08: what moments and zero density say.}
Harper (2013), Theorem~1, author p.~1, gives
$\int_T^{2T}|\zeta(1/2+it)|^{2k}dt\ll_kT(\log T)^{k^2}$ for fixed
$k\ge0$ under RH.  Its proof does not supply the conjectural leading
constant for every $k$.  Simoni\v c (2023), Theorem~2, equations (7),(8),
provides the explicit effective Mertens constants in NT-L07.
For NT-L08, if $\sigma>1/2$, every nontrivial zero covered by $A$ or
conditionally $M$ has real part $1/2<\sigma$, so the indicated counting
set is empty.  At $\sigma=1/2$ this reasoning says nothing about
multiplicity.  Trivial zeros from gamma factors and zeros of removed
finite Euler factors are not silently included in this definition.

\paragraph{NT-Z01 and NT-Z02: the exact pair-correlation range.}
For zeta zero ordinates $0<\gamma,\gamma'\le T$, counted with
multiplicity, put
\[
 F(\alpha,T)=\left(\frac{T}{2\pi}\log T\right)^{-1}
 \sum_{\gamma,\gamma'}T^{i\alpha(\gamma-\gamma')}
                    \frac4{4+(\gamma-\gamma')^2}.
\]
Montgomery (1973), \S1, unnumbered theorem, printed p.~181, proves under
RH the formula in NT-Z01 uniformly on each $0\le\alpha\le1-\epsilon$.
Corollary~2, p.~182, gives the proportion $2/3+o(1)$ of simple zeros.
The extension of the pair-correlation formula beyond that range is
explicitly a conjecture in that paper.  Adopting RH removes the theorem's
hypothesis; it does not replace its range by all $\alpha$.

\paragraph{NT-Z03--NT-Z05: low zeros of the prime-modulus Dirichlet family.}
For prime $q$ average over the $q-2$ nonprincipal characters, with scaled
ordinates $x_{\chi,j}=\gamma_{\chi,j}\log q/(2\pi)$.
Hughes--Rudnick, \emph{Linear statistics of low-lying zeros of
L-functions} (2003), Theorem~3.1, author p.~5, gives the unitary mean for
admissible even test functions with Fourier support in $[-2,2]$.
Theorem~4.1, p.~7, gives variance
$\int_{-1}^1|u|\widehat f(u)^2du$ with support in $[-1,1]$.
Theorem~5.1, p.~9, gives the Gaussian moments only for
$m<2/\alpha$ when the support is in $[-\alpha,\alpha]$.
These identities are unconditional with complex $\gamma$ allowed;
$A$ supplies the real-ordinate interpretation.  They are not an
unrestricted all-moment Gaussian theorem for one fixed test function.

In that source's ordering on p.~1, $\gamma_{\chi,1}$ is the first
\emph{nonnegative} ordinate, so a central zero is included with value zero.
With exactly this convention, Corollary~8.2, p.~20, gives
$\liminf_q\min_{\chi\ne\chi_0}x_{\chi,1}\le1/4$ under $A$.
Theorem~8.3, pp.~20--21, gives a positive proportion for
$x_{\chi,1}<\beta$ if $\beta\ge0.633$.
No strictly-positive-zero variant was inferred by discarding central
zeros.  The geometric comparison is with finite-field family monodromy
and its symmetry group; it does not identify the unrestricted
number-field limiting distribution.

\paragraph{NT-Z06 and NT-Z07: the additional gates.}
No argument in the cited sources proves full GUE, simplicity of every
zero, or central nonvanishing for every quadratic Dirichlet function
from zero location alone.  We record the specific missing correlation
or nonvanishing estimates, without claiming a logical independence
theorem.  For Iwaniec--Luo--Sarnak, \emph{Low lying zeros of families of
L-functions} (2000), the standing assumption on printed p.~56 explicitly
includes RH for the GL$_2$ functions in the family as well as Dirichlet
GRH.  Theorem~1.1, pp.~58--59, gives support $(-2,2)$ for fixed even
weight and squarefree level tending to infinity.  $A$ supplies its
Dirichlet part.  A universal Hecke statement does not identify a general
non-CM GL$_2$ newform with a rank-one Hecke character.  Thus the GL$_2$
hypothesis remains an additional input in this invocation.

\paragraph{NT-Z08 and NT-Z09: two exact classical equivalents.}
Rodgers--Tao, \path{arXiv:1801.05914v5}, Theorem~1, author p.~3,
proves $\Lambda\ge0$ for the classical de Bruijn--Newman constant.
The definition and RH equivalence in its \S1 give
$A\Rightarrow\Lambda\le0$, hence $\Lambda=0$.
For the deformation as normalized there, real-rootedness holds exactly
for $t\ge\Lambda$, yielding the two assertions in NT-Z08.
Lagarias, \path{arXiv:math/0008177v2}, Theorem~1.1, author p.~1,
is the exact locator for NT-Z09's divisor-sum inequality.
These are cited classical consequences, not additional new proofs of RH.

\subsection{Fixed-character cancellation and the local zero count}

\paragraph{NT-L09 and NT-L10. A full Perron transfer.}
Here is the elementary analytic form used for both rational and ideal
M\"obius sums.  Suppose
$G(s)=\sum_{n\ge1}a_nn^{-s}$ is absolutely convergent for
$\operatorname{Re}s>1$, $|a_n|\le d_D(n)$ for a fixed positive integer
$D$, and $G$ extends holomorphically to $\operatorname{Re}s>1/2$.
Assume that for every $\delta,\eta>0$,
\[
 G(\sigma+it)\ll_{\delta,\eta}(1+|t|)^\eta
 \quad\text{uniformly for }\tfrac12+\delta\le\sigma\le2.
\tag{NT-Perron-input}
\]
Then $\sum_{n\le x}a_n\ll_\epsilon x^{1/2+\epsilon}$ for every
$\epsilon>0$.

We prove the assertion, recording the truncation error as well as the
contour displacement.  It suffices to take $0<\epsilon<1/4$ and
half-integral $x\ge4$, because replacing any real $y$ by
$\lfloor y\rfloor+1/2$ leaves the partial sum unchanged and changes its
size by a bounded factor; the bounded remaining range is absorbed into
the constant.  Put
$c=1+1/\log x$, $T=x^2$, $\delta=\epsilon/4$ and
$\eta=\epsilon/8$.  The classical truncated Perron kernel gives
\[
 \sum_{n\le x}a_n=
 \frac1{2\pi i}\int_{c-iT}^{c+iT}G(s)\frac{x^s}{s}\,ds+R,
\]
with
\[
 |R|\ll\sum_{n\ge1}|a_n|(x/n)^c
           \min\left(1,\frac1{T|\log(x/n)|}\right).
\tag{NT-Perron-error}
\]
This is the usual Perron formula; an exact classical locator is
Titchmarsh, Lemma~3.12 and its proof, printed pp.~61--62.  The kernel
estimate follows by moving the integral of $y^s/s$ to the left or right
according to $y>1$ or $y<1$, taking the residue at zero in the first
case, and integrating the horizontal tails; the alternative trivial
bound handles $|\log y|\le T^{-1}$.

For $n\le x/2$ or $n\ge2x$ the denominator $|\log(x/n)|$ is bounded
below by $\log2$, so that part of (NT-Perron-error) is
\[
 \ll x^cT^{-1}\sum_{n\ge1}d_D(n)n^{-c}
 =x^cT^{-1}\zeta(c)^D\ll xT^{-1}(\log x)^D.
\]
For $x/2<n<2x$, use $d_D(n)\ll_{D,\eta}x^\eta$ and
$|\log(x/n)|\gg |x-n|/x$.  Since $x$ is half-integral,
$\sum_{x/2<n<2x}|x-n|^{-1}\ll\log x$.  Also $(x/n)^c\ll1$,
uniformly for these $x,n$.  This part of the error is therefore
$O(x^{1+\eta}T^{-1}\log x)$.  Both parts are smaller than
$x^{1/2+\epsilon}$.

Move the vertical line to $\sigma_0=1/2+\delta$.  The integrand has no
pole in the intervening rectangle: $G$ is holomorphic there and
$\sigma_0>0$.  On either horizontal segment,
(NT-Perron-input) bounds the integral by
$O(x^cT^{\eta-1})=O(x^{-1+2\eta})$.
On the new vertical segment it gives
\[
 \ll x^{\sigma_0}\left(1+\int_1^Tt^{\eta-1}\,dt\right)
 \ll_\eta x^{1/2+\delta}T^\eta
 =O(x^{1/2+\epsilon/2}).
\]
The bounded part $|t|\le1$ is harmless because $\sigma_0>0$ and $G$
is bounded on its compact segment.  This proves the asserted partial-sum
bound with room in the exponent.

Apply this lemma first to $G(s)=1/L(s,\chi)$.
Its coefficients are $a_n=\mu(n)\chi(n)$, including the usual zero
values at integers not coprime to the conductor.  Under $A$, the
reciprocal is holomorphic on $\operatorname{Re}s>1/2$; the principal
L-function's pole produces a zero of the reciprocal, not a pole.
NT-L02 supplies (NT-Perron-input), so the twisted bound follows.
Character orthogonality, for $(a,q)=1$, is the finite identity
\[
 \sum_{\substack{n\le x\\n\equiv a\ ({\rm mod}\ q)}}\mu(n)
 =\frac1{\varphi(q)}\sum_{\chi\ ({\rm mod}\ q)}
       \overline{\chi(a)}\sum_{n\le x}\mu(n)\chi(n).
\]
Imprimitive functions differ by finitely many Euler factors whose zeros
lie on $\operatorname{Re}s=0$, and their reciprocals are bounded in any
closed half-strip to the right of $1/2$.  Thus each term satisfies the
same fixed-character argument.  A finite sum proves the fixed-$q$
progression bound, with no uniformity in $q$ claimed.

For a unitary Hecke character, expand
\[
 \frac1{L(s,\eta)}
 =\sum_{(\mathfrak a,\mathfrak f)=1}
         \mu_K(\mathfrak a)\eta(\mathfrak a)(N\mathfrak a)^{-s}
 =\sum_{n\ge1}a_nn^{-s}.
\]
For each $n$, $|a_n|$ is at most the number of ideals of norm $n$.
If $[K:\mathbb Q]=D$, comparison of the rational-prime Euler factors
of $\zeta_K$ with $(1-p^{-s})^{-D}$ gives this number at most $d_D(n)$.
Indeed there are at most $D$ prime ideals above $p$, each of residue
degree at least one, and replacing every local factor by a degree-one
factor can only increase its nonnegative coefficients.  Holomorphy and
(NT-Perron-input) follow conditionally on $M$ from NT-L03, including
pure norm twists.  The lemma proves NT-L10.  Averaging over the finite
character group of any fixed ray class group then projects onto one ray
class exactly as above; all sums are over ideals prime to its modulus.
This completes every claimed cancellation statement.

\paragraph{NT-Z10. The fixed infinite-order Hecke zero count.}
Let $\eta$ be any fixed unitary Hecke character of a number field $K$.
Its number of nontrivial zeros, counted with multiplicity, in
\[
 0\le\operatorname{Re}s\le1,\qquad |\operatorname{Im}s-T|\le1
\]
is $O_{K,\eta}(\log(|T|+3))$, unconditionally.  This assertion includes
infinite-order characters; its constant is not asserted uniform in
their infinity types.

Here is a full local argument.  On any fixed vertical strip the finite
Hecke L-function has a polynomial upper bound in $|t|$, away from its
possible pole.  In the present argument this follows directly from the
unconditional first part of Simoni\v c's Theorem~1, equation (2), on
$\sigma\ge1/2$, and from the functional equation and Stirling's formula
on the reflected strip.  The Selberg-class hypotheses for a non-pure-norm
character were verified in NT-L03 without using GRH; a pure norm twist
is a translate of the Dedekind function.  Consequently constants $C,A>0$
exist such that, for large $|T|$,
\[
 \max_{|s-(2+iT)|\le4}|L(s,\eta)|\le C(|T|+3)^A.
\tag{NT-J-upper}
\]
For a pure norm twist the only pole has a fixed imaginary coordinate, so
choose $|T|$ sufficiently large that this disk does not contain it.
In all other cases the function is entire.

The Euler product at the centre gives a lower bound that a general
Schwartz Mellin transform need not have:
\[
 |L(2+iT,\eta)|
 \ge\prod_{\mathfrak p\nmid\mathfrak f}(1+(N\mathfrak p)^{-2})^{-1}
 \ge\prod_{\mathfrak p}(1-(N\mathfrak p)^{-2})
 =\zeta_K(2)^{-1}>0.
\tag{NT-J-lower}
\]
These comparisons follow term by term from $|\eta(\mathfrak p)|=1$
and $1/(1+u)\ge1-u$ for $0\le u<1$, and persist on taking limits of
the absolutely convergent products.

Let $n_T(3)$ count zeros in the disk of radius $3$ about $2+iT$.
Jensen's formula on the radius-$4$ disk, or its limiting-radius version
if a zero lies on the outer circle, implies
\[
 n_T(3)\log(4/3)
 \le\log\max_{|s-(2+iT)|\le4}|L(s,\eta)|
                -\log|L(2+iT,\eta)|.
\]
The right side is $O_{K,\eta}(\log(|T|+3))$ by (NT-J-upper) and
(NT-J-lower).  Every point with $0\le\operatorname{Re}s\le1$ and
$|\operatorname{Im}s-T|\le1$ is at distance at most $\sqrt5<3$
from the centre, so its zeros are included in $n_T(3)$.  The remaining
bounded interval of $T$ meets only finitely many zeros, by meromorphic
analyticity and the absence of an identically zero function; increasing
the constant covers it.  This proves the claim.

The finite-order conductor-explicit counterpart is
Greni\'e--Molteni2019, Appendix~A, Theorem~A.1, author p.~26.
NT-Z10 instead proves the fixed-character form needed for the local
arithmetic divisor, and supplies the infinite-order version of root
MT-M04.  The argument uses the lower bound at its Jensen centre and
polynomial vertical growth.  It does not repair a statement quantified
over every function in an arbitrary Schwartz Mellin-transform space.

\subsection{An explicit uniform cyclotomic specialization}

\paragraph{NT-P11 and NT-P12.}
For $q\ge3$ and $(a,q)=1$, define
\[
 \psi(x;q,a)=\sum_{\substack{p^m\le x\\p^m\equiv a\ ({\rm mod}\ q)}}\log p,
 \qquad
 \pi(x;q,a)=\#\{p\le x:p\equiv a\ ({\rm mod}\ q)\}.
\]
Put $e_q=1$ if $q\equiv2\pmod4$ and $e_q=0$ otherwise.
Then $A$ implies
\begin{align*}
 \left|\psi(x;q,a)-\frac{x}{\varphi(q)}\right|
 &\le\sqrt{x}\left[
 \left(\frac{\log x}{2\pi}+2\right)\log q+
 \frac{\log^2x}{8\pi}+2\right]+e_q\log x,\qquad x\ge1,\\
 \left|\pi(x;q,a)-\frac{\operatorname{Li}_2(x)}{\varphi(q)}\right|
 &\le\sqrt{x}\left[
 \left(\frac1{2\pi}+\frac3{\log x}\right)\log q+
 \frac{\log x}{8\pi}+\frac1{4\pi}+\frac6{\log x}\right]+e_q,
 \qquad x\ge2.
\end{align*}

For the full proof, take $L=\mathbb Q(\zeta_q)$ and $K=\mathbb Q$.
The Galois group is $(\mathbb Z/q\mathbb Z)^\times$, of order
$n_L=\varphi(q)$, and its conjugacy classes are singletons.
The cyclotomic discriminant formula gives
\[
 \Delta_L=q^{\varphi(q)}
    \prod_{p\mid q}p^{-\varphi(q)/(p-1)},\qquad
 \log\Delta_L\le\varphi(q)\log q.
\]
At every prime $p\nmid q$, arithmetic Frobenius acts on $\zeta_q$ by
$\zeta_q\mapsto\zeta_q^p$.  Thus its $m$th power is the element
represented by $p^m$ modulo $q$.  The use of arithmetic or geometric
Frobenius merely replaces the labelled class by its inverse; the estimate
is identical after that relabelling.  The Dirichlet factorization of
$\zeta_L$ makes RH for it a consequence of $A$.
Apply NT-C01 and NT-C02 with $\delta_C=1/\varphi(q)$ and insert the
discriminant upper bound; all displayed coefficients multiplying
$\log\Delta_L$ are positive on the stated ranges.

If $q$ is odd or $4\mid q$, the primes unramified in $L$ are exactly
those not dividing $q$, so the resulting sums are the progression sums
just defined.  If $q\equiv2\pmod4$, write $q=2m$ with $m$ odd.
Then $L=\mathbb Q(\zeta_m)$ and the prime $2$ is unramified, whereas it
is excluded from a reduced class modulo $q$.  For odd primes the class
condition modulo $m$ is equivalent to that modulo $q$, since reduction
$(\mathbb Z/q\mathbb Z)^\times\to(\mathbb Z/m\mathbb Z)^\times$ is
an isomorphism.  The unramified Chebotarev sum can therefore differ only
by the powers of $2$ whose Frobenius powers lie in the designated class.
Their total weight is at most
$\lfloor\log x/\log2\rfloor\log2\le\log x$.
For the prime-counting sum the discrepancy is at most one.  Applying the
triangle inequality proves both formulas, including this degeneration.

The small moduli are transparent separately:
\begin{align*}
 \psi(x;1,0)&=\psi(x),& \pi(x;1,0)&=\pi(x),\\
 \psi(x;2,1)&=\psi(x)-\lfloor\log x/\log2\rfloor\log2,&
 \pi(x;2,1)&=\pi(x)-\mathbf1_{x\ge2}.
\end{align*}
For $1\le x<2$ the floor is zero.  Hence NT-P01 and NT-P03 give the
corresponding bounds in their stated ranges, with the explicit powers-of-two
correction retained for modulus two.  These are estimates on the open
cyclotomic arithmetic curve, and the additional term is exactly its
removed-prime contribution.

\subsection{A specified quadratic family and central nonvanishing}

\paragraph{NT-Z11. The exact one-level-density input.}
Let $\mathcal D$ consist of the nonzero odd squarefree integers, of both
signs.  For $d\in\mathcal D$, $8d$ is a fundamental discriminant and
$\chi_{8d}$ is its primitive quadratic character.  Put
$\ell_X=\log(X/(2\pi e))$.  Let $w$ be an even, nonnegative, nonzero
Schwartz function and write
\[
 W_w(X)=\sum_{d\in\mathcal D}w(d/X),\qquad
 \mathcal Z_{w,\phi}(X)=\frac1{W_w(X)}
  \sum_{d\in\mathcal D}w(d/X)
       \sum_{\rho=1/2+i\gamma}\phi(\gamma\ell_X/(2\pi)),
\]
where the inner sum counts every nontrivial zero with multiplicity.
Use the Fourier convention
$\widehat\phi(u)=\int_{\mathbb R}\phi(t)e^{-2\pi iut}\,dt$.
For a real even Schwartz function $\phi$ with
$\operatorname{supp}\widehat\phi\subset(-2,2)$, $A$ implies
\[
 \lim_{X\to\infty}\mathcal Z_{w,\phi}(X)
 =\widehat\phi(0)-\frac12\int_{-1}^1\widehat\phi(u)\,du.
\tag{NT-Z11}
\]
This is the leading term of Fiorilli--Parks--S\"odergren,
\path{arXiv:1601.06833}, Theorem~1.1, author p.~2, with definitions
(1.1)--(1.2) on the same page.  Its proof uses Theorem~3.5, author p.~12, and is concluded on p.~14.
The GRH input in the prime-sum step is Dirichlet GRH
(Lemma~2.5, author p.~7), which $A$ supplies.  No unrestricted Fourier
support, individual central nonvanishing, or nonabelian GRH is being
asserted here.

\paragraph{NT-Z12. A full nonvanishing deduction, including deweighting.}
Under $A$, at least $15/16$ of this specified family have a nonzero
central value in the following precise sense:
\[
 \liminf_{X\to\infty}
 \frac{\#\{d\in\mathcal D:|d|\le X,\ L(1/2,\chi_{8d})\ne0\}}
      {\#\{d\in\mathcal D:|d|\le X\}}\ge\frac{15}{16}.
\tag{NT-Z12}
\]
The same lower bound holds for the weighted proportions defined by any
of the fixed weights $w$ in NT-Z11.  This is a deduction from the cited
classical density theorem, not a claim of a new nonvanishing theorem.
Here are all approximation and counting details.

Let $m_d$ be the order of vanishing at $1/2$.  The completed primitive
quadratic function has root number $+1$.  One way to see the sign without
choosing a Gauss-sum square root is to divide the symmetric completed
Dedekind functional equation of $\mathbb Q(\sqrt{8d})$ by that of
$\mathbb Q$: the factorization
$\zeta_{\mathbb Q(\sqrt{8d})}=\zeta L(s,\chi_{8d})$ and the real or
complex gamma factors identify the quotient, up to a nonzero constant,
with the completed quadratic function.  The gamma duplication formula
handles the imaginary quadratic case.  Thus this entire function is
even about $s=1/2$.  Its gamma factor is finite and nonzero at the
centre, so $m_d$ is an even nonnegative integer.  In particular
\[
 2\mathbf1_{L(1/2,\chi_{8d})=0}\le m_d.
\tag{NT-central-parity}
\]
These functional equations are the standard Hecke equations; a precise
source for the completed convention is Ramakrishnan--Valenza,
Theorem~7-19, printed p.~279.

Choose real even smooth functions $g_j$ compactly supported in $(-1,1)$
such that $\int g_j=1$ and
$g_j\to g=\tfrac12\mathbf1_{[-1,1]}$ in $L^2(\mathbb R)$.
For example, smooth even cutoffs of $[-1+1/j,1-1/j]$, normalized by
their integrals, have these properties.  Define
\[
 f_j(t)=\int g_j(u)e^{2\pi iut}\,du,\qquad
 \phi_j(t)=|f_j(t)|^2.
\]
Then $\phi_j$ is real, even, nonnegative and Schwartz,
$\phi_j(0)=1$, and
$\operatorname{supp}\widehat\phi_j\subset(-2,2)$.
The support assertion follows because
$\widehat\phi_j=g_j*\widetilde g_j$, where
$\widetilde g_j(u)=\overline{g_j(-u)}$.
Plancherel gives $f_j\to f$ in $L^2$, with
\[
 f(t)=\frac{\sin(2\pi t)}{2\pi t}.
\]
Moreover, Cauchy--Schwarz yields
\[
 \|\phi_j-|f|^2\|_1
 \le\|f_j-f\|_2(\|f_j\|_2+\|f\|_2)\longrightarrow0.
\]
Consequently their Fourier transforms converge uniformly.  Since
\[
 \widehat{|f|^2}(u)=
 \begin{cases}\tfrac12-\tfrac14|u|,&|u|\le2,\\0,&|u|>2,
 \end{cases}
\]
we obtain
\[
 \lim_{j\to\infty}
 \left(\widehat\phi_j(0)-\frac12
                   \int_{-1}^1\widehat\phi_j(u)\,du\right)
 =\frac12-\frac12\cdot\frac34=\frac18.
\tag{NT-test-limit}
\]
Under $A$ every ordinate in NT-Z11 is real.  The nonnegativity of
$\phi_j$ and $\phi_j(0)=1$ therefore imply
\[
 \frac1{W_w(X)}\sum_{d\in\mathcal D}w(d/X)m_d
 \le\mathcal Z_{w,\phi_j}(X).
\]
First take $\limsup_{X\to\infty}$ for a fixed $j$, use NT-Z11, and
then let $j\to\infty$.  The result and (NT-central-parity) give
\[
 \limsup_{X\to\infty}
 \frac{\sum_{d\in\mathcal D}w(d/X)
               \mathbf1_{L(1/2,\chi_{8d})=0}}{W_w(X)}\le\frac1{16}.
\tag{NT-weighted-vanishing}
\]
This proves the weighted assertion without applying the density theorem
directly to a nonsmooth Fourier transform or to support touching $2$.

For completeness, the ordinary count follows by a separate elementary
argument.  If $A_+(Y)$ counts positive odd squarefree integers up to $Y$,
the squarefree indicator gives
\begin{align*}
 A_+(Y)&=\sum_{\substack{k\le\sqrt Y\\k\ {\rm odd}}}\mu(k)
       \#\{n\le Y/k^2:n\ {\rm odd}\}\\
 &=\frac Y2\sum_{k\ {\rm odd}}\frac{\mu(k)}{k^2}+O(\sqrt Y)
 =\frac4{\pi^2}Y+O(\sqrt Y).
\end{align*}
Here the omitted tail is $O(Y\sum_{k>\sqrt Y}k^{-2})=O(\sqrt Y)$,
and the Euler product is
$\sum_{k\ {\rm odd}}\mu(k)/k^2=8/\pi^2$.
Thus $N(X):=\#\{d\in\mathcal D:|d|\le X\}
 =8X/\pi^2+O(\sqrt X)$.
Partial summation of this counting formula gives, for every fixed
smooth compactly supported $w$,
\[
 W_w(X)=\frac4{\pi^2}X\int_{\mathbb R}w(t)\,dt+O_w(\sqrt X).
\]
Indeed apply Stieltjes integration to the positive count $A_+$ and
its reflected negative count; the integrated error is bounded by a
constant depending on $\int_0^\infty\sqrt t\,|w'(t)|dt$.

Fix $\delta>0$.  Take such an even smooth $w$ equal to $1$ on
$[-1,1]$, between $0$ and $1$ everywhere, and supported in
$[-1-\delta,1+\delta]$.  The number of vanishing members with
$|d|\le X$ is bounded above by its weighted counterpart.  Dividing by
$N(X)$ and using (NT-weighted-vanishing) gives a $\limsup$ of at most
$\frac1{16}\cdot\frac12\int w\le(1+\delta)/16$.
Letting $\delta\downarrow0$ proves (NT-Z12).

The characters in this statement are quadratic rank-one Artin objects.
Their varying conductors describe a family of double covers of open
arithmetic curves.  The symplectic expression in NT-Z11 matches the
finite-field hyperelliptic-family comparison discussed immediately after
Theorem~1.1 of the cited paper.  This comparison does not enlarge the
analytic support range or turn the $15/16$ bound into nonvanishing for
each member.  The ordinary proportion proved here concerns the
$\chi_{8d}$ family exactly; it has not been silently relabelled as all
fundamental discriminants.

% END number_theory

% BEGIN geometry_programme
\section{Geometry, motives, and the Grothendieck programme}
\label{sec:geometry}

The letters A and M have the dependency meaning fixed in the introduction.
In particular, accepting the stipulated endpoint A is not a certification of
every intermediate statement in the note, and assuming M does not certify
the functional analysis in its proposed proof. The findings below are new
formulations proved in this run; no claim of priority over the literature is
intended. Earlier findings are used with their original identifiers.

\subsection{The transfer results already established}
\label{ag:transfer}

\paragraph{AG-AT (the exact Artin--Tate domain).}
For a finite-type scheme over $\mathbb Z$ with exact all-prime point counts
\[
 \#X(\mathbb F_{p^r})=
 \sum_j p^{jr}\operatorname{Tr}(F_p^r\mid V_j^{I_p}),
\]
where the $V_j$ are finitely many virtual finite-image representations,
the identity
\[
 \zeta_X(s)=\prod_j L(V_j,s-j)
\]
is unconditional. This is ZETA\_TWO ZT01, in
\path{RETURN_ZETA_TWO_final_IMPLICATIONS.md}, Section 1,
extending BRIDGES F1. If the representations have abelian image, A gives
noninteger global divisor support on the lines $\Re s=j+1/2$.
For arbitrary finite image, M gives the corresponding meromorphic
divisor-support assertion through integer Brauer induction (ZT02).
Negative exponents can produce poles; Artin holomorphy is an extra question.

\paragraph{AG-DED and AG-LOCAL (ramification and models).}
The permutation representation on the embeddings of a number field has
inertia invariants giving the Euler factors of $\operatorname{Spec}\mathcal
O_K$ at ramified primes as well. Thus every Dedekind zeta function is covered
by M, and abelian fields over $\mathbb Q$ are covered by A (ZT03).
If a factorization holds only away from a finite set $S$, the actual factors
at $S$ must be restored (ZT04). The counterexample ZT05,
\[
 \zeta_{\operatorname{Spec}\mathbb Z[1/p]}(s)
       =(1-p^{-s})\zeta(s),
\]
has additional zeros $2\pi i k/\log p$, $k\ne0$. Such local divisor lines
are not shifted global critical lines. This is a statement about the
integral model, not a failure of the generic-fibre motivic normalization.

\paragraph{AG-CM (the potentially CM domain).}
ZETA\_TWO CM01 supplies the all-weight extension to pure objects in the
absolute-Hodge category generated by Artin motives, Tate motives, and
potentially CM abelian varieties. Its cited CM descent input gives a signed
product of Hecke functions of one weight; norm twisting makes these
characters unitary. The required M therefore includes infinite-order
unitary characters. The result is meromorphic divisor RH, and does not
declare absolute-Hodge projectors algebraic. The exact prior locator is
\path{RETURN_ZETA_TWO_final_IMPLICATIONS.md}, Section 2, CM01,
with its reported Anderson Section 5.7, Corollary 5.7.2 and Remark 5.7.3.
The original source has now been recovered and checked: G.~W.~Anderson,
\emph{Cyclotomy and an extension of the Taniyama group}, Compositio
Math.~57 (1986), 153--217. The category is defined in Section 5.3,
printed p.~178 (PDF p.~27); Theorem 6, Section 5.6, printed pp.~179--180
(PDF pp.~28--29), identifies its one-dimensional characters with algebraic
Hecke characters, including the inertia and Frobenius comparisons.
Corollary 5.7.2 and Remark 5.7.3 are both on printed p.~181 (PDF p.~30).
The former gives integer exponents, with no positivity claim; the latter
gives the common weight. Its internal reference to ``Corollary 5.8.2'' is
a printed cross-reference error, not a new result used here.

\paragraph{AG-CM-REC and AG-CM-PERIOD (classical conclusions already
stronger than zero location).}
In this precise potentially CM absolute-Hodge category, equality of the
whole $L$-functions implies isomorphism of objects: Anderson,
Proposition 5.7.5, printed pp.~181--182 (PDF pp.~30--31). This is
unconditional. The proof uses semisimplicity and Corollary 5.7.4's
identification, for a weight-zero object, of the pole order at $1$ with
the dimension of its invariant part. Equality of the local factors makes
the tensor products with a dual agree away from finitely many primes;
their local factors are nonzero and nonpolar at $1$. A pole of
$L(M\otimes M^\vee,s)$ therefore gives a nonzero morphism
$M\longrightarrow N$, and induction on the common rank proves the
claim. The proof was read in full, rather than inferred from GRH.
This is a recognition theorem for semisimple motives; it supplies no
effective finite bound for general infinite CM monodromy.

Anderson's Section 5.8, Theorem 7, printed p.~182 (PDF p.~31), also
proves Deligne's period conclusion
$L(M,0)/c^+(M)\in\mathbb Q$ for the stated \emph{tractible, critical}
objects: those for which the $L$-function is a single algebraic Hecke
$L$-function over a totally real or CM field. This is a separate
classical special-value theorem, using Siegel and Blasius, and is not
enlarged to arbitrary motives by A or M. It complements, rather than
replaces, the cycle/height and integral comparison gates below.

\paragraph{AG-WEIGHT and AG-NONCM (the residual geometric factors).}
With geometric Frobenius,
\[
 \operatorname{wt}N(n)=\operatorname{wt}N-2n,
 \qquad L(N(n),s)=L(N,s+n).
\]
The line for weight $w$ is $\Re s=(w+1)/2$. Positive weight is not the
remaining obstruction: CM objects of positive weight are covered by M,
whereas $\operatorname{Sym}^2h^1(E)(1)$ for a non-CM elliptic curve still
has a non-CM factor despite having weight zero. ZT06 reduces a finite-type
scheme to finitely many smooth projective geometric factors and actual
exceptional fibres. ZT07 proves that, under A, an off-line elliptic zero
in $1/2<\Re s<3/2$ cannot be cancelled by the two Tate factors in the
good-model formula
\[
 \zeta_{\mathcal E}(s)
   =\zeta^S(s)\zeta^S(s-1)/L^S(E,s).
\]
These are retained reductions, not a proof of RH for the non-CM factor.
Framework locators carried from ZETA\_TWO are SGA 4 XVII Sections 5.1--5.3;
SGA 5 XV Section 3, no.~2, unnumbered theorem, equation (1);
EGA IV$_3$ (8.8.2); and Deligne, Weil I, Theorem 1.6, and Weil II,
Theorem 3.3.1.

\subsection{AG-FROB: a finite recognition theorem for Artin--Tate objects}
\label{ag:frob}

Let $L/\mathbb Q$ be a nontrivial finite Galois extension, with group $G$,
absolute discriminant $D_L$, and degree $n_L$. Put
\[
 B_L=(4\log D_L+\tfrac52 n_L+5)^2.
\]
Let
\[
 U=\bigoplus_j V_j(-j),\qquad W=\bigoplus_j W_j(-j),
\]
where all sums are finite and $V_j,W_j$ are actual finite-dimensional
characteristic-zero representations of $G$. Their Euler polynomials at
unramified primes are understood as those of the associated compatible
Artin--Tate systems; a single auxiliary prime may instead be chosen greater
than $B_L$.

\textbf{Statement.} Assume RH for $\zeta_L$ and $\zeta_{\mathbb Q}$.
If
\[
 \det(1-TF_p\mid U)=\det(1-TF_p\mid W)
 \quad\text{for every }p\le B_L\text{ unramified in }L,
\]
then $V_j\simeq W_j$ for every $j$, and hence $U\simeq W$.
For rational coefficients this is also an isomorphism in the rational
Artin--Tate category. A supplies the hypothesis when $L/\mathbb Q$ is
abelian; M supplies it in general. When $L=\mathbb Q$, the one prime $p=2$
suffices.

\begin{proof}
At an unramified prime, every eigenvalue in $V_j(-j)$ has the form
$p^j\omega$, where $\omega$ is a root of unity. Therefore its complex
absolute value is $p^j$. Distinct integers $j$ give distinct absolute
values, so equality of the total Euler polynomials separates uniquely into
equality of the polynomials in each weight. It follows that
$\operatorname{Tr}(F_p\mid V_j)=\operatorname{Tr}(F_p\mid W_j)$ for every
$j$ and every prime in the hypothesis.

Bach--Sorenson, \emph{Explicit bounds for primes in residue classes},
Math. Comp. 65 (1996), Theorem 5.1, printed p.~1729 (PDF p.~13), gives
an unramified prime $p\le B_L$ in every conjugacy class of $G$.
The precise analytic assumptions are stated in their Section 2,
printed p.~1720: zero-freeness of $\zeta_L$ and $\zeta_{\mathbb Q}$ in
$\Re s>1/2$. The reduction from cyclic to general Galois extensions is
Corollary 3.3, p.~1721. Thus the characters of $V_j$ and $W_j$ agree on
every conjugacy class. In characteristic zero finite-group representations
are semisimple, and equal characters give isomorphic representations.

For completeness, this last assertion also descends to a non-algebraically
closed characteristic-zero coefficient field $E$. The space of
$G$-equivariant maps is the solution space over $E$ of finitely many linear
equations. If it contains an invertible map after scalar extension, its
determinant polynomial is nonzero on that $E$-vector space. Since $E$ is
infinite it has a nonzero value at an $E$-point. That point is an invertible
$G$-equivariant map over $E$. For $L=\mathbb Q$, the eigenvalues $2^j$
already recover all multiplicities from the polynomial at $2$.
\end{proof}

In Grothendieck's language this is effective recognition inside a specified
semisimple Artin--Tate realization category. It extends
ALPOGE\_BRIDGES AQ-FINITE, whose sharper abelian conductor bound remains
available. It does not construct algebraic projectors for an arbitrary
geometric realization. Moreover, recognition of an Artin motive is weaker
than recognition of a finite cover as a scheme; the connected Gassmann
example in CX-GEO-004's companion finding makes this distinction explicit.
The forbidden-prime variant is handled there by a quadratic compositum.

\paragraph{AG-COVER: equal Artin motives need not give equal covers.}
Put $L=\mathbb Q(\sqrt2,\sqrt3)$ and let $K_1,K_2,K_3$ be its three
quadratic subfields. The finite schemes
\[
 X=\operatorname{Spec}L\amalg\operatorname{Spec}\mathbb Q
       \amalg\operatorname{Spec}\mathbb Q,
 \qquad
 Y=\coprod_{i=1}^3\operatorname{Spec}K_i
\]
are not isomorphic, but their rational Artin motives, and all their local
Artin Euler factors, are equal. Consequently
\[
 \zeta_L(s)\zeta_{\mathbb Q}(s)^2
       =\zeta_{K_1}(s)\zeta_{K_2}(s)\zeta_{K_3}(s).
\]
This identity includes the factors at ramified primes.

\begin{proof}
Write $G=\operatorname{Gal}(L/\mathbb Q)=C_2\times C_2$, with its three
subgroups $H_i$ of order two. The character of
$\mathbb Q[G]\oplus\mathbb Q^2$ takes value $6$ at the identity and
$2$ at every other element. Each nonidentity element belongs to exactly
one $H_i$, so the character of $\bigoplus_i\mathbb Q[G/H_i]$ has the
same values. Characteristic-zero semisimplicity and the descent argument
in AG-FROB give an isomorphism of rational representations. Isomorphic
representations have isomorphic inertia invariants with their residual
Frobenius action, proving equality of every local factor and hence the
zeta identity in the initial half-plane and by continuation. On the other
hand the degrees of the connected components are $4,1,1$ for $X$ and
$2,2,2$ for $Y$, precluding a scheme isomorphism.
\end{proof}

The same example over $\mathbb Z$ uses the respective rings of integers;
over $\mathbb Z[1/6]$ these are finite \emph{\'etale} covers. Thus even
complete Euler data forget the multiplication on an Artin algebra. A
applies to all fields here because they are abelian over $\mathbb Q$,
but does not restore the forgotten multiplication. The independent
transitive seven-element example CX-GEO-005 treats the connected
$G$-set distinction without asserting an arithmetic realization.

\subsection{AG-LAMBDA: the precise closure supplied by tensor operations}
\label{ag:lambda}

\textbf{Statement.} The finite direct sums of abelian Artin characters with
integer Tate twists are closed under tensor products, duals, exterior and
symmetric powers, and direct summands in characteristic zero. All their
$L$-functions have the A divisor-support conclusion. Replacing abelian
finite image by arbitrary finite image gives the corresponding M conclusion.
The same operations preserve the potentially CM absolute-Hodge category
used in CM01, so its M conclusion applies there as well.

\begin{proof}
Over an algebraic closure of the coefficient field an abelian Artin--Tate
object is a sum of lines $\chi_a(-j_a)$. A tensor product or a monomial
appearing in a symmetric or exterior power is a line
$(\prod_a\chi_a^{e_a})(-\sum_a e_aj_a)$. The character still has finite
abelian image and the Tate index is an integer. Duality replaces the
character and index by their inverses and negatives. Direct summands remain
in the same semisimple category, and descent preserves the asserted
factorization. The nonabelian assertion follows because all these
operations still factor through a finite group and commute with Tate
twists; apply ZT02 to their resulting Artin representations. Finally, the
potentially CM category is the tensor, dual, and summand closure in CM01.
Symmetric and exterior powers are the images of the usual idempotents in
the symmetric-group action on tensor powers, so belong to it too.
\end{proof}

The placement is SGA 6 V Section 2 and Section 7, especially (7.1.1),
and the determinant/$\lambda$-ring discussion carried in ZETA\_TWO ZT08
and SV01. These closures do not generate an arbitrary non-CM $GL_2$
motive from CM objects; enlarging a list of tensor operations is not the
missing analytic theorem.

\subsection{AG-SIGN: central signs with the necessary normalization}
\label{ag:sign}

\textbf{Statement.} Let $F$ be a nonzero holomorphic function on a
neighbourhood of $[c,\infty)$, real-valued there. Suppose $F$ has no zeros
on $(c,\infty)$ and is positive at one sufficiently large real point.
If $r=\operatorname{ord}_{s=c}F$, then
\[
 \frac{F^{(r)}(c)}{r!}>0.
\]
If, moreover, $F(c+z)=\varepsilon F(c-z)$ near $0$, with
$\varepsilon\in\{1,-1\}$, then $(-1)^r=\varepsilon$.
For a real, standardly normalized, holomorphic Hecke or CM $L$-function,
the absence of real zeros to the right of the centre follows from M
together with its usual analytic continuation and Euler-product theory.
The analogous Dirichlet statement uses A whenever no pole is present.

\begin{proof}
A continuous real nonvanishing function has constant sign on the connected
interval $(c,\infty)$. The normalization makes that sign positive.
Write $F(c+t)=a_rt^r+O(t^{r+1})$, with $a_r\ne0$. For $t>0$ the
left-hand side is positive, so $a_r>0$ after division by $t^r$ and passage
to the limit. Comparing the leading terms in the functional equation gives
$a_r=\varepsilon(-1)^r a_r$, hence the parity assertion.
\end{proof}

\paragraph{AG-HEIGHT (what this proves for a supplied height formula).}
If $F(c)=0$, then $F'(c)\ge0$: it is positive when $r=1$ and zero when
$r>1$. Consequently, an independently established identity
$H(Z)=C F'(c)$ with a specified constant $C>0$ yields $H(Z)\ge0$.
The elementary deduction is complete; the existence, domain, and
normalization of the height identity are extra inputs. It does not assert
that every homologically trivial cycle has such a formula, or that its
$L$-function has simple central vanishing. For an identity involving the
first nonzero derivative, the sign is strictly positive. For a formula with
an additional sign, that sign must be retained.

\paragraph{AG-SIGN-NORM (a counterexample to the unnormalized inference).}
Zero locations and functional-equation parity alone do not determine the
sign: $F(s)=s-c$ and $-F(s)$ have exactly the same zeros and the same odd
functional equation, but their central derivatives have opposite signs.
This is a counterexample to the stated analytic inference, not an example
of two standardly normalized arithmetic $L$-functions.

These findings refine \path{OURS_HODGE_BRIDGE_SKETCH_20261121.md},
Section 3.5. Their archive placement is cote 44, scans 2, 4 and 6,
Theorems 4--5, and scans 8--9, the letter to Serre of 12 August 1964,
as reported in that attached sketch. No original scan has been inferred
from this sign argument. The remaining arithmetic Hodge-index gate is
the cycle/height comparison for the stated class of cycles, not a sign
convention suppressed in a formula.

\subsection{AG-CRITICAL: why off-centre critical integers gain little}
\label{ag:critical}

\textbf{Statement.} Suppose $N$ is a pure, integral-weight, explicitly
Hecke-factored motive of weight $w$, with its usual functional equation
and archimedean factors. At an endpoint of the unitary strip, require each
factor in the specified signed Hecke factorization (or its dual, when the
functional equation is used) to be nonpolar. This factorwise requirement
prevents a denominator pole from producing a zero. Let $m$ be critical in
Deligne's sense: $L_\infty(N,s)$ and $L_\infty(N^\vee,1-s)$
are both regular at $s=m$. If $m\ne(w+1)/2$, then the standard
zero-free half-plane and boundary for unitary Hecke $L$-functions,
together with the functional equation, already give $L(N,m)\ne0$.
M is not needed for this conclusion.

The precise unitary boundary theorem used here is
D.~Ramakrishnan--R.~J.~Valenza, \emph{Fourier Analysis on Number Fields}
(1999), Section 7.7, Theorem 7-28, printed p.~289: for every number
field and every unitary idele class character, there is no zero at
$1+it$. Section 7.4, Theorem 7-19, printed p.~279, supplies the completed
continuation and functional equation, with its pure norm-character pole
exceptions retained. These are classical inputs independent of M.

\begin{proof}
After unitary normalization set $z=m-w/2$. If $w$ is even, $z$ is an
integer; if $w$ is odd, $z$ is a half-integer. An integer $m$ with
$0<z<1$ therefore has $z=1/2$, and is precisely the excluded centre.
For every other $m$ one has $z\ge1$ or $z\le0$. In the first case
standard unitary Hecke nonvanishing in $\Re z\ge1$, away from a pole,
gives the result for each factor and hence its finite signed product.
In the second case the functional equation relates the completed value
to its dual at $1-z\ge1$. Criticality ensures that both archimedean
factors are finite; gamma functions have no zeros. The nonzero dual
value therefore gives a nonzero original value. The factorwise nonpolar
condition excludes a Tate pole or its inverse being misread as a finite
nonzero special value. For a product with nonnegative exponents, the
usual no-pole condition implies this condition.
\end{proof}

Thus the strongest new real-axis information from M lies between the
off-centre integer points and at signs of central leading coefficients;
it does not evaluate regulators, periods, central ranks, or integral
torsion. This is the special-value distinction in ZETA\_TWO Section 4
and in the Hodge sketch's edge-versus-centre table. The relevant motivic
conventions are Deligne, \emph{Valeurs de fonctions $L$ et p\'eriodes
d'int\'egrales}, Sections 1.1--1.3; the period conjectures numbered 1.8
and 2.8 are separate inputs.

\subsection{AG-DET: Hecke versions of the square-variable determinant}
\label{ag:det}

The following formulation records exactly which analytic data are needed.
Let $\Lambda$ be a nonzero entire function of order at most one, with
\[
 \Lambda(c+z)=\varepsilon\Lambda(c-z),\qquad
 \varepsilon\in\{1,-1\},
\]
and let $r=\operatorname{ord}_{s=c}\Lambda$. Assume that every other zero
is $c+i\gamma$, with $\gamma$ real and nonzero. Choose one positive
$\gamma$ from each pair, with multiplicity, and write the sequence as
$\gamma_1,\gamma_2,\ldots$ in nondecreasing order, so $\gamma_1$
denotes the least positive zero when the sequence is nonempty.
Define the removable entire function
\[
 H(t)=t^{-r}\Lambda(c+it),\qquad
 D(w)=H(\sqrt w)/H(0).
\]
No choice of square-root branch affects $D$.

\textbf{Statement.} The diagonal operator
$B=\operatorname{diag}(\gamma_j^{-2})$ is positive and trace class, and
\[
 D(w)=\prod_j(1-w\gamma_j^{-2})=\det(I-wB).
\]
In particular this applies under M to standard entire selfdual completed
Hecke functions and to the entire completed $L$-function of a CM elliptic
curve, with its centre $c=1$. The usual order, functional equation,
and completion are additional classical inputs, not assertions supplied
by the word ``GRH''. Poles must first be removed in a separately specified
entire completion; arbitrary meromorphic Artin quotients are not included.

\begin{proof}
Comparison of the first nonzero Taylor coefficients shows
$\varepsilon=(-1)^r$. Hence $H(-t)=H(t)$, so $D$ is entire by the even
power series for $H$. Its order is at most one half: a bound of order
$1+\delta$ for $H(t)$ gives a bound of order $(1+\delta)/2$ for $D(w)$.
Its zeros, with multiplicity, are exactly $\gamma_j^2$.
The convergence exponent of the zeros of an entire function of order at
most one is at most one. In particular
$\sum_j|\gamma_j|^{-2}<\infty$. Equivalently the diagonal positive
operator $B$ has finite trace norm. Hadamard factorization for $D$, whose
order is strictly less than one, has no nonconstant exponential factor
and uses genus-zero factors. Since $D(0)=1$, its constant factor is one.
This gives the displayed product. The Fredholm determinant of a diagonal
trace-class operator is that same absolutely locally uniformly convergent
product, proving the identity. An exact analytic locator for the
factorization used here is C.~McMullen, \emph{Advanced Complex Analysis},
Math 213a course notes, version 16 November 2026, Theorem 3.16,
printed p.~92 (PDF p.~93); Theorems 3.10--3.12 give the zero-exponent
bound used just above. This is classical background, not a new claim
dependent on that revision date.
\end{proof}

This extends the prior zeta realization BR-TR-01 to a precisely specified
Hecke domain. Its relationship to cote 1 scans 79--80 and 111--112 is
the trace/Fredholm theorem reported in
\path{OURS_INTAKE_RESEARCH_SA23_20261120.md}, Section 4.1,
geo-H-lidskii and geo-H-genus. Those sheets are not themselves an
arithmetic RH claim. The operator here is defined from the zeros; the
construction has not supplied a cohomology theory or an operator defined
independently from prime data.

\subsection{AG-MOMENT and AG-HANKEL: recoverable positive spectral data}
\label{ag:moment}

Assume the hypotheses of AG-DET and infinitely many distinct positive
zeros. Put $\lambda_j=\gamma_j^{-2}$ and
\[
 s_k=\operatorname{Tr}(B^k)=\sum_j\lambda_j^k\quad(k\ge1).
\]
The logarithmic determinant expansion is
\[
 -\log D(w)=\sum_{k\ge1}\frac{s_k}{k}w^k
 \qquad (|w|<\gamma_1^2).
\]

\textbf{AG-MOMENT.} For every $k\ge1$,
\[
 s_k^{-1/k}\le\gamma_1^2\le\frac{s_k}{s_{k+1}}.
\]
The left bounds increase to $\gamma_1^2$; the right bounds decrease to it.
If $m_1$ is the multiplicity of the least positive zero, then
$\lim_{k\to\infty}s_k\gamma_1^{2k}=m_1$.

\begin{proof}
Let $a=\max_j\lambda_j=\gamma_1^{-2}$. Then
$s_k\ge a^k$ and $s_{k+1}\le a s_k$, proving the two inequalities.
The $\ell^p$ norm monotonicity gives
$s_{k+1}^{1/(k+1)}\le s_k^{1/k}$. One direct proof is to divide by
$s_k^{1/k}$: the resulting nonnegative numbers have $k$-th powers
summing to one, are each at most one, and thus have $(k+1)$-th powers
summing to at most one. Hence the left bounds increase.
Cauchy--Schwarz applied to
$\lambda_j^{k/2}$ and $\lambda_j^{(k+2)/2}$ gives
$s_{k+1}^2\le s_ks_{k+2}$, so the right bounds decrease.

There are finitely many indices with $\lambda_j=a$, since $\sum\lambda_j$
is finite. For $k\ge1$,
$0\le(\lambda_j/a)^k\le\lambda_j/a$, whose sum is finite.
Dominated convergence shows
$s_k/a^k\to m_1$. It follows that both $s_k^{1/k}$ and
$s_{k+1}/s_k$ tend to $a$, which proves both limits and the
multiplicity formula.
\end{proof}

\textbf{AG-HANKEL.} For all integers $h\ge1$ and $d\ge0$, the matrix
\[
 (s_{h+i+j})_{0\le i,j\le d}
\]
is positive definite. Its determinant is therefore strictly positive.

\begin{proof}
For a nonzero vector $a=(a_0,\ldots,a_d)$ and
$P(X)=\sum_{i=0}^d a_iX^i$, its Hermitian quadratic form equals
\[
 \sum_{j\ge1}\lambda_j^h|P(\lambda_j)|^2.
\]
This converges because the $\lambda_j$ are bounded and
$\sum\lambda_j^h<\infty$. All terms are nonnegative. A nonzero
polynomial has only finitely many distinct zeros, whereas the positive
$\lambda_j$ have infinitely many distinct values. At least one term
is positive. This proves strict positive definiteness.
\end{proof}

These consequences strengthen the two-by-two moment inequalities carried
in BR-TR-02. They describe actual determinant coefficients available in
the cote 1 analytic framework, while leaving the arithmetic/geometric
realization gate unchanged. They assert no simplicity and no
random-matrix statistics.

\subsection{AG-SIGNED: a finite moment witness for a surviving pole}
\label{ag:signed}

Consider a finite signed product of specified entire selfdual functions
of order at most one, with a common centre $c$, each satisfying the
zero-location hypotheses of AG-DET. Let $r\in\mathbb Z$ be the net
central order. Divide out the central Laurent factor and normalize its
nonzero constant term. Its square-variable quotient has the form
\[
 D(w)=\prod_{\lambda\in\mathcal L}(1-w\lambda)^{m_\lambda},
 \qquad m_\lambda\in\mathbb Z\setminus\{0\},\quad
 \lambda>0,\quad \sum_\lambda |m_\lambda|\lambda<\infty,
 \qquad D(0)=1.                                      \tag{AG-S1}
\]
Here $\mathcal L$ has no positive accumulation point, and the signed
multiplicities are taken \emph{after} cancellation. The form (AG-S1)
follows by applying AG-DET to each factor. More generally the proposition
below applies whenever (AG-S1) is independently supplied. Put
\[
 s_k=\sum_\lambda m_\lambda\lambda^k\quad(k\ge1),\qquad
 H_d=(s_{1+i+j})_{0\le i,j\le d}.
\]
These moments are the local logarithmic coefficients
$-\log D(w)=\sum_{k\ge1}s_kw^k/k$ near $0$.

\textbf{Statement.} All the matrices $H_d$ are positive semidefinite
if and only if $m_\lambda>0$ for every $\lambda\in\mathcal L$.
Equivalently, the normalized quotient has no noncentral poles. The
central order condition $r\ge0$ must be checked separately. If a
negative multiplicity survives, a finite matrix $H_d$ already detects it,
with the explicit witness described below.

\begin{proof}
Set $a=\sup\mathcal L$ and introduce the finite signed measure
\[
 \mu=\sum_\lambda m_\lambda\lambda\,\delta_\lambda
 \quad\hbox{on }[0,a].
\]
The empty support case is immediate. For every complex polynomial
$P(X)=\sum_{i=0}^d b_iX^i$, absolute convergence gives
\[
 \sum_{i,j=0}^d\overline b_i s_{1+i+j}b_j
       =\int |P(x)|^2\,d\mu(x).
\]
Positive multiplicities therefore imply all the required inequalities.

Conversely, suppose $m_{\lambda_0}<0$. The point $\lambda_0>0$ is
isolated in the compact set $\{0\}\cup\mathcal L$. Let
\[
 \delta=\operatorname{dist}
   (\lambda_0,(\{0\}\cup\mathcal L)\setminus\{\lambda_0\})>0,
 \qquad R>a+\lambda_0,
 \qquad q=1-\delta^2/R^2<1.
\]
For each nonnegative integer $N$, the polynomial
\[
 P_N(x)=\left(1-\frac{(x-\lambda_0)^2}{R^2}\right)^N
\]
has degree at most $2N$, takes the value $1$ at $\lambda_0$, and has
$|P_N(\lambda)|\le q^N$ at every other point of $\mathcal L$.
Writing $W=\sum_{\lambda\ne\lambda_0}|m_\lambda|\lambda$, we obtain
\[
 \int |P_N|^2\,d\mu
   \le -|m_{\lambda_0}|\lambda_0+Wq^{2N}.
\]
Choose $N$ with $Wq^{2N}<|m_{\lambda_0}|\lambda_0$; take $N=0$
if $W=0$. Then the coefficient vector of $P_N$ is a negative
quadratic-form witness for the matrix $H_{2N}$, of size $2N+1$.
This proves the converse and its finite-witness assertion. In (AG-S1)
a negative net multiplicity is exactly a pole at $w=\lambda^{-1}$,
which proves the analytic equivalence. Removing the central Laurent
factor erased its order, explaining the separate condition on $r$.
\end{proof}

Geometrically this strengthens the cote 1 determinant bridge in the
meromorphic setting. A or M can first localize an identified divisor;
positivity of these finite moment matrices is then an exact additional
condition preventing cancellation from leaving poles. It is not supplied
by zero location alone. In particular, universal Hecke GRH plus Brauer
induction does not prove this positivity, and hence does not prove Artin
holomorphy. The criterion is for the specified quotient: symmetrizing
an arbitrary nonselfdual function with its dual can cancel a pole between
the two factors, so holomorphy of that product must not be attributed to
each unsymmetrized factor. The local character-matrix test in MT-ART02
is a complementary criterion at an individual Artin divisor point.

\subsection{AG-EXT: the precise information that the determinant forgets}
\label{ag:ext}

The earlier findings BR-MIX-01 and ZETA\_TWO SV01--SV02 are retained.
At a good prime an adapted matrix for a mixed object is block triangular,
so its determinant equals that of its weight-graded object. The nonsplit
Kummer extension from the unit $2$ over $\mathbb Z[1/2]$ and its split
Tate object therefore have the same incomplete $L$-function. The omission
of ramified factors is essential: inertia invariants need not preserve a
short exact sequence. These findings refute recovery of extensions from
semisimplified good Euler data; they do not assert that all richer
arithmetic data agree.

The exact independent background locator is Deligne--Goncharov,
\emph{Groupes fondamentaux motiviques de Tate mixte},
arXiv:math/0302267v2 (27 February 2003), (2.1.3), printed p.~14
(PDF p.~16):
\[
 \operatorname{Ext}^1(\mathbb Q(0),\mathbb Q(n))=
 \begin{cases}
 0&n\le0,\\
 \mathcal O_S^\times\otimes\mathbb Q&n=1,\\
 K_{2n-1}(k)\otimes\mathbb Q&n\ge2.
 \end{cases}
\]
Propositions 2.2--2.3 on the same page identify the corresponding graded
free motivic Lie algebra, while Section 2.4 explicitly distinguishes a
free presentation from a canonical set of generators. These statements
are unconditional background. They do not calculate integral torsion or
the pro-$\ell$ kernels in the prompt's cote 149 programme.

\subsection{AG-KUMMER: an integral obstruction on the actual cote 149 sheets}
\label{ag:kummer}

The original Montpellier PDF was inspected here, not merely the intake.
It has a cover page: PDF p.~80 has repository footer 79, and so on.
At PDF p.~80/footer 79, equations (13)--(17) relate an inertia-normalizer
extension, its $H^2$ class, and Kummer cohomology. PDF p.~81/footer 80,
(18)--(19), identifies the relevant line bundle as the normal bundle of
the section. PDF p.~82/footer 81, (20)--(21), and PDF p.~85/footer 84,
(22), describe the $H^1$ and completed-unit ambiguity. These are the
actual geometric objects on the cited sheets. Their link to LM6
A.44--A.45 is supplied by ORDER item 31; no theorem about an irregular
prime or $K_8(\mathbb Z)$ is printed on these pages.

Here is a precise integral consequence, independent of A and M. Let $S$
be a scheme on which $\ell$ is invertible, put $P=\operatorname{Pic}(S)$,
and let $\operatorname{Br}'(S)=H^2_{\text{ét}}(S,\mathbb G_m)$.
Write $\widehat P_\ell=\varprojlim_nP/\ell^nP$ and
$T_\ell A=\varprojlim_n A[\ell^n]$, with transition maps multiplication
by $\ell$. There are exact sequences
\begin{align*}
 0&\longrightarrow\widehat P_\ell
   \longrightarrow\varprojlim_nH^2_{\text{ét}}(S,\mu_{\ell^n})
   \longrightarrow T_\ell\operatorname{Br}'(S)\longrightarrow0,\\
 0&\longrightarrow\Gamma(S,\mathcal O_S^*)^{\wedge}_\ell
   \longrightarrow\varprojlim_nH^1_{\text{ét}}(S,\mu_{\ell^n})
   \longrightarrow T_\ell P\longrightarrow0.
\end{align*}
We deliberately use inverse-limit cohomology: identifying it with a
chosen definition of continuous cohomology requires its own
$\varprojlim{}^1$ check. If $P$ is finitely generated, the second right
term is zero. A line bundle $\mathcal L$ has zero compatible integral
$\ell$-adic Chern class if and only if its class in $P$ is torsion of
order prime to $\ell$.

\begin{proof}
The Kummer exact sequence gives
\[
 0\longrightarrow P/\ell^nP
 \longrightarrow H^2(S,\mu_{\ell^n})
 \longrightarrow\operatorname{Br}'(S)[\ell^n]\longrightarrow0,
\]
and the analogous degree-one sequence with units modulo $\ell^n$th
powers and $P[\ell^n]$. The exact source is Stacks Project,
Lemma 59.28.1, tag 03PK, and its immediately following cohomology
sequence. The transition maps on the left terms are surjective.
Consequently a compatible sequence on the right lifts compatibly in
the middle: after choosing a lift at level $n$, correct any lift at
level $n+1$ by a lift of its discrepancy in the left term. This recursive
construction proves right-surjectivity after taking inverse limits;
left-exactness is immediate. It proves both displayed sequences.

Now write the finitely generated group $P$ as
$\mathbb Z^a\oplus P[\ell^\infty]\oplus P_{\ell'}$, with finite torsion
summands. Its $\ell$-adic completion is
$\mathbb Z_\ell^a\oplus P[\ell^\infty]$, and the kernel of completion
is exactly $P_{\ell'}$. Also no nonzero compatible infinitely
$\ell$-divisible torsion sequence can occur in a finitely generated
group, so $T_\ell P=0$. The injection in degree two then proves the
Chern-class criterion.
\end{proof}

\paragraph{AG-KUMMER-2: a fully explicit integral/rational separation.}
Let $K=\mathbb Q(\sqrt{-35})$, $\omega=(1+\sqrt{-35})/2$, and
$S=\operatorname{Spec}\mathbb Z[\omega,1/2]$. The invertible ideal
$\mathfrak p=(3,\omega)$ determines a line bundle $\mathcal L$ whose
compatible integral $2$-adic Chern class is nonzero, although its image
after tensoring with $\mathbb Q_2$ is zero.

\begin{proof}
The ring of integers is $\mathbb Z[\omega]$ and its discriminant is
$-35$. The minimal polynomial $X^2-X+9$ is irreducible modulo $2$,
so $(2)$ is the only prime over $2$, and is principal. Modulo $3$ it is
$X(X-1)$; thus the primes over $3$ are
$\mathfrak p=(3,\omega)$ and $\overline{\mathfrak p}=(3,\omega-1)$.
Since $N(\omega)=9$ and $\omega$ is a unit modulo
$\overline{\mathfrak p}$, unique ideal factorization gives
$\mathfrak p^2=(\omega)$. The ideal $\mathfrak p$ is not principal:
an integral generator would have norm $3$, whereas
\[
 N(a+b\omega)=a^2+ab+9b^2=(a+b/2)^2+35b^2/4
\]
cannot be $3$ for integers $a,b$.

Minkowski's bound is $2\sqrt{35}/\pi<4$, so every ideal class has a
representative of norm $1$, $2$, or $3$. There is no ideal of norm $2$;
the two of norm $3$ have the same class of order two. Hence the class
group is $\mathbb Z/2$. The precise classical bound is Milne,
\emph{Algebraic Number Theory}, v3.08 (19 July 2020), Theorem 4.3,
PDF p.~72. Inverting $2$ quotients the class group by the classes of
the primes over $2$, which here are principal, so
$\operatorname{Pic}(S)=\mathbb Z/2$ as well. To see this last assertion
directly, fractional ideals of the localization are obtained by forgetting
the valuation at $(2)$; a relation introduced in the localization differs
from a principal ideal only by a power of that principal prime.

AG-KUMMER now injects this nonzero class into the inverse limit of
$H^2(S,\mu_{2^n})$. It has order two there, so rationalization kills it.
The total space of $\mathcal L$ has a zero section with normal bundle
$\mathcal L$, making the example a literal normal-bundle obstruction of
the kind appearing in cote 149. No claim identifying this particular
scheme's group-cohomology extension class is needed.
\end{proof}

This proves a concrete limit of rational information: a rationally zero
normal Chern class need not be integrally zero. A supplies Dirichlet GRH
for the quadratic field here, but changes none of this calculation.
The finite-level Kummer test is therefore a genuine extra gate for an
integral inertia-splitting argument. It does not refute the original
LM6 statements, whose comparison maps must still be supplied.

\subsection{AG-C34-TYPE and AG-C34-REPAIR: resolving the commutator notation}
\label{ag:c34repair}

The official cote 34 PDF has a cover: PDF p.~53 is manuscript/footer
p.~52, and PDF p.~64 is p.~63. Theorem 1.1 on the former page fixes a
field $k$ finitely generated over its prime field, a normal geometrically
connected $k$-scheme $X$ locally of finite type, a prime
$\ell\ne\operatorname{char}k$, and an open subgroup
$\pi\subseteq\pi_1(X)$. It defines $\pi_0$ to be the image in
$G_k$ and $I=\pi\cap\pi_1(X_{\bar k})$, in
\[
 1\longrightarrow I\longrightarrow\pi\longrightarrow\pi_0
 \longrightarrow1.
\]
Its $I(\ell)$ is the maximal abelian pro-$\ell$ quotient of $I$.
The theorem states that $I(\ell)$ is finitely generated over
$\mathbb Z_\ell$ and that its $\pi_0$-coinvariants are finite.
The proof was located and read in Sections 2--3, PDF pp.~70--78
(manuscript pp.~69--77), rather than guessed from the intervening
applications. Section 2.4, PDF p.~72, makes the dual $H^1$-invariant
reduction; Section 2.5, PDF pp.~73--74, uses a curve compactification,
its cyclotomic boundary, and the finite-field curve theorem after
spreading out. Sections 3.1--3.3, PDF pp.~75--78, reduce normal $X$ to
a generic hyperplane curve using surjectivity on fundamental groups.
These are classical geometric inputs; A and M do not enter them.

In the crossed-out Remark 1.17 on PDF p.~64, $H$ is a quotient of $I$
whose kernel is $\pi$-stable, and the readable final hypothesis is that
$H$ is a solvable pro-$\ell$ group. The preceding proposed
$\ell$-adic Lie-group language is itself struck out. The remark writes
\[
 1\longrightarrow H\longrightarrow\pi'\longrightarrow\pi_0
 \longrightarrow1
\]
and places the closure of $[\pi_0,H]$ as an open subgroup of $\pi'$.
This is a tentative, struck manuscript assertion, not a theorem.

\paragraph{AG-C34-TYPE (the literal ambient-group obstruction).}
For an infinite $\pi_0$, every commutator of a lift of an element of
$\pi_0$ with an element of $H$ lies in $H$, and its closure cannot be
open in $\pi'$. Indeed the quotient by $H$ is the infinite group
$\pi_0$, whereas a subgroup open in a profinite group has finite index.
For a concrete example satisfying the manuscript's geometric hypotheses,
take $X=\mathbb G_m/\mathbb F_q$ and $\ell\nmid q$. Its geometric
pro-$\ell$ quotient $H=\mathbb Z_\ell(1)$ is abelian and hence
solvable; the rational point $1$ gives a section. Arithmetic Frobenius
acts as multiplication by $q$, so the action commutators generate
$(q-1)\mathbb Z_\ell(1)$. This is open in $H$ but is not open in
$\pi'$, whose quotient is $G_{\mathbb F_q}\simeq\widehat{\mathbb Z}$.
Thus the literal ambient group in the inherited Q03 cannot be retained.
For nonabelian $H$ there is a second issue: $\pi_0$ acts canonically
only by outer automorphisms unless a section or lifts are specified.

\paragraph{AG-C34-REPAIR (the intrinsic corrected statement).}
Let $H$ be \emph{any} pro-$\ell$ quotient of $I$ with $\pi$-stable
kernel, and let $P$ be the resulting quotient extension of $\pi_0$ by
$H$. Define $K$ to be the closed subgroup generated by all
$[p,h]$, $p\in P$, $h\in H$. Then $K$ is open in $H$, and
\[
 H/K\simeq (H^{\mathrm{ab}})_{\pi_0},
 \qquad [H:K]\le |I(\ell)_{\pi_0}|<\infty.       \tag{AG-C34.1}
\]
In particular the solvability and Lie-group hypotheses are unnecessary
for this precise corrected conclusion.

\begin{proof}
Because $H$ is normal in $P$, conjugating a commutator $[p,h]$ by
$a\in P$ gives $[apa^{-1},aha^{-1}]$. Hence $K$ is normal in $P$.
It contains the closed commutator subgroup of $H$, so $H/K$ is abelian;
all conjugation by $P$ is trivial on it. Conjugation on
$H^{\mathrm{ab}}$ factors canonically through $\pi_0$, irrespective
of any section. Killing its action differences is exactly the additional
quotient by the images of the commutators $[p,h]$. This proves the
isomorphism in (AG-C34.1).

The surjection $I\twoheadrightarrow H$ induces an equivariant
surjection $I(\ell)\twoheadrightarrow H^{\mathrm{ab}}$ by the
universal property of the maximal abelian pro-$\ell$ quotient.
Taking coinvariants remains surjective: it is simply quotienting both
compact abelian groups by the closed subgroups generated by their
action differences. Theorem 1.1 therefore supplies a surjection from
the finite group $I(\ell)_{\pi_0}$ to $H/K$. This proves the bound.
Since $K$ is closed, finite index in the compact group $H$ implies
openness (the finite Hausdorff quotient is discrete).
\end{proof}

This repairs the archived Q03 as a theorem already implied by its
stated abelian predecessor. It does not identify the unqualified
symbol $[\pi_0,H]$ in every possible intended meaning. A smaller
commutator subgroup formed from a chosen section, or openness in a
different ambient group, needs a separately stated question. In
particular the repair is not described as an automatic proof of all
interpretations of the crossed-out remark.

\paragraph{AG-C34-INDEX (an explicit Frobenius refinement).}
More generally, let $M=H^{\mathrm{ab}}$ be a finitely generated
$\mathbb Z_\ell$-module with continuous $\Gamma$-action, let $T$ be
its finite torsion subgroup, and let $\Lambda=M/T$. Suppose some
$F\in\Gamma$ has no eigenvalue $1$ on
$\Lambda\otimes\mathbb Q_\ell$. Then
\[
 |\operatorname{coker}(F-1\mid M)|
   =|T^F|\ell^{v_\ell\det(1-F\mid\Lambda\otimes\mathbb Q_\ell)},
 \qquad
 |M_\Gamma|\le
 |T^F|\ell^{v_\ell\det(1-F\mid\Lambda\otimes\mathbb Q_\ell)}.
                                                               \tag{AG-C34.2}
\]
Equality in the second assertion holds if $\Gamma$ is topologically
generated by $F$.

\begin{proof}
The assumption gives $\Lambda^F=0$. Apply the snake lemma to
$0\to T\to M\to\Lambda\to0$ with vertical maps $F-1$.
It yields
\[
 0\longrightarrow T/(F-1)T\longrightarrow M/(F-1)M
 \longrightarrow\Lambda/(F-1)\Lambda\longrightarrow0.
\]
For an endomorphism of the finite group $T$, kernel and cokernel have
the same cardinality. Smith normal form over $\mathbb Z_\ell$ gives
the order of the last group as the displayed power of $\ell$.
This proves the equality. The continuous image $(F-1)M$ is compact
and hence closed; the full coinvariant quotient is a quotient of it.
If the powers of $F$ are dense in $\Gamma$, every action difference
is a limit of differences of powers, all of which belong to
$(F-1)M$, so the two quotients agree.
\end{proof}

For a smooth proper geometrically connected curve $C/\mathbb F_q$,
let $H$ be its full geometric pro-$\ell$ fundamental group and $J$ its
Jacobian. Then $H^{\mathrm{ab}}=T_\ell J$, compatibly with Galois
action, and the corrected commutator quotient is
\[
 H/\overline{[P,H]}
       \simeq T_\ell J/(F-1)T_\ell J
       \simeq J(\mathbb F_q)[\ell^\infty].             \tag{AG-C34.3}
\]
For the fundamental-group comparison, an exact modern primary locator
is J.~S.~Milne, \emph{Jacobian Varieties}, Proposition 9.1,
pp.~30--32, together with \emph{Abelian Varieties}, Remarks 15.3
and 15.5, pp.~28--29. To check the last isomorphism directly, use
$0\to T_\ell J\to V_\ell J\to J(\overline{\mathbb F}_q)
[\ell^\infty]\to0$. The fixed subgroup of $T_\ell J$ is zero:
$J(\mathbb F_q)$ is finite, so a compatible indefinitely divisible
sequence of its $\ell$-power torsion points vanishes. Thus $F-1$ is
invertible on $V_\ell J$, and the snake lemma gives (AG-C34.3).
Consequently the index is exactly $\ell^{v_\ell\#J(\mathbb F_q)}$.
This finite-field calculation even uses less than the full Weil bound.
In higher-dimensional smooth proper situations purity supplies the
corresponding no-eigenvalue-$1$ test, but finite-field purity is already
classical and no new global GRH input is involved.

\subsection{Archive and run consequences, with the original scope retained}
\label{ag:archive}

\paragraph{AG-C1.} Cote 1/79--80 and 111--112 gives the trace/Fredholm
framework, and SA23 geo-H-lidskii/geo-H-genus is its reported modern
reconstruction. A yields the earlier zeta positivity realization;
M yields AG-DET for the stated Hecke functions. AG-MOMENT and AG-HANKEL
give further consequences in that framework. A geometric operator defined
from primes remains a separate construction.

\paragraph{AG-C16.} Cote 16, scans 14--18, no.~7, concerns the integral
Lefschetz correspondence ring. Scans 44--50 contain the Coates letter's
conditions A, C and index positivity; the attached Hodge sketch uses a scan
numbering one higher than the GCT\_NSRH intake. These are algebraicity and
integrality problems. Neither A nor M constructs the required
correspondence or bounds its denominator. The sketch's explicit route is
the rank comparison on $X\times X$, not a claim that zero location
reverses Grothendieck's implication to the Weil conjectures.

\paragraph{AG-C34.} Original PDF p.~74, manuscript p.~73, (2.5.5),
has now been visually checked: its multiplicative $1$ denotes the trivial
invariant group, equivalently additive $0$. The finite-field curve theorem
is the stated input, so inherited Q04 is unchanged. The original
definitions and full proof route for Theorem 1.1 have also been read.
AG-C34-TYPE identifies the literal obstruction in Q03, and
AG-C34-REPAIR proves the coherent intrinsic correction. AG-C34-INDEX
gives its finite Frobenius/Jacobian index. No assertion of new purity or
of a consequence of A/M is made.

\paragraph{AG-C44.} Cote 44 Theorems 4--5 and the Serre letter, at the
locators in AG-SIGN, concern polarization and heights. AG-SIGN corrects
the central-derivative bridge but leaves cycle existence and the
height formula explicit. The supplied Hodge sketch also cites a separate
2026 CM-Hodge theorem for its abelian positivity consequences; those
claims are not consequences of either A or M and are carried with
their original dependency rather than independently certified here.

\paragraph{AG-C149 and AG-GT.} ORDER\_tonight, item 31, points to cote
149 scans 79--85, LM6 A.44--A.45, integral bridge B3, and the
irregular-prime/$\ell=2$ questions. GRH does not remove the integral
cohomology comparison or the specified class-group component gate.
The semi-graded Euler information in AG-EXT explains why the endpoint
zero theorems do not evaluate that extension data. In SA22 run 1,
A6-node-character is the quadratic Artin character times the cyclotomic
character; A gives effective detection of that finite quadratic sign.
M additionally covers finite nonabelian splitting fields in effective
Frobenius tests such as AG-FROB. These improvements do not reverse the
recorded refutation of run 2 A:inertia-maximality: finite-index inertia
subgroups have the same weights. Nor do they identify the profinite GT
kernel from the cyclotomic quotient. The existing supplied claim about
rational Deligne--Drinfeld freeness (CC-GT03, external 2026 input) is
carried separately from the full profinite Ihara question R-GTHAT.

\paragraph{AG-PURITY and AG-WM.} Finite-field purity and the usual
Grothendieck--Lefschetz trace formula are already classical. The full
geometric weight--monodromy statement remains represented in these
materials by BR-WM-02's finite exterior-power trace criterion; its
remaining inequalities are geometric hypotheses, not consequences of
global Dirichlet or Hecke zero locations. Likewise the supplied archive
questions about standard conjectures require algebraic cycles or
correspondences. No formal logical independence from A or M is asserted;
the precise point is that the displayed deductions do not supply those
inputs.

% END geometry_programme

% BEGIN complexity
\section{Complexity, effective algebra, and finite geometric tests}
\label{sec:complexity}

Here $A_D$ denotes the rational Dirichlet-GRH input designated by the task as
Alp\"oge's theorem. The letter $M$ denotes Masullo's all-number-field Hecke-GRH
claim, and every use of $M$ remains conditional. The input size of an arithmetic
problem includes the circuits, explicitly listed variables, rational constants
in binary, and any promised field representation. A geometric object without
an encoding has no associated complexity bound.

\subsection{Classical consequences and inherited results}

The following catalogue entries preserve the original finding identifiers of
the earlier runs. They are not advertised as new findings of this session.
An exact locator to an earlier proof is distinguished from a fresh reading of
that proof. In particular, the earlier complete file \path{CX_GRH.tex} is
not in tonight's attachment: its \path{IMPLICATIONS}, \path{LEDGER},
\path{QUESTIONS}, and targeted intake are available. Statements that depend
on its elaborate new finite-field constructions are carried from those records;
this session does not manufacture an independent verification of the absent
proof text.

\paragraph{CX-ALG-001 and CX-ALG-002: primality.}
Bach, \emph{Explicit bounds for primality testing and related problems}
(1990), Theorem 2, p.~372, and its application on p.~373 give the cutoff
$2(\log N)^2$ for a compositeness witness. Schoof, \emph{Four primality testing
algorithms} (2008), Theorem 2.2, pp.~105--106, states and proves the Miller test
with this cutoff using only quadratic Dirichlet characters. Therefore $A_D$
discharges its analytic input. Writing $N-1=2^s m$ with $m$ odd, test the usual
conditions $a^m=1$ or $a^{2^i m}=-1$ modulo $N$ for every base through the
cutoff, also checking nontrivial gcds. Even inputs are handled separately.
There are $O((\log N)^2)$ bases and each modular exponentiation takes a
polynomial number of bit operations. With $b=\lceil\log_2N\rceil$ and
multiplication cost $\mathcal M(b)$, the elementary implementation uses
$O(b^3\mathcal M(b))$ bit operations. This is the transfer already recorded
as CX\_GRH CX-N01 and ALPOGE\_GCT C-GRH-02.

The statement $\mathrm{PRIMES}\in\mathrm P$ was already unconditional:
Agrawal--Kayal--Saxena, \emph{PRIMES is in P}, Annals 160 (2004), Theorems 4.1
and 5.1, with their 2019 erratum. The new input validates this particular
short-witness test; it does not first place primality in $\mathrm P$.
Geometrically these are computations in $\mathbf G_m$ over $\mathbf Z/N\mathbf Z$
and, for AKS, finite cyclotomic algebras. The archival companion is cote 28,
scan 3 (group-action axioms), as mapped in the CX\_GRH return, not freshly
inspected from an archive scan in this session.

\paragraph{CX-ALG-003--CX-ALG-004: generators and nonresidues.}
Bach's Theorem 3, p.~372, gives units below $3(\log N)^2$ generating
$(\mathbf Z/N\mathbf Z)^\times$. Applied to the proper subgroup of $r$th
powers in $\mathbf F_p^\times$, it gives a polynomial search for a specified
prime $r\mid p-1$. Membership is tested by exponentiation to $(p-1)/r$;
the other factors of $p-1$ need not be known. This is CX\_GRH CX-N02.
The original Bach PDF was not freshly recovered here (the publisher returned
403); its exact locator is carried from the preceding checked run. Schoof's
primary source independently supplies the quadratic special case used above.
In the Kummer sequence this is finding a nonzero class in
$\mathbf F_p^\times/(\mathbf F_p^\times)^r$, or a nontrivial
$\boldsymbol\mu_r$-torsor. The established companion locator is SGA 4 IX 3.2.

\paragraph{CX-ALG-006--CX-ALG-008: primitive roots and extension fields.}
CX\_GRH CX-N03 constructs a certified generator of $\mathbf F_p^\times$ when
the full factorization of $p-1$ is supplied, by constructing each primary
component and multiplying them. Shoup, \emph{Searching for Primitive Roots
in Finite Fields} (1992), author version Theorem 3, p.~4, and Section 6,
p.~13, gives a least primitive root of order
$O(r^4(\log r+1)^4(\log p)^2)$, where $r=\omega(p-1)$,
and hence $O((\log p)^6)$ under the
rational Dirichlet hypothesis. This is a guaranteed short candidate list
without a supplied factorization, as CX\_GRH CX-N04 explicitly records.
A candidate list does not identify or certify a particular primitive root.

The stronger extension-field construction CX\_GRH CX-N13 supplies a specified
$r$th nonresidue in a supplied $\mathbf F_{p^m}$ in
$\operatorname{poly}(m\log p)$ time when $r\mid p-1$ is prime. Its quadratic
case with even $m$ already has an unconditional construction. The theorem is
inherited from the full targeted check in
\path{OURS_INTAKE_CX_GRH_20261121.md}, Section 1. A fresh source audit
found a nonvanishing gap in one cited auxiliary formula; CX-ALG-016 gives
the exact counterexample, and CX-ALG-017--CX-ALG-018 give a fully checked
repair and re-proof of the needed construction. The defining
irreducible polynomial is part of the input. These are computations with
Kummer covers inside an explicitly represented finite \`etale algebra.

\paragraph{CX-ALG-009--CX-ALG-010: polynomial factorization over finite fields.}
The inherited theorem CX\_GRH CX-N15 has bound
\[
                 (d^{\log d}\log q)^{O(1)}
\]
for a degree-$d$ polynomial over a supplied $\mathbf F_q$. Under $M$ its
analytic input is satisfied by $\zeta_{\mathbf Q}$ and the particular
$\zeta_{E_\nu}$ used in the construction, where
$E_\nu=\mathbf Q(\zeta_{2\nu},p^{1/\nu})$ and $\nu\leq d^2$.
Its exact component locators are Adleman--Lenstra (1986), Theorem 1, p.~350,
Algorithm A, pp.~350--351; Bach--von zur Gathen--Lenstra (2001), Theorem 2,
p.~7; and the 2008 Ivanyos--Karpinski--R\'onyai--Saxena version, Main Theorem,
p.~3, Lemma 4.2, p.~14, Proposition 4.6, pp.~16--17.
The fresh primary published IKRS article confirms the quasipolynomial
factorization-versus-automorphism distinction (Theorems 1.1 and 1.3), and
its unconditional noncommutative zero-divisor result (Theorem 1.7).

The displayed bound remains quasipolynomial in $d$. The residual $A_D$-only
problem in CX\_GRH CX-U02 is the high-valuation branch
$p>d$, $r\nmid p(p-1)$ and
$v_r(p^{\operatorname{ord}_r(p)}-1)>1$: construct the maximal coefficient
Teichm\"uller unit needed by that descent, or replace the descent. No new
solution of that gate is asserted. Geometrically factorization computes
primitive idempotents, hence connected components, of a finite \`etale
algebra (SGA 1 V Section 7). It is a different computational problem from
factoring an integer with unknown prime divisors.

\paragraph{CX-ALG-011--CX-ALG-013: arithmetic Picard groups.}
Bach (1990), Theorem 4, p.~375, and its corollary, p.~376, give the familiar
$12(\log|D_K|)^2$ class-group generator bound, using class-character Hecke
$L$-functions over $K$. Thus $M$ supplies the premise. This is CX\_GRH
CX-N07 and ALPOGE\_BRIDGES CC-F03/F03b. If the Hilbert class field is
abelian over $\mathbf Q$, $A_D$ suffices by Dirichlet factorization; the
weaker condition that $K/\mathbf Q$ be abelian does not establish this.
A factor base is not the full lattice of relations between its ideal
classes. Relation collection, saturation, regulator information and their
running times remain separate inputs in CX\_GRH CX-N08 and CX-U03.
In Grothendieck's language the class group is
$\operatorname{Pic}(\operatorname{Spec}\mathcal O_K)$ and principal divisors
supply the relation lattice. The established companion locators are SGA 4
IX Theorem 3.3, SGA 6 V 3.10 and VI 3.2--3.3, and cote 2 scan 5 as recorded
by CX\_GRH.

\paragraph{CX-ALG-019: a complete class-group algorithm retains randomness.}
For a supplied negative fundamental discriminant $-d$, Biasse--Erukulangara,
\emph{A proof of the conjectured run time of the Hafner--McCurley class
group algorithm}, Main Theorem, author PDF p.~4, Algorithm 4, p.~15,
and Corollary 1, p.~16, give a Las Vegas computation of
$\operatorname{Cl}(\mathcal O_{\mathbf Q(\sqrt{-d})})$ with expected bound
\[
 \exp\!\left(\left(\frac3{\sqrt8}+o(1)\right)
                         \sqrt{\log d\log\log d}\right).
\]
Their Conjecture 1, p.~4, explicitly states number-field Hecke-GRH; hence
$M$ supplies the hypothesis. The proof also uses the class-group Cayley
graph expansion of Jao--Miller--Venkatesan (2009), Theorem 1.1 and
Corollary 1.3. Rational quadratic Dirichlet-GRH alone does not supply all
these class-character estimates. This is the complete algorithmic transfer
already recorded as CX\_GRH CX-N12, not a new algorithm of this session.

The source's bounded trial returns a certified answer or FAILURE, with
success probability $1-d^{-1+o(1)}$ and the displayed time bound. Repeating
on detected failure gives the stated expected Las Vegas bound. The
certification in Algorithm 4 compares the order $h'$ of the group presented
by the collected relations with $2h^*$, where $h^*\le h<2h^*$ and $h$
is the actual class number. Since $h\mid h'$ and $h'>0$, the condition
$h'<2h^*$ forces $h'=h$. Thus the restart does not accept an incorrect
relation lattice. The full earlier proof and its detailed source check are
credited to CX-N12; the primary theorem, analytic-input definition,
Algorithm 4 and Corollary 1 were freshly inspected here.
The runtime is subexponential in $\log d$, not polynomial in that input
length, and the algorithm remains randomized. Its geometric output is the
finite Picard group, with actual principal-divisor relations rather than
only a factor base. The prior companions are SGA 4 IX Theorem 3.3,
SGA 6 V 3.10 and VI 3.2--3.3, and cote 2 scan 5.

\paragraph{CX-ALG-020: small-prime walks in ray class groups.}
Jao--Miller--Venkatesan (2009), Theorem 1.1, author PDF p.~3,
Remark 1.2(a) and Corollary 1.3, p.~4, give the following exact transfer.
For degree-$n$ $K$, modulus $\mathfrak m$, and
$Q=|D_K|N\mathfrak m$, take the narrow ray class group $G$ and the
multiset of degree-one prime ideals of norm at most $x$, coprime to
$\mathfrak m$, together with their inverses. Conditional on $M$, its
Cayley adjacency eigenvalues satisfy
\[
 \lambda_1=2\operatorname{li}(x)+O(n\sqrt x\log(xQ)),\qquad
 |\lambda_\chi|=O(n\sqrt x\log(xQ))\quad(\chi\ne1).
\]
The constants here are absolute. For fixed $n$ and $B>2$,
$x\ge(\log Q)^B$, and sufficiently large $Q$, a walk of length
$t\ge C_{n,B}\log|G|/\log\log Q$ reaches any fixed subset $T$ with
probability at least $|T|/(2|G|)$. Quotients inherit the bound.
The precise hypothesis is GRH for the Hecke characters actually occurring;
for $K=\mathbf Q$, $A_D$ suffices. Random walks remain randomized.
Geometrically these are bounded-prime correspondences on arithmetic
Picard/ray-class quotients, complementing CX\_GRH CX-N12; the Picard
companions are SGA 4 IX Theorem 3.3 and cote 2 scan 5.

\paragraph{CX-ALG-022: the ordinary elliptic isogeny graph.}
The same paper's Theorem 1.5, p.~5, applies conditionally on $M$ to ordinary
elliptic curves over a fixed $\mathbf F_q$, with fixed point count and
the same geometric endomorphism order. The graph of isogenies of prime
degree less than $(\log(4q))^B$, $B>2$ fixed, has the stated expansion
bound. Theorem 1.6, pp.~6--7, transfers a discrete-logarithm solver working
on a fraction $\mu$ of that level to every curve on it with an expected
$\mu^{-1}\operatorname{poly}(\log q)$ overhead. These are conditional
classical transfers; they do not supply such a solver. The Hecke input is
for imaginary-quadratic ring-class characters, rather than an assertion
about a non-CM elliptic curve's global L-function. The geometric link is
the CM ideal-class action (Proposition 3.1, p.~11) and Deuring lifting
(Theorem 3.3, p.~12); ORDER item 20 supplies the abelian-scheme/isogeny
programme, alongside the inherited Picard comparison of cote 2 scan 5.

\paragraph{CX-HN-001--CX-HN-004: the exact Nullstellensatz input.}
Koiran's DIMACS 96-27, Section 2, Theorem 2, pp.~5--6, proves
$\mathrm{HN}\in\mathrm{AM}\subseteq\Pi_2^{\mathrm P}$ under its GRH input.
The August 8, 1997 erratum corrects the lower estimate in Theorem 8 and a
prime-product expression in Theorem 10; it explicitly preserves the main
statements. The primary author-hosted report and the entire erratum were
retrieved in this session. Section 5, Theorem 9 and Corollary 1, printed p.~12 (PDF page~13),
show that the analytic estimate is for the bounded primitive solution
field itself. A Galois closure is not required. These are the earlier
CX\_GRH CX-K01--K02; the circuit and homogeneous nonzero forms are also
stated in GCT V Theorem 2.10, p.~25.

Consequently $M$ gives the unrestricted transfer, whereas $A_D$ gives the
bounded abelian-witness promise of CX\_GRH CX-K03. All the arithmetic size
bounds in that promise matter: field degree, logarithmic discriminant and
denominators have the stated single-exponential bounds in the input length.
The universal replacement by an abelian solution field was already refuted
by $x^3-2$ in CX\_GRH CX-K04 and ALPOGE\_GCT C-GRH-03. This refutes that
specific replacement, not every possible future $A_D$-only proof.
The geometric object is a finite \`etale spread of the solution field over
an open of $\operatorname{Spec}\mathbf Z$; degree-one Frobenius orbits give
rational reductions. The companion locators are SGA 1 V Section 7 and
cote 98 scans 11--12 as carried in CX\_GRH.

\paragraph{CX-NN-001--CX-NN-004: normalization maps.}
The preceding CX\_GRH CX-K05 refines GCT V Theorem 3.4, pp.~27--28, to
$\mathrm{FP}^{\Sigma_2^{\mathrm P}}$ under $M$ for equation-given,
standard-graded algebras with explicit ambient variables and polynomial
output bits. The source proof's prefix-existence queries already have
$\Sigma_2$ form. Its unconditional polynomial-space statement remains
available. The resulting map
$\operatorname{Spec}R\to\mathbf A^{\dim R}$ is finite; the companion is
EGA II 6.1.4--6.1.5, and cote 98 scan 44 in the prior crosswalk.

The historical EXPH transfers for succinct small separating systems are
GCT V Theorems 4.10, p.~34, and 5.12, carried in CX\_GRH CX-K06. The
optimal-cardinality h.s.o.p. bounds of Theorems 4.1, p.~30, and 5.5 retain
their efficient defining-equation premise. Their output conventions cannot
be silently changed. Earlier unconditional improvements already exist:
Forbes--Shpilka, ECCC TR17-163, Theorem 7.1, p.~21, constructs the relevant
hitting set in polynomial space for unary $n,s,r$; its introductory
Theorems 1.2--1.3 record normalization consequences. Guo--Saxena--Sinhababu
(2019), Theorems 1.5--1.6, are recorded in CX\_GRH CX-K07. These are not
new consequences of tonight's analytic inputs.

\subsection{New formulations with complete proofs}

\paragraph{CX-ALG-005: deterministic prime-field square roots.}
\emph{Statement.} Under $A_D$, given an odd prime $p$ and $a\in\mathbf F_p$,
one can find all roots of $X^2-a$ or certify their absence in deterministic
time polynomial in $\log p$.

\emph{Proof.} The case $a=0$ is immediate. Euler's criterion decides whether
$a$ is a square. If it is not, return no roots. Otherwise find a quadratic
nonresidue by the deterministic search of CX-ALG-004. Apply the classical
Tonelli--Shanks algorithm with this nonresidue. More specifically, Algorithm
3.34 of Menezes--van Oorschot--Vanstone, \emph{Handbook of Applied Cryptography},
p.~100, has only one random step, its choice of the nonresidue in Step 2.
All its remaining steps are modular exponentiation, gcd inversion and a
loop of length at most $v_2(p-1)\leq\log_2p$; Fact 3.35 gives a polynomial
bit bound for those operations. Replacing Step 2 by the bounded deterministic
search preserves polynomial time and the algorithm's proved correctness.
The resulting roots are $r,-r$. Completing the square factors every quadratic
with nonzero leading coefficient; in characteristic two the separate field
case is trivial for $\mathbf F_2$. This does not provide a square-root oracle
modulo an integer whose factors are unknown.\hfill$\square$

\paragraph{CX-ALG-014: a compositeness witness need not give the usual gcd factor.}
\emph{Statement.} The rule ``take a Miller witness and compute
$\gcd(a,N)$ or $\gcd(a^m\pm1,N)$'' does not always output a nontrivial divisor.

\emph{Proof.} Take $N=35$, so $N-1=2\cdot17$, and $a=2$. Since
$2^{12}=1\pmod{35}$, one has $2^{17}=32\pmod{35}$. This is neither $1$ nor
$-1$ modulo $35$, so $2$ is a Miller witness. Yet
\[
 \gcd(2,35)=\gcd(32-1,35)=\gcd(32+1,35)=1.
\]
This is a counterexample to precisely the displayed extraction rule.
It is not a lower bound for general integer factorization and does not
rule out other algorithms.\hfill$\square$

\paragraph{CX-ALG-015: integer factorization with a supplied exponent multiple.}
\emph{Statement.} Given integers $N\ge2$ and $T\ge1$, with the promise
\[
       \lambda'(N):=\operatorname{lcm}_{p\mid N}(p-1)\mid T,
\]
the complete factorization of $N$ is computable deterministically in time
polynomial in $\log N+\log T$ under $A_D$. This is the input-length form
of Miller (1976), Lemma 5 and Theorem 3, p.~310. In particular, a supplied
value of $\varphi(N)$ or $\lambda(N)$ meets the promise. Neither value is
asserted to be available for an arbitrary $N$.

\emph{Proof of the stated input-length form.} Remove powers of two, test
primality, and detect perfect powers by integer-root computations. These
are polynomial-time operations. For a perfect power, recurse on its base
and multiply the resulting exponents. It remains to split an odd composite
with at least two distinct prime divisors. Write $T=2^s t$ with $t$ odd.
For every $a\leq\lceil3(\log N)^2\rceil$, first compute $\gcd(a,N)$;
return a proper nontrivial gcd if found, and otherwise test, for units $a$,
\[
             \gcd\bigl(a^{2^j t}-1,N\bigr),\qquad 0\leq j<s.
\]
Powers are computed modulo $N$ by repeated squaring.

To prove that some test splits $N$, put
$e_p=v_2(p-1)$ and $e=\max_{p\mid N}e_p$. If $e_p=e$ and $e_q<e$ for
two divisors, the quadratic character $a\mapsto(a/p)$ is a nontrivial
character of $(\mathbf Z/N\mathbf Z)^\times$. If all $e_p=e$, use
$a\mapsto(a/p)(a/q)$ for distinct divisors $p,q$ instead. This too is
nontrivial, by the Chinese remainder theorem. Bach's generating bound
CX-ALG-003 implies that some tested unit has character $-1$.
In the first case its reduction modulo $p$ has multiplicative order with
2-part $2^e$, whereas the 2-part modulo $q$ is smaller. In the second
case one Legendre symbol is $-1$ and the other is $1$, with the same
consequence after interchanging $p,q$. Since $t$ is a multiple of the odd
part of each $p-1$, one obtains
\[
       a^{2^{e-1}t}=-1\pmod p,\qquad
       a^{2^{e-1}t}=1\pmod q.
\]
The corresponding gcd is divisible by $q$ and not by $p$, so is proper.
The promise gives $s\ge e$, hence this test occurs in the list. This
argument uses only reductions modulo the distinct primes; repeated prime
factors in $N$ cause no problem.

Every divisor $D$ encountered in the recursion still satisfies
$\lambda'(D)\mid T$. A recursion tree has $O(\log N)$ leaves and nodes
after perfect-power compression. Each node uses polynomially many bases
and at most $\log_2T$ squarings per base, so the total bit cost is
polynomial in the displayed input length. Geometrically the known exponent
annihilates the multiplicative groups of the reduced residue fibers;
quadratic characters separate two unknown fibers and a gcd recovers their
different supports. This is a precise promise case of the integer-factoring
gate in CX-ALG-014, not an algorithm for producing the exponent promise.
\hfill$\square$

\paragraph{CX-ALG-016: the fixed-generator resolvent needs a nonzero test.}
\emph{Statement.} The literal fixed-generator formula in
Bhargava--Ivanyos--Mittal--Saxena (2017), Theorem 3.1, author PDF p.~3,
can be zero. This does not refute their nonresidue-construction conclusion;
CX-ALG-017 supplies the missing nonzero choice.

\emph{Proof.} Let $q=7$, $r=3$, $d=3$, and
$E=\mathbf F_7[X]/(X^3-2)$, with $\alpha=X\bmod(X^3-2)$.
The nonzero cubes in $\mathbf F_7$ are $1,-1$, so this cubic has no root
and is irreducible. Take the primitive cube root $\zeta=4$. Then
$\alpha^7=4\alpha$ and $\alpha^{49}=2\alpha$, whence
\[
 L=\alpha+4\alpha^7+2\alpha^{49}
   =(1+16+4)\alpha=0.
\]
Consequently $L^{(7^3-1)/3}=0$, whereas the stated right-hand side is
$\zeta^{-1}=2$. The displayed formula and its hypotheses were checked
against the rendered original PDF page, not merely its text extraction.
The invalid step is division by $L$ when inferring
$L^{q^k-1}=\zeta^{-1}$ from $L^{q^k}=\zeta^{-1}L$.
\hfill$\square$

\paragraph{CX-ALG-017: a basis-search repair for the resolvent.}
\emph{Statement.} Let $E=\mathbf F_q[X]/(f)$ be a supplied field of degree
$d=rk$, where $r$ is prime. Suppose a primitive $r$th root of unity
$\zeta\in\mathbf F_q$ is supplied, and assume $r\mid q-1$ for odd $r$,
or $4\mid q-1$ when $r=2$. Then an $r$th nonresidue in $E$, and its
nonresidue norm in $\mathbf F_q$, can be found deterministically in time
polynomial in $d\log q$. No analytic hypothesis is needed for this statement.

\emph{Proof.} Write $Q=q^k$ and define the $\mathbf F_q$-linear map
\[
       P(z)=\sum_{i=0}^{r-1}\zeta^i z^{Q^i}.
\]
As a polynomial in $z$, it is nonzero and has degree
$Q^{r-1}<Q^r=|E|$. Thus it cannot vanish at every element of $E$.
Since it is linear, it cannot vanish on every element of the power basis
$1,\alpha,\ldots,\alpha^{d-1}$. Scan this basis and choose
$\beta=P(\alpha^j)\ne0$. Cyclically shifting the sum and using
$z^{Q^r}=z$ gives $\beta^Q=\zeta^{-1}\beta$, hence
$\beta^{Q-1}=\zeta^{-1}$.

Put $S=1+Q+\cdots+Q^{r-1}$. For odd $r$, write $Q=1+rc$ and expand
modulo $r^2$ to get
\[
 S\equiv r+rc\frac{r(r-1)}2\equiv r\pmod{r^2}.
\]
For $r=2$, the hypothesis gives $S=1+Q\equiv2\pmod4$.
In both cases $S/r\equiv1\pmod r$. Therefore
\[
 \beta^{(q^d-1)/r}
   =\left(\beta^{Q-1}\right)^{S/r}=\zeta^{-1}\ne1.
\]
This proves that $\beta$ is an $r$th nonresidue. Its norm satisfies
\[
 \left(N_{E/\mathbf F_q}\beta\right)^{(q-1)/r}
      =\beta^{(q^d-1)/r}=\zeta^{-1},
\]
so is also a nonresidue. There are $d$ basis candidates and $r\le d$
terms per test; all exponent bit lengths are $O(d\log q)$.
Field arithmetic in the supplied representation and these exponentiations
therefore take polynomial time.\hfill$\square$

For the inherited CX\_GRH CX-N13, obtain $\zeta$ under $A_D$ from a
prime-field nonresidue. If odd $r\mid m$, apply the theorem with
$q=p,d=m$. If $r\nmid m$, a prime-field nonresidue $a$ remains a
nonresidue upstairs, since
\[
 a^{(p^m-1)/r}
   =\left(a^{(p-1)/r}\right)^{1+p+\cdots+p^{m-1}}
   =\left(a^{(p-1)/r}\right)^m\ne1.
\]
The remaining quadratic-even case is proved next. Thus the prior
construction is preserved with an explicit nonvanishing justification.
In the Kummer-cover language, the repair chooses a nonzero vector in the
specified Frobenius eigenspace before passing to its torsor class; a
particular chosen generator need not have a nonzero projection.

\paragraph{CX-ALG-018: the even-degree quadratic case, with no GRH.}
\emph{Statement.} For odd $p$ and even $m$, a supplied
$E=\mathbf F_p[X]/(f)$ of degree $m$ admits a deterministic
$\operatorname{poly}(m\log p)$ construction of a quadratic nonresidue.
This is a fresh complete re-proof of the even-degree clause of inherited
CX\_GRH CX-N13, not a new claim of priority over BIMS17 Corollary 1.3.

\emph{Proof.} Put $Q=p^{m/2}$ and $\alpha=X\bmod f$.
If $p\equiv1\pmod4$ or $4\mid m$, then $Q\equiv1\pmod4$.
The element $u=\alpha^Q-\alpha$ is nonzero, since $\alpha$ has degree
$m$ and is not in the subfield of degree $m/2$. It satisfies $u^Q=-u$.
Consequently
\[
       u^{(p^m-1)/2}=(-1)^{(Q+1)/2}=-1,
\]
so it is the required nonresidue.

It remains that $p\equiv3\pmod4$ and $m\equiv2\pmod4$.
Compute the subfield $K=\mathbf F_{p^2}\subseteq E$ as the kernel of
the $\mathbf F_p$-linear map $z\mapsto z^{p^2}-z$.
Choose $v\in K\setminus\mathbf F_p$ from a computed basis and set
$u=v^p-v\ne0$. Then $u^p=-u$ and $c=u^2\in\mathbf F_p$ is a
nonresidue: a square $c$ would force $u\in\mathbf F_p$.
Because $-c$ is a square, the exponent formula
$t=(-c)^{(p+1)/4}$ gives $t^2=-c$. Thus $i=u/t$ satisfies $i^2=-1$,
and $K=\mathbf F_p(i)$ is explicitly represented.

Start with $z=i$, of order four. Its norm to $\mathbf F_p$ is one.
Whenever the current $z=a+bi$ has norm one, its order is a power of two
at least four; hence $b\ne0$, as the 2-part of $\mathbf F_p^\times$ has
order two. The nonzero numbers
\[
                 \delta_+=(a+1)/2,\qquad \delta_-=(a-1)/2
\]
have product $-b^2/4$, a nonsquare. Exactly one is a square in
$\mathbf F_p$. Determine which by Euler's criterion, choose its square
root $x=\delta^{(p+1)/4}$, and set $y=b/(2x)$.
Using $a^2+b^2=1$ verifies directly that
$(x+yi)^2=a+bi=z$. The norm of this root is $1$ if
$\delta=\delta_+$ and $-1$ if $\delta=\delta_-$.
Replace $z$ by this root, thereby doubling its order, until the norm is
$-1$.

This loop terminates after $O(\log p)$ steps. Indeed the norm-one subgroup
has order $p+1$, whereas the 2-part of $K^\times$ has order
$2^{v_2(p^2-1)}=2^{1+v_2(p+1)}$. An element's order cannot continue
doubling inside the former subgroup. At termination
\[
        z^{(p^2-1)/2}
          =N_{K/\mathbf F_p}(z)^{(p-1)/2}=-1.
\]
The extension $E/K$ has odd degree $m/2$, so this $z$ remains a nonresidue
in $E$. All operations are modular exponentiation or linear algebra in
the supplied $m$-dimensional field, with polynomial bit cost.
\hfill$\square$

The quadratic square-root algorithm then applies in $E$ with this known
nonresidue. In characteristic two the separate square-root problem is
inverse Frobenius, $a\mapsto a^{2^{m-1}}$, and there is no quadratic
nonresidue. These are finite-\'etale and Kummer computations; SGA 1 V
Section 7 and SGA 4 IX 3.2 are the same companions as for CX-ALG-008.

\paragraph{CX-ALG-021: the genus quotient only needs rational Dirichlet-GRH.}
\emph{Statement.} For an imaginary quadratic field $K$, put
$C=\operatorname{Cl}(\mathcal O_K)$ and $G=C/C^2$. The small-prime
spectral and mixing bounds in CX-ALG-020 hold on $G$ under $A_D$, with
$n=2$, $\mathfrak m=(1)$ and $Q=|D_K|$. No Hecke-GRH assumption for
the remaining, non-genus characters of $C$ is needed. If $C$ has exponent
two, this covers the whole class-group graph.
Moreover, the images of prime ideals of prime norm at most
\[
                         (8\log|D_K|+15)^2
\]
generate $C/C^2$ under $A_D$.

\emph{Proof.} Let $H$ be the Hilbert class field of $K$.
By class field theory $\operatorname{Gal}(H/K)=C$. Complex conjugation
preserves $H$: the conjugate of an unramified abelian extension of $K$
is again such an extension, and the maximal one is unique. Its action on
$C$ is inversion, since
$\mathfrak a\overline{\mathfrak a}=(N\mathfrak a)$ is principal.
Actual complex conjugation has order two and restricts nontrivially to
$K$, so
\[
       \operatorname{Gal}(H/\mathbf Q)=C\rtimes\langle c\rangle,
                    \qquad cac^{-1}=a^{-1}.
\]
The subgroup $C^2$ is normal. Hence $F=H^{C^2}$ is Galois over
$\mathbf Q$, and its Galois group is
$(C/C^2)\times\langle c\rangle$, since inversion is trivial on a group
of exponent two. In particular $F/\mathbf Q$ is abelian.

Every character $\chi$ of $C/C^2$ extends to two one-dimensional
characters $\chi_+,\chi_-$ of $\operatorname{Gal}(F/\mathbf Q)$, by
assigning $c$ the values $1,-1$. Directly on the two cosets of
$\operatorname{Gal}(F/K)$,
\[
       \operatorname{Ind}_{\operatorname{Gal}(F/K)}^{
                       \operatorname{Gal}(F/\mathbf Q)}\chi
                         =\chi_+\oplus\chi_-.
\]
Artin induction of L-functions and abelian class field theory over
$\mathbf Q$ therefore express the class-character Hecke function
$L_K(s,\chi)$ as a product of the two associated Dirichlet L-functions.
This includes $\chi=1$, for which the product is $\zeta_K$.
The input $A_D$ gives precisely the zero-location premises for all these
functions.

For the explicit generator bound, fix any nontrivial character $\chi$ of
$C/C^2$. Its kernel in $C$ defines a quadratic unramified extension
$F_\chi/K$ inside $H$. The preceding conjugation argument shows that
$F_\chi/\mathbf Q$ is biquadratic. Its absolute discriminant is
$|D_K|^2$ by the relative-discriminant tower formula and unramifiedness.
Apply Bach--Sorenson Theorem 5.1 to the nonidentity element of
$\operatorname{Gal}(F_\chi/K)$, regarded inside
$\operatorname{Gal}(F_\chi/\mathbf Q)$. Its zeta hypothesis follows from
$A_D$, since this degree-four field is abelian over $\mathbf Q$.
The resulting unramified rational prime satisfies
\[
            p\le(4\log(|D_K|^2)+2.5\cdot4+5)^2
                    =(8\log|D_K|+15)^2.
\]
Its Frobenius is trivial on $K$, so it splits there, and either prime
ideal above it has $\chi$-value $-1$. If the prime ideals in the displayed
range generated a proper subgroup of $C/C^2$, a nontrivial character of
that elementary two-group would be trivial on the subgroup. The prime just
constructed contradicts this. This proves the generator assertion without
paying for the degree of the entire genus field.

The character-by-character proof of JMV Theorem 1.1, author p.~9,
uses only the L-function of the individual character being estimated.
The spectrum of the quotient graph consists exactly of characters
factoring through $C/C^2$ (Remark 1.2(a), p.~4). Thus all its eigenvalue
bounds hold under the premises just established. Lemma 2.1, pp.~7--8,
and the proof of Corollary 1.3, p.~10, give the claimed mixing bound
with $|G|=|C/C^2|$. If $G$ is trivial, the statement is immediate.
\hfill$\square$

This is the genus-level part of the arithmetic Picard action. It sharpens
the input split in CX-ALG-013: the full Hilbert class field may be
nonabelian over $\mathbf Q$, while its genus quotient is abelian and
already covered by $A_D$. It does not give the missing estimates for the
other class characters or a general deterministic class-group algorithm.

\paragraph{CX-GEO-001: existential constructible feasibility.}
\emph{Statement.} Conditional on $M$, deciding whether a quantifier-free
Boolean circuit of rational polynomial equalities and inequations has a
complex solution belongs to $\mathrm{AM}$. Its emptiness problem belongs to
$\mathrm{coAM}\subseteq\Sigma_2^{\mathrm P}$. Polynomial expressions are
given by arithmetic circuits and all original variables are explicitly listed.

\emph{Proof.} Represent each atomic equality $f(x)=0$ by a new Boolean
variable $b$ and a new variable $u$, with equations
\[
 b(b-1)=0,\qquad bf(x)=0,\qquad (1-b)(u f(x)-1)=0.
\]
If $b=1$, these require $f(x)=0$. If $b=0$, they require $f(x)\ne0$
and an inverse $u=f(x)^{-1}$. Conversely any assignment to $x$ extends
with $b$ equal to the truth value of $f(x)=0$; when $b=1$, any $u$ works.
An inequation is the negation of its equality atom.
For each Boolean gate, introduce its output variable $t$ and use
\[
 t-uv=0\quad(\mathrm{AND}),\qquad
 t-u-v+uv=0\quad(\mathrm{OR}),\qquad
 t+u-1=0\quad(\mathrm{NOT}),
\]
where $u,v$ here denote that gate's Boolean inputs. Impose $t(t-1)=0$
for each output if desired, and impose that the final output equals $1$.
Induction through the acyclic Boolean circuit proves that the resulting
polynomial system is feasible exactly when the original formula is true
at some $x$. The construction shares each original arithmetic circuit;
it does not expand a Boolean formula into disjunctive normal form.
Its size and rational-constant bit lengths increase polynomially.
Apply CX-HN-001 to the resulting system and take complements for emptiness.
The reduction concerns equations and inequations over $\mathbf C$, not
ordered-field inequalities. Geometrically it realizes the given constructible
locus as the image of an explicitly presented affine scheme and preserves
nonemptiness.\hfill$\square$

\paragraph{CX-GEO-002: radical membership.}
\emph{Statement.} Conditional on $M$, for
$I=(f_1,\ldots,f_s)\subseteq\mathbf Q[x_1,\ldots,x_n]$ and $g$ given
by rational arithmetic circuits, radical membership $g\in\sqrt I$ is in
$\mathrm{coAM}$ and nonmembership is in $\mathrm{AM}$.

\emph{Proof.} By the Hilbert Nullstellensatz over $\mathbf C$,
$g\notin\sqrt{I\mathbf C[x]}$ exactly when there is a point $x$ with
$f_i(x)=0$ for every $i$ and $g(x)\ne0$. This is equivalent to feasibility of
\[
               f_i(x)=0\quad(1\leq i\leq s),\qquad y g(x)-1=0.
\]
For completeness, the rational and complex membership tests agree. If
$g^m\in I\mathbf C[x]$, a representation by a finite list of polynomial
multiples gives a finite linear system with rational coefficients for those
multipliers' coefficients. A rational linear system having a complex solution
has a rational solution, by row reduction. Thus $g^m\in I$. The converse
follows by extension of scalars. The new system has polynomial size, so
CX-HN-001 gives the asserted classes. In scheme language the question is
whether the principal open $D(g)$ in $\operatorname{Spec}(\mathbf Q[x]/I)$
has a geometric point. Nilpotents matter only through the radical in this
particular problem.\hfill$\square$

\paragraph{CX-GEO-003: reduced containment and equality.}
\emph{Statement.} Conditional on $M$, inclusion and equality of affine
algebraic sets specified by rational arithmetic-circuit equations are in
$\mathrm{coAM}$. This statement does not assert the same bound for
scheme-theoretic ideal inclusion.

\emph{Proof.} Put $X=V(f_1,\ldots,f_s)$ and
$Y=V(g_1,\ldots,g_t)$. Then $X\nsubseteq Y$ exactly when
\[
        f_i(x)=0\quad(1\leq i\leq s),\qquad
        \sum_{j=1}^t u_jg_j(x)=1                              \tag{CX.1}
\]
is feasible. Indeed, the last equation is impossible if all $g_j(x)$ vanish;
if one does not vanish, its inverse supplies a solution with all other
$u_j=0$. For $t=0$ the sum is zero, making (CX.1) impossible, correctly
expressing $Y=\mathbf A^n$. For an empty $f$-list there is no restriction
on $x$, also correctly. Thus noninclusion is in $\mathrm{AM}$ and
inclusion is in $\mathrm{coAM}$.

Equality can be handled by the same single-reduction method. Introduce
$b$ with $b(b-1)=0$. Impose
\[
\begin{aligned}
 b f_i(x)&=0,& (1-b)g_j(x)&=0,\\
 b\left(\sum_j u_jg_j(x)-1\right)&=0,&
 (1-b)\left(\sum_i v_if_i(x)-1\right)&=0.
\end{aligned}
\]
For $b=1$ this is a point of $X\setminus Y$; for $b=0$ it is a point
of $Y\setminus X$. Hence inequality is one HN instance of polynomial size.
Finally, the ideals $(x)$ and $(x^2)$ have the same complex zero set but
are unequal. They show exactly why this reduction tests underlying reduced
sets rather than the structure sheaves or ideal inclusions.
For GCT this test applies when defining equations of the orbit closures
are actually supplied; producing such equations remains a separate gate.
\hfill$\square$

\paragraph{CX-GEO-004: a finite Frobenius test with prescribed excluded primes.}
\emph{Statement.} Let $L/\mathbf Q$ be finite Galois, let $G$ be its Galois
group, and let $E_i=L^{H_i}$ for subgroups $H_i\leq G$ ($i=1,2$).
Let $\chi_i$ be the permutation character of $G/H_i$. Set $n=[L:\mathbf Q]$
and $D=|D_L|$. Suppose the right halves of the critical strips for
$\zeta_L$ and $\zeta_{\mathbf Q}$ are zero-free, as supplied by $M$, or by
$A_D$ if $L/\mathbf Q$ is abelian. If $\chi_1\ne\chi_2$, there is a prime
unramified in $L$ such that
\[
       p\leq B_L=(4\log D+\tfrac52 n+5)^2                  \tag{CX.2}
\]
and the numbers of degree-one prime ideals of $E_1$ and $E_2$ above $p$
differ. Equivalently the numbers of $\mathbf F_p$-points of the two integer
models differ.

For any nonempty finite set $S$ of rational primes put
$d=\prod_{p\in S}p$, $F=\mathbf Q(\sqrt d)$, and $L'=LF$.
The same assertion holds with $p\notin S$ and bound $B_{L'}$, under the
corresponding hypotheses for $\zeta_{L'}$ and $\zeta_{\mathbf Q}$.
If $n'=[L':\mathbf Q]$ and $D'=|D_{L'}|$, an explicit upper estimate is
\[
 n'\leq2n,\qquad \log D'\leq2\log D+n\log|D_F|
                     \leq2\log D+n\log(4d),               \tag{CX.3}
\]
so the entirely explicit larger bound
\[
        \bigl(8\log D+4n\log(4d)+5n+5\bigr)^2            \tag{CX.4}
\]
also suffices. If $S=\varnothing$, use (CX.2) directly.

\emph{Proof.} Choose $\sigma\in G$ such that
$\chi_1(\sigma)\ne\chi_2(\sigma)$. Bach--Sorenson,
\emph{Explicit Bounds for Primes in Residue Classes} (1996), Theorem 5.1,
p.~1729, gives an unramified prime whose Frobenius conjugacy class is that
of $\sigma$, with bound (CX.2). Their Section 2, p.~1720, specifies the
needed zero-free right-half assumption for the top-field zeta function and
$\zeta_{\mathbf Q}$; no Artin-holomorphy conjecture is required.
For an unramified prime the Frobenius orbits on $G/H_i$ correspond to the
prime ideals above $p$, and the orbit lengths are their residue degrees.
Its fixed points therefore count precisely the degree-one primes.
This proves the first assertion.

For the exclusion refinement, every prime of $S$ ramifies in $F$. For odd
$p\mid d$ this follows from the quadratic discriminant formula; if $2\in S$,
then $d$ is even and $D_F=4d$, so 2 ramifies as well. Consequently every
prime in $S$ ramifies in $L'$, since its ramification index in the subextension
$F/\mathbf Q$ divides its index upstairs. Choose any lift of $\sigma$ to
$\operatorname{Gal}(L'/\mathbf Q)$ and apply the same theorem there. The
prime obtained is unramified in $L'$, hence avoids $S$, and its restriction
to $L$ has the required character difference.

We justify (CX.3) without an unstated discriminant estimate. If $F\subseteq L$,
then $L'=L$ and the inequalities are immediate. Otherwise $L$ and $F$ are
linearly disjoint because $F$ is quadratic. If $\omega$ is an integral basis
generator for $\mathcal O_F$, the order $\mathcal O_L[\omega]\subseteq
\mathcal O_{L'}$ has relative discriminant generated by $D_F$. The relative
field discriminant divides the order discriminant. The tower formula gives
$D'\leq D^2|D_F|^n$. Also $n'=2n$ and $|D_F|\leq4d$.
Substitution in (CX.2) yields (CX.4). If $L/\mathbf Q$ is abelian, $LF$ is
also abelian, so the stated $A_D$ alternative survives the exclusion step.
The exceptional source case $L=\mathbf Q$ cannot have unequal characters
for the two subfields in this statement, so causes no gap.\hfill$\square$

\emph{Consequences and limits.} Agreement at all the tested primes implies
$\chi_1=\chi_2$. Hence the corresponding characteristic-zero permutation
representations are isomorphic and their Artin $L$-functions, including
ramified factors formed on inertia invariants, agree. In particular
$\zeta_{E_1}=\zeta_{E_2}$. This is a finite trace certificate, not necessarily
an efficient algorithm in a short presentation of $E_i$: a common Galois
closure may have large degree. Nor is equality of permutation characters
an isomorphism theorem for finite covers, as the next result proves.
The Grothendieck framework is SGA 1 V Section 7 and the Frobenius trace
formalism; cote 98 scans 11--12 are the finite-$G$-set archival companion
carried from the CX\_GRH crosswalk.

\paragraph{CX-GEO-005: equal characters do not identify transitive $G$-sets.}
\emph{Statement.} Let $V=\mathbf F_2^3$, $G=\operatorname{GL}(V)$,
$X=V\setminus\{0\}$ and $Y=V^*\setminus\{0\}$ with the natural and dual
actions. Then $X$ and $Y$ are transitive seven-element $G$-sets with identical
permutation characters but are not isomorphic as $G$-sets.

\emph{Proof.} For $g\in G$, the number of fixed nonzero vectors is
$2^{\dim\ker(g-I)}-1$. The corresponding number of fixed nonzero covectors
is $2^{\dim\ker((g^{-1})^{\mathsf T}-I)}-1$. The matrices
$(g^{-1})^{\mathsf T}-I$ and $g^{\mathsf T}-I$ have the same rank, since
they differ by multiplication by an invertible matrix and sign. Transpose
preserves rank, so both fixed-point counts are equal.

Suppose the two $G$-sets were isomorphic. The stabilizer $H$ of a nonzero
vector, say $e_1$, would stabilize some nonzero covector and therefore its
kernel, a plane $W\subset V$. Such an invariant plane does not exist.
If $e_1\in W$, its image $W/\langle e_1\rangle$ is a line in
$V/\langle e_1\rangle$, invariant under all of
$\operatorname{GL}_2(\mathbf F_2)$ because $H$ induces that full group.
No line is invariant under that group. If $e_1\notin W$, projection identifies
$W$ with the two-dimensional quotient, so $W$ is a complement to
$\langle e_1\rangle$. But $H$ contains all transvections
$(a,v)\mapsto(a+u(v),v)$ with $u$ a linear functional on that quotient.
Choosing $u(v)=1$ for a nonzero projected vector in $W$ moves its lift by
$e_1$, which cannot stay in $W$. Both cases contradict invariance.
Thus no $G$-set isomorphism exists.\hfill$\square$

This is a representation-theoretic counterexample, needing no GRH and no
inverse-Galois assertion. For any connected finite \`etale $G$-torsor that
realizes this group, its two associated connected covers exhibit the same
limitation. No particular unverified realization over $\mathbf Q$ is used.
The example distinguishes the Galois category from its linearized
representation category; the latter is what Frobenius traces recognize.

\paragraph{CX-GEO-006: orbit membership with an explicit action.}
\emph{Statement.} Conditional on $M$, suppose an algebraic group $G$ and its
action on an affine space have polynomial-length rational circuit descriptions,
with every inverted denominator explicitly listed. Given rational points
$v,w$, deciding whether $w\in G(\mathbf C)v$ belongs to $\mathrm{AM}$.

\emph{Proof.} Introduce coordinates for $g$, impose the supplied equations
and inequations defining $G$, and impose the coordinate equations
$g\cdot v=w$. For example, in $\operatorname{GL}_n$ an additional variable
$z$ and equation $z\det(g)-1=0$ enforce invertibility; the determinant has
a polynomial-size arithmetic circuit. Inverted denominators of a rational
action are handled by the same inverse-variable equations. The resulting
constructible-feasibility instance has polynomial size and is equivalent to
orbit membership by definition of the action. Apply CX-GEO-001.
Orbit-closure membership is not represented by this system: points in the
boundary need not equal $g\cdot v$ for any $g$. If closure equations are
provided, CX-HN-001 and CX-GEO-003 apply to those equations; without them
the equation-access problem persists.\hfill$\square$

\paragraph{CX-GEO-007: polynomially parameterized orbit closures already admit a PSPACE test.}
\emph{Statement.} Let the ambient coordinates of $V$ be explicitly listed.
Suppose the action on a supplied rational vector $v$ is given by
division-free arithmetic circuits with rational coefficients for all its
coordinates. In either of the following
cases, membership of a rational point $w$ in the complex orbit closure is
in $\mathrm{PSPACE}$, without either GRH input: (a) the group is
$\operatorname{GL}_n$ and the map $g\mapsto g\cdot v$ extends by the supplied
circuits to a polynomial map on all $\operatorname{Mat}_n$; (b) the group is
$\operatorname{SL}_n$ and the supplied action coordinates are polynomial
in its matrix entries. This statement charges the explicit ambient
coordinate count in the input size.

\emph{Proof.} Guo--Saxena--Sinhababu (2019), Theorem 1.3, p.~5, together
with Lemma 4.1 and Theorem 4.2, pp.~15--16, gives a polynomial-space test
for whether zero lies in the Zariski closure of the image of a polynomial
map specified by arithmetic circuits: this is their approximate polynomial
satisfiability problem. In case (a), $\operatorname{GL}_n$ is dense in
$\operatorname{Mat}_n$. For a polynomial map $F$ on the latter, every closed
set containing $F(\operatorname{GL}_n)$ has closed inverse image containing
a dense subset, hence contains $F(\operatorname{Mat}_n)$. Their image
closures are equal. Apply the cited test to $F-w$.

For case (b), construct a polynomial map from an affine space into
$\operatorname{SL}_n$ whose image contains its Gaussian-elimination big
cell. Products of elementary unipotent matrices parameterize arbitrary
lower and upper unitriangular matrices. A diagonal matrix of determinant
one is a product of adjacent diagonal blocks
$\operatorname{diag}(t,t^{-1})$: for diagonal entries $d_1,\ldots,d_n$,
use $t_i=d_1\cdots d_i$ in positions $(i,i+1)$ for $1\leq i<n$.
Such a block belongs to the image of a
six-factor elementary-unipotent word, because, in its $2$ by $2$ block,
\[
 E_{12}(t)E_{21}(-t^{-1})E_{12}(t)
 E_{12}(-1)E_{21}(1)E_{12}(-1)
       =\operatorname{diag}(t,t^{-1}).
\]
Use independent parameters in each of the six positions of this fixed
word. The resulting map is polynomial and always has determinant one;
the displayed specializations show that it contains every such diagonal
block in its image. Place these words between the lower and upper
unitriangular words. Its image contains every $LDU$ decomposition in the
big cell, which is a dense open subset of $\operatorname{SL}_n$.
The number of elementary factors and parameters is $O(n^2)$ and the
matrix product has polynomial-size circuits. Composing with the supplied
action on $v$ gives a polynomial map with the same image closure as the
orbit map, by the same closed-inverse-image argument. Subtract $w$ and
apply the cited PSPACE algorithm. For $n=1$, the special-linear action is
trivial and membership reduces to $w=v$.\hfill$\square$

This gives an equation-free route for this exact orbit-closure encoding.
It does not give an AM bound from $M$ for approximate satisfiability, and
it does not turn an exponentially long list of ambient coefficients into
a polynomial-length input. Those are separate questions from the theorem.

\paragraph{CX-GEO-008: a finite Frobenius generating certificate.}
\emph{Statement.} Let $L/\mathbf Q$ be finite Galois with group $G$ and
assume the zeta hypotheses of CX-GEO-004. For each unramified rational
prime $p\leq B_L$, choose any representative $g_p$ of its Frobenius class.
The elements $g_p$ generate $G$. More strongly, a subgroup meeting every
Frobenius class observed in this range equals $G$. A sublist of at most
$\lfloor\log_2|G|\rfloor$ representatives suffices to generate $G$.
For a prescribed excluded set $S$, replace the bound by $B_{L'}$ from
CX-GEO-004 and use only primes outside $S$. As before, $A_D$ suffices for
abelian $L/\mathbf Q$, and $M$ supplies the general hypothesis.

\emph{Proof.} Let $H<G$ be proper and let $X=G/H$. Double counting gives
\[
 \sum_{g\in G}|\operatorname{Fix}_X(g)|
      =\sum_{x\in X}|\operatorname{Stab}_G(x)|=|G|.
\]
The identity fixes $|X|>1$ points. If every element fixed at least one,
the left side would exceed $|G|$. Hence there is $g$ with no fixed coset.
Since $g$ fixes $aH$ exactly when $a^{-1}ga\in H$, its conjugacy class
misses $H$. The explicit Chebotarev theorem used in CX-GEO-004 produces
that class at an unramified prime $p\leq B_L$. With excluded primes,
choose a lift to $\operatorname{Gal}(L'/\mathbf Q)$ and apply the same
argument at bound $B_{L'}$; every prime of $S$ ramifies in $L'$ and is
therefore avoided. Thus no proper subgroup can meet every observed class,
and in particular no proper subgroup contains all the chosen $g_p$.

Starting from the trivial subgroup, choose a listed $g_p$ outside the
subgroup generated so far until it equals $G$. Each proper enlargement
has index at least two, so after $k$ selections the generated subgroup has
order at least $2^k$. Consequently $k\leq\lfloor\log_2|G|\rfloor$.
For $G=1$ the empty list suffices.\hfill$\square$

In SGA 1 V Section 7 this is a finite quotient of the arithmetic
fundamental group, with a finite list of Frobenius conjugacy classes
certifying that no proper subgroup contains the image. The bound is
relative to the specified extension and its discriminant; it is not a
uniform presentation of the entire profinite fundamental group. The
archival companion is the finite-$G$-set discussion of cote 98 scans
11--12 recorded in the preceding CX\_GRH crosswalk, and the work target
is ORDER item 30's effective finite Galois-image question.

\paragraph{CX-GEO-009: bounded Laurent degenerations give an AM promise.}
\emph{Statement.} Let $F:\mathbf A^n\to\mathbf A^m$ and
$w\in\mathbf Q^m$ be specified by division-free rational arithmetic
circuits of total size $s$. Supply in unary an upper bound $D\ge1$ for
their syntactic degrees and an integer $B\ge0$.
Conditional on $M$, it is in $\mathrm{AM}$ to decide whether there are
Laurent polynomials
\[
 \beta_i(\varepsilon)=\sum_{j=-B}^{B}c_{i,j}\varepsilon^j,
 \qquad c_{i,j}\in\mathbf C,
\]
such that every coordinate of $F(\beta)-w$ belongs to
$\varepsilon\mathbf C[\varepsilon]$.
Consequently image-closure membership is in $\mathrm{AM}$ on the promise
that every positive instance admits such an arc with $D$ and $B$ bounded
polynomially in the original encoding length. This includes the polynomial
orbit parameterizations of CX-GEO-007 with the same arc promise.

\emph{Proof.} Use the coefficients $c_{i,j}$ as unknowns. At each circuit
gate $g$ of syntactic degree $d_g$, maintain arithmetic circuits for the
coefficients of its substituted Laurent polynomial in the range
$[-Bd_g,Bd_g]$. Input-variable gates have their prescribed coefficient
unknowns; a rational constant has only a constant coefficient. At an
addition or subtraction gate, add or subtract coefficient circuits.
At a multiplication gate, take the usual coefficient convolutions.
Syntactic degrees add at multiplication, so no exponent outside the stated
range can occur. Discard gates unused by the outputs; then each remaining
degree is at most $D$.

There are at most $2BD+1$ coefficient positions per gate, and each product
uses at most $(2BD+1)^2$ pairs. Sharing all earlier coefficient circuits
therefore gives a polynomial-size arithmetic-circuit construction in
$s,n,m,B,D$. For every output impose that all coefficients of exponents
$-BD,\ldots,0$ in $F(\beta)-w$ vanish. These are polynomial equations in
the $n(2B+1)$ unknown coefficients. The resulting HN instance has polynomial
encoding size and is feasible exactly when the displayed bounded Laurent
arc exists. Koiran's theorem gives the claimed AM bound under $M$.
The construction also covers $B=0$, when it is ordinary exact feasibility.

Any such arc puts $w$ in the image closure by GSS Lemma 4.1 and Theorem
4.2, pp.~15--16. Conversely the promise provides a bounded arc for every
point of that closure under consideration, proving the promise corollary.
All oracle queries concern this explicit coefficient system; its bounded
solution field is the exact number-field zeta input. An $A_D$ alternative
requires the bounded abelian-witness promise for this coefficient system,
not merely for the original rational points $v,w$.
\hfill$\square$

The generic GSS bound on Laurent precision can be exponential; writing
that bound in binary would not justify the polynomial-size coefficient
construction. This result therefore isolates a concrete smaller gate for
the promised families: a polynomially bounded degeneration, rather than
access to all equations of the orbit closure. It does not assert that
arbitrary GCT orbit closures meet that promise. The geometric objects are
finite Laurent jets of a parameterized degeneration; the group-action
framework is SGA 3 I 4.6.3.1 and 4.7.1--4.7.2, with cote 28 scan 3 as the
archival companion inherited from CX\_GRH. The work target is ORDER item
28's Koiran/GCT route.

\subsection{Gates retained rather than silently discharged}

\paragraph{CX-LIM-001: general derandomization and GCT.}
The exact remaining premises of the cited routes are uniform hitting sets
or independent circuit-hardness hypotheses. GCT V Theorem 4.7, p.~33,
assumes a specific permanent border-hardness bound. Its Lemma 4.8, on the
same page, extracts a hitting set from a sufficiently explicit separating
system. Neither construction is furnished by a statement about zeros of
Dirichlet or Hecke $L$-functions. This records the boundary of these proofs;
it is not a proof that no logical implication from GRH to general
derandomization could ever exist. No $\mathrm{BPP}=\mathrm P$, integer
factorization in $\mathrm P$, or $\mathrm{VP}\ne\mathrm{VNP}$ is claimed.

\paragraph{CX-LIM-002: positivity and purity.}
The earlier ALPOGE\_GCT P-F7 and P-F10 and PLETH PLE-F5/F7 already show
why pure geometry, nonnegative Hilbert data, efficient counting and actual
orbit-closure multiplicities are different inputs. Their established
counterexamples and corrected positivity-index implications are carried,
not rebranded. The new tests CX-GEO-002--CX-GEO-003 can verify specified
reduced-locus properties once equations are given. They do not construct a
uniform positive multiplicity formula or a strict asymptotic quotient-rank
inequality. Those are the surviving GCT gates.

\paragraph{CX-LIM-003: the separate counting premise.}
The attached CX\_GRH return also assumed an independent theorem denoted
$E$ and drew ETH/counting consequences in CX-E01--CX-E05. Tonight's two
analytic inputs $A_D$ and $M$ do not contain $E$. Those rows therefore do
not become consequences of these two inputs by being present in the same
archive. This is an input classification, not a new assertion about ETH.

% END complexity

% BEGIN root
\section{The two methods: precise transfers and repairs}
\label{sec:methods}

\subsection*{MT-A01. What the adopted Dirichlet theorem says adelically}
The supplied declarations \path{riemannHypothesis},
\path{LFunction_ne_zero_of_half_lt_re}, and
\path{LFunction_eq_zero_re} in \path{ALPOGE_Main.lean} have the
definitions printed in \path{ALPOGE_Statements.lean}.  We use these as the
cited input $A_D$, as requested.  The supplied archive contains these exported
declarations and the scan report, but not the imported tree
\path{Supplements.Kernels.MainTheoremsAll}; a kernel build of that tree is
not a check performed in this run.  This distinction concerns our evidence,
not the stipulated status of the input theorem.

Every unitary idele-class character of $\mathbb Q$ has the form
$\chi_0|\cdot|^{it}$, where $\chi_0$ is a finite-order Dirichlet character
and $t\in\mathbb R$.  Consequently $A_D$ gives the rank-one Hecke GRH over
$\mathbb Q$ for these characters, not just its finite-order specialization.
Indeed the norm map has a continuous section and compact kernel, which for
$\mathbb Q$ is $\widehat{\mathbb Z}^{\times}$ after the sign convention at
infinity.  A continuous character of this profinite group into the circle has
finite image: its image is a compact quotient of a profinite group, hence
totally disconnected, whereas an infinite closed subgroup of the circle is
the circle.  It therefore factors through some
$(\mathbb Z/N\mathbb Z)^\times$.  The remaining unitary character of
$\mathbb R_{>0}$ is $x^{it}$.  In the convention
$L(s,\chi_0|\cdot|^{it})=L(s+it,\chi_0)$, zeros are translated only in
their imaginary coordinate.  This proves the assertion.  Geometrically,
the finite-order part is a rank-one representation of the arithmetic
fundamental group; the continuous norm twist is an analytic normalization,
not an additional finite-monodromy motive.  The result does not cover
arbitrary characters over a nontrivial number field.

\subsection*{MT-M01. A nonzero counterexample to the draft's general zero-count lemma}
Masullo, Lemma 9.3.2.4, PDF p.~425, claims a local $O(\log |T|)$ zero count
for the Mellin transform of every $\varphi\in\mathcal S_E(G)$.
We checked the image of the original PDF page, including the displayed
Jensen inequality and its last inference.  The quantifier is present in the
PDF; this is not an inference from the model transcription.

Take $G=\mathbb R_{>0}$, trivial character, Haar measure $dt/t$, and set
\[
 g_a(u)=\frac{1}{\sqrt{4\pi a}}\exp\left(-\frac{u^2}{4a}\right),
 \qquad a=1+i,\qquad
 \varphi(e^u)=\frac{g_{1+i}(u)-g_{1-i}(u)}{2i}.
\]
The square roots use the branch with positive real part.  This is a real,
nonzero smooth function.  For every real $b$ and every nonnegative integer
$j$, the $j$th derivative of $e^{bu}g_{1\pm i}(u)$ is a polynomial in $u$
times $e^{bu-u^2/(4(1\pm i))}$.  Its absolute value is bounded by a
polynomial times $e^{bu-u^2/8}$, hence is rapidly decreasing.  Thus
$t^b\varphi(t)$ is Schwartz in the logarithmic coordinate for every $b$;
this is exactly the condition of Definition D.3.3.1, PDF p.~804 (also
\S1.1.2, PDF p.~11).  In particular $\varphi\in\mathcal S_E(G)$.

For $\operatorname{Re}a>0$, the Gaussian integral and analytic continuation
in $a$ give
\[
 \int_{\mathbb R}g_a(u)e^{su}\,du=e^{as^2},\qquad s\in\mathbb C.
\]
For completeness, the identity first follows by completing the square for
positive real $a$ and real $s$; both sides are entire in $s$ by Gaussian
domination on compact sets, and holomorphic in $a$ on
$\operatorname{Re}a>0$, where $\operatorname{Re}(1/a)>0$.  The identity
theorem proves the claimed continuation.  Therefore
\[
 \mathcal M\varphi(s)=e^{s^2}\sin(s^2).
\]
Its zeros on the positive imaginary axis are
$i\sqrt{k\pi}$, $k=1,2,\ldots$, and are simple.  For $I=[-1,1]$ and
$T\ge2$, their number in $I+i[T,T+1]$ is
\[
 \#\{k:T^2/\pi\le k\le(T+1)^2/\pi\}
   =\frac{2T+1}{\pi}+O(1).
\]
This is not $O(\log T)$.  The function is nonzero, so removing the zero
function from the lemma does not repair it.  The proof on p.~425 needs a
\emph{lower} bound for the modulus at a Jensen centre, whereas rapid decay
supplies upper bounds.  No statement about the zeros of a Hecke L-function
is contradicted by this example.

\subsection*{MT-M02. The displayed approximate identities are proper kernel projections}
Theorem 9.3.2.6, equation (9.3.2.10), PDF p.~426, defines subspaces by
intersecting kernels of all jet evaluations up to height $T$, then claims
their orthogonal projections tend strongly to the identity.  The printed
formula and the asserted inclusion in its proof have both been checked on
the page image.

Here is a counterexample even with a single zero, so the zero-count issue
in MT-M01 is irrelevant.  Put
\[
 g(u)=(4\pi)^{-1/2}e^{-u^2/4},\qquad
 \psi(e^u)=-g'(u).
\]
Then $\psi\in\mathcal S_E(\mathbb R_{>0})$, and integration by parts gives
$\mathcal M\psi(z)=z e^{z^2}$.  Its only zero is $0$, with multiplicity one,
strictly inside the strip for $I=[-1,1]$.  Choose Sobolev index $0$.
The draft's Hilbert norm is then
\[
 \|f\|_H^2=\int_{\mathbb R}|f(e^u)|^2e^{2|u|}\,du.
\]
The functional $\delta_0(f)=\int_{\mathbb R}f(e^u)\,du$ is continuous,
because Cauchy--Schwarz bounds its modulus by
$\|f\|_H(\int e^{-2|u|}\,du)^{1/2}=\|f\|_H$.
It is nonzero, since $\delta_0(g)=1$.  For every $T>0$ the subspace displayed
in (9.3.2.10) is the same proper closed subspace $\ker\delta_0$.
Its orthogonal projector cannot converge to the identity.  In particular,
the Riesz vector $r(e^u)=e^{-2|u|}$ has $\|r\|_H=1$ and is killed by every
one of these projectors.  It belongs to the indicated completion because
$C_c^\infty(\mathbb R)$ is contained in $\mathcal S_E$ in logarithmic
coordinates and is dense in this weighted $L^2$ space: truncate first and
then smooth on compact intervals, where the positive weight is bounded
above and below.  The ranges decrease when more evaluation
conditions are added, whereas the proof on p.~426 uses the opposite
inclusion.

\subsection*{MT-M03. The exact projection repair}
Let $H$ be a Hilbert space, let $(\ell_j)_{j\ge1}$ be continuous linear
functionals, and put
\[
 V_N=\bigcap_{j\le N}\ker\ell_j,\qquad V=\bigcap_{j\ge1}\ker\ell_j.
\]
Then $P_{V_N}\to P_V$ strongly.  The complementary projectors converge
strongly to $P_{V^\perp}$, the projector onto the closed span of the Riesz
vectors of the $\ell_j$.  Thus they give an approximate identity on the
spectral quotient $H/V\simeq V^\perp$.

Here is a full proof.  For fixed $x\in H$, put $y_N=P_{V_N}x$.  For
$M\ge N$, nestedness gives $P_{V_M}y_N=y_M$ and hence
\[
 \|y_N-y_M\|^2=\|y_N\|^2-\|y_M\|^2.
\]
The nonincreasing nonnegative sequence $\|y_N\|^2$ converges, so $(y_N)$ is
Cauchy.  Write its limit as $y$.  For each fixed $N$, all sufficiently late
$y_M$ lie in the closed space $V_N$, so $y\in V$.  Also
$x-y_N\perp V$ for every $N$; passing to the limit gives $x-y\perp V$.
Therefore $y=P_Vx$.  Taking complements proves the second assertion.
Finally the orthogonal complement of the closed span of the Riesz vectors
is precisely $V$, which proves the last identification.  Exhaustion by
height instead of an integer enumeration gives the same proof.

This repairs a functional-analytic direction error.  It does not identify
the common kernel with the closed adelic theta image, prove a Schauder
basis, establish nuclearity of the actual convolution operators, or justify
the geometric-kernel decomposition.  Those are separate statements and
must be checked in the quotient on which they are used.  In Grothendieck's
trace language this is the distinction between an annihilator and its
spectral quotient, before any trace or Fredholm determinant is taken.

\subsection*{MT-M04. An arithmetic replacement for the false general lemma}
For every fixed unitary Hecke character, the number of nontrivial zeros
in a unit-height window is $O_{K,\eta}(\log(|T|+3))$, unconditionally.
The full proof for arbitrary infinity type is NT-Z10 above: polynomial
vertical growth and a uniform lower bound at the Euler-product centre
give Jensen's estimate.  It includes pure norm twists, whose single
translated pole is excluded before applying Jensen at large height.
For a fixed primitive finite-order character a conductor-explicit
primary locator is Greni\'e--Molteni,
\emph{An explicit Chebotarev density theorem under GRH}, Appendix A,
Theorem A.1, author PDF p.~26: that appendix's zero count is unconditional.
Subtract the counts at consecutive heights.  The main term has derivative
$O(\log(|T|+3))$ for a fixed field and character and each remainder has the
same bound.  The arithmetic restriction therefore closes the local
zero-count gate, including infinite-order characters, without using GRH.
It supplies a route around MT-M01 for arithmetic divisors, but does not
establish Lemma 9.3.2.5 for arbitrary Schwartz convolution ideals.  The
constant in the infinite-order statement is fixed-character, not a
claimed uniform archimedean-conductor constant.

\subsection*{MT-ART01. A nonmonomial GL\texorpdfstring{$_2$}{2} consequence: odd Artin forms}
Assume $M_H$.  Let $\rho:G_{\mathbb Q}\to\operatorname{GL}_2(\mathbb C)$
be continuous, irreducible, odd and of finite image.  Then its Artin
L-function is entire and all of its nontrivial zeros lie on
$\operatorname{Re}s=1/2$.

The holomorphy input is already a theorem: Khare--Wintenberger,
\emph{Serre's modularity conjecture (I)}, Corollary 10.2(ii), author PDF
p.~21 (with Theorem 10.1(ii), pp.~20--21), associates $\rho$ to a weight-one
cuspidal newform.  Its standard L-function is entire.  Let $L$ be the finite
Galois extension cut out by $\rho$.  Brauer induction writes its character
as an integer linear combination of characters induced from
one-dimensional characters of subgroups of $\operatorname{Gal}(L/\mathbb Q)$.
Artin induction and class field theory therefore express $L(s,\rho)$ as
a finite product of integer powers of finite-order Hecke L-functions over
the corresponding fixed fields.  In the open critical strip, $M_H$ makes
each factor holomorphic and nonzero away from the critical line.  The
product, including its negative exponents, is consequently holomorphic and
nonzero there.  Combining this with the independent holomorphy theorem
proves GRH for $L(s,\rho)$.

This includes projective image $A_5$: no monomial expression with only
nonnegative exponents is required.  It is a new transfer in this catalogue,
not a new proof of the classical modularity theorem.  In motivic terms it
concerns odd rank-two Artin motives and their weight-one automorphic
realizations.  It gives no GRH for the weight-two L-function of an arbitrary
non-CM elliptic curve.  Nor does it prove holomorphy of arbitrary even
icosahedral or higher-dimensional Artin representations.  The inherited
Brauer divisor-support statements in RETURN\_ZETA\_TWO are the general
input behind this specialization.

\subsection*{MT-ART02. The missing positivity is a finite matrix at each zero}
Let $L/K$ be a finite Galois extension with group $G$, and let
$0<\operatorname{Re}s_0<1$.  By Brauer meromorphy define
\[
 n_\chi=\operatorname{ord}_{s=s_0}L(s,\chi,L/K)\in\mathbb Z,
 \qquad\Theta_{s_0}=\sum_{\chi\in\operatorname{Irr}(G)}n_\chi\chi.
\]
Form the $|G|$ by $|G|$ matrix
\[
 K_{g,h}=\Theta_{s_0}(g h^{-1}).
\]
Then all Artin L-functions for $L/K$ are holomorphic at $s_0$ if and only
if $K$ is positive semidefinite.  The matrix is Hermitian because
$\chi(g^{-1})=\overline{\chi(g)}$ and $n_\chi$ is real.

To prove the equivalence, consider the central convolution operator
$T=|G|^{-1}K$ on $\mathbb C[G]$.  On the matrix-coefficient space of an
irreducible character $\chi$, Schur orthogonality makes $T$ scalar with
eigenvalue $n_\chi/\chi(1)$ (interchanging $\chi$ and its dual under the
opposite convolution convention only permutes the list).  Each such space
has dimension $\chi(1)^2$ and together they exhaust the regular
representation.  Thus $K\succeq0$ is equivalent to $n_\chi\ge0$ for every
$\chi$, which is exactly absence of a pole of every factor.

The restrictions of $\Theta_{s_0}$ to cyclic subgroups are already actual
characters.  Indeed, for a cyclic subgroup $C$ and a one-dimensional
character $\eta$ of $C$, Frobenius reciprocity and Artin induction give
\[
 \langle\Theta_{s_0}|_C,\eta\rangle_C
   =\operatorname{ord}_{s=s_0}L(s,\eta,L/L^C)\ge0,
\]
because the right-hand side is an order of a holomorphic Hecke function in
the strip.  This proves positivity on every such restriction; it does not
prove the full matrix is positive.  Also
\[
 \Theta_{s_0}(1)=\sum_\chi n_\chi\chi(1)
   =\operatorname{ord}_{s=s_0}\zeta_L(s).
\]
The point of the criterion is the exact remaining object: a central
positive kernel in the Artin character ring.  $M_H$ removes all tests off
the critical line, while neither zero location nor cyclic positivity
supplies the full critical-line kernel.  An Alp\"oge-type positivity
argument would have to produce this kernel or an equivalent positive
decomposition, with the actual Artin orders as coefficients.

\subsection*{MT-ART04. All induced-linear positivity tests still leave a gap}
There is a virtual character which passes every positivity inequality
obtained by inducing a one-dimensional character of a subgroup, but is not
an actual character.  This refutes an abstract reconstruction rule; we do
not claim that this virtual character occurs as the orders of actual
Artin L-functions at a point.

Let $G=A_5$, and let $\chi$ be its four-dimensional augmentation character
for the natural action on five letters.  Define
\[
             \Theta=\operatorname{reg}_G-5\chi.
\]
The natural permutation action is doubly transitive, so its permutation
character has inner-product norm squared two.  Its invariant subspace is
one-dimensional; the orthogonal complement therefore has character
$\chi$ with norm squared one, proving that $\chi$ is irreducible.
The coefficient of $\chi$ in $\Theta$ is $4-5=-1$, so $\Theta$ is not an
actual character and its full matrix from MT-ART02 is not positive.

Nevertheless, for every $H\le G$ and every linear character $\eta$ of $H$,
\[
       \left\langle\Theta,\operatorname{Ind}_H^G\eta\right\rangle_G
                    \ge0.                                      \tag{MT.1}
\]
Here are all the details.  The conjugacy classes of $A_5$ have sizes
$1,15,20,12,12$, as follows by counting the double transpositions,
three-cycles and the two five-cycle classes.  A normal subgroup is a union
of these classes containing the identity, and its size divides $60$.
No proper nontrivial sum of these class sizes containing $1$ divides $60$;
thus $A_5$ is simple.  An action on cosets of a proper subgroup of index
$2,3,$ or $4$ would be an injective homomorphism into the corresponding
symmetric group, which is impossible by order.  Lagrange's theorem leaves
only the possible proper indices
\[
                    k\in\{5,6,10,12,15,20,30,60\}.
\]
Put $m=\langle\chi|_H,\eta\rangle_H$.  Since $\chi|_H$ has dimension four,
$0\le m\le4$.  Frobenius reciprocity says that
$\operatorname{Ind}_H^G\eta$, of dimension $k$, contains $m$ copies of
$\chi$, so $4m\le k$.  These two bounds imply $k-5m\ge0$ for every index
in the list except possibly $k=12$.  At $k=12$, $|H|=5$, and $H$ acts
regularly on the five letters.  The restriction of $\chi$ is the sum of
the four distinct nontrivial characters of $C_5$, so $m\le1$ and again
$k-5m\ge0$.  Finally
\[
 \left\langle\operatorname{reg}_G,
       \operatorname{Ind}_H^G\eta\right\rangle=k,
\]
proving (MT.1) for proper $H$.  For $H=G$, simplicity and nonabelianity
make the trivial character the only linear character, and its coefficient
in $\Theta$ is one.  This proves (MT.1) in all cases.

For comparison with the methods, Artin induction and Hecke holomorphy
give exactly inequalities of the form (MT.1) for an actual local Artin
order character.  The example proves that these inequalities alone do not
force a positive character, even before any question of zero location.
Thus a transfer of determinant positivity from the Alp\"oge method would
need additional arithmetic information, not just all the existing
induced-linear nonnegative orders.  The elementary exact computation
\path{checks/artin_cone_certificate.py} checks the negative matrix
quadratic form: with $v(g)=\chi(g)$ one has $v^*Kv=-900$.

\subsection*{MT-ART03. The classical simple-zero reduction}
Stark, \emph{Some effective cases of the Brauer--Siegel theorem},
Inventiones Mathematicae \textbf{23} (1974), Theorem 3, p.~144, proves
that if $\operatorname{ord}_{s_0}\zeta_L\le1$ and $s_0\ne1$, then all
Artin factors for $L/K$ are holomorphic at $s_0$.  We cite this classical
result rather than count another proof as a finding.  Applied with the
divisor-support theorem under $M_H$, it confines any still-possible pole
inside the strip to a multiple zero of the top-field zeta on the critical
line.  No simplicity assertion for these Dedekind zetas is being made.

\subsection*{MT-Q01. Exact path quadrics in every characteristic}
For $k\ge2$ let
\[
 V_k=\left\{(x_1,\ldots,x_k):\sum_{i=1}^{k-1}x_ix_{i+1}=0\right\}
       \subset\mathbb A^k_{\mathbb Z},\qquad
 r=\lfloor k/2\rfloor,\quad d=\lceil k/2\rceil.
\]
Then for every finite field $\mathbb F_q$, including characteristic two,
\[
 \#V_k(\mathbb F_q)=q^{k-1}+(q-1)q^{d-1}.                 \tag{MT.2}
\]
In the relative Grothendieck ring of varieties over $\mathbb Z$ one has
\[
                 [V_k]=\mathbb L^{k-1}+\mathbb L^d-\mathbb L^{d-1}.
                                                               \tag{MT.3}
\]
Here $\mathbb L=[\mathbb A^1_{\mathbb Z}]$; the identity is an identity of
classes, not a claim of a direct-sum splitting of all integral realizations.

To prove these formulas, group the odd coordinates into $x\in\mathbb A^d$
and the even coordinates into $y\in\mathbb A^r$.  The quadratic form
becomes $x^{\mathsf T}Ay$, where the integral $d$ by $r$ matrix $A$ has an
$r$ by $r$ triangular minor with diagonal entries one.  It has a left
inverse over $\mathbb Z$ and its columns extend to a unimodular basis of
$\mathbb Z^d$.  Thus after an invertible integral change of the $x$
coordinates the equation is $\sum_{j=1}^r x_jy_j=0$, with $d-r$ free
$x$ coordinates.  Over $y=0$ its fibre is $\mathbb A^d$.  Over
$\mathbb A^r\setminus\{0\}$, the map of vector bundles
$\mathcal O^d\to\mathcal O$ given by $(y_1,\ldots,y_r,0,\ldots,0)$ is
surjective, and its kernel is a vector bundle of rank $d-1$.
It is Zariski locally trivial on the opens where a coordinate $y_j$ is
invertible.  Additivity and the vector-bundle relation give
\[
 [V_k]=\mathbb L^d+(\mathbb L^r-1)\mathbb L^{d-1},
\]
which is (MT.3).  Counting the same fibres over $\mathbb F_q$ proves
(MT.2) directly, without a purity theorem.

Consequently the exact error term in the unweighted orthogonality identity
of the supplied Fable sketch is
\[
             q\,\#V_k(\mathbb F_q)=q^k+(q-1)q^d.
\]
For a fixed nonzero additive-character parameter its quadratic sum is
$q^d$, by summing first over $x$.  This gives square-root size when $k$ is
even; for odd $k$ the extra factor $q^{1/2}$ reflects the one-dimensional
radical.  It explains the sign and the exact constant for this family.
The general weighted operator estimate in
\path{OURS_FABLE_PROGRAM_SUGGESTIONS_20261120_EXPLAINED.md},
\S3, remains a separate inherited result.  This calculation does not prove
the sign or norm assertions for arbitrary sheaves in Alp\"oge's steps 3
and 5, nor identify this elementary model with all the constructions in
the omitted Lean dependency tree.

\subsection*{MT-Q02. Its arithmetic zeta function and shifted zero lines}
The exact all-prime formula in MT-Q01 gives
\[
 Z(V_k/\mathbb F_q,T)
     =\frac{1-q^{d-1}T}{(1-q^{k-1}T)(1-q^dT)},
 \qquad
 \zeta_{V_k}(z)
     =\frac{\zeta(z-k+1)\zeta(z-d)}{\zeta(z-d+1)}.       \tag{MT.4}
\]
Indeed apply (MT.2) over every $\mathbb F_{q^m}$ in
$Z=\exp(\sum_{m\ge1}\#V_k(\mathbb F_{q^m})T^m/m)$ and use
$-\log(1-aT)=\sum_{m\ge1}a^mT^m/m$.  Multiplying the resulting factors
over rational primes proves the second identity first in a right
half-plane of absolute convergence, and then meromorphically.
Under $A_D$ the nontrivial divisor of this function is supported on
\[
    \operatorname{Re}z=k-\tfrac12,\qquad
    \operatorname{Re}z=d+\tfrac12,\qquad
    \operatorname{Re}z=d-\tfrac12.
\]
Divisor orders are signed and cancellations are kept exactly as in
(MT.4); the denominator's zeros are possible poles, not additional zeros.
The shifted pole of $\zeta$ at $1$ and shifted trivial zeros are recorded
by the same explicit formula, separately from this nontrivial-divisor
assertion.  This is an all-characteristic, concrete Artin--Tate example
linking the path model of the method with arithmetic zeta functions and
the Grothendieck ring.  It illustrates why a mixed-weight scheme has
several shifted lines rather than a single universal critical line.

\subsection*{MT-L01. Damped Euler products and Adams operations}
Let $\rho$ be a finite-image complex representation of $G_K$ with no
invariant vectors, of dimension $d$.  At each finite place $v$ put
$A_v=\rho(\operatorname{Frob}_v)|_{V^{I_v}}$.  Its eigenvalues have
absolute value one.  For $|c|\le1$ define, initially for $\operatorname{Re}s>1$,
\[
 D_c(s,\rho)=\prod_v\det(1-cA_v(Nv)^{-s})^{-1}.
\]
If $L(s,\rho)$ is holomorphic and nonzero on
$H=\{\operatorname{Re}s>1/2\}$, then $D_c$ has a holomorphic nonvanishing
continuation to $H$, given by
\begin{align*}
 D_c(s,\rho)&=\exp\{c\log L(s,\rho)+E_c(s)\},\\
 E_c(s)&=\sum_v\sum_{k\ge2}
   \frac{c^k-c}{k}\operatorname{tr}(A_v^k)(Nv)^{-ks}.
\end{align*}
The branch of $\log L$ is fixed by the Euler series in
$\operatorname{Re}s>1$.

For a full proof, expand $-\log\det(1-X)$ in powers of $X$ in the initial
half-plane.  The terms with $k=1$ cancel in
$\log D_c-c\log L$, giving the displayed $E_c$.  On a compact subset with
$\operatorname{Re}s\ge1/2+\varepsilon$, its terms in absolute value sum to
at most
\[
 2d\sum_v\sum_{k\ge2}(Nv)^{-k(1/2+\varepsilon)}
 \le \frac{2d}{1-2^{-1/2-\varepsilon}}
       \sum_v(Nv)^{-1-2\varepsilon}<\infty.
\]
Normal convergence, also uniform in $|c|\le1$, proves that $E_c$ is
holomorphic in $s$.  A nonvanishing holomorphic function on the simply
connected half-plane $H$ has a holomorphic logarithm, uniquely normalized
by its existing branch on the right.  The exponential above is therefore
holomorphic, nonzero, and agrees with the product where it converges.
This proves the continuation.

Under $M_H$, Brauer induction supplies the needed holomorphic nonvanishing
on $H$ for $\rho$ with no invariants: all Hecke factors have no zeros
there, and their orders at $1$ combine to
$-\dim V^{G_K}=0$.  This does not require Artin holomorphy on the boundary
line.  Under $A_D$ the same argument applies to sums of nonprincipal
Dirichlet characters over $\mathbb Q$.

The traces $\operatorname{tr}(A_v^k)$ are the local Adams-power characters;
the proof separates the first prime term from all higher Adams terms,
whose series converge in the larger half-plane.  This is a concrete
$\lambda$-ring link, rather than a claim that positivity of coefficients
follows from the formal operations.  For $c=1/m$ the phenomenon is closely
related to Gun--Hazra--Sahu, \emph{On holomorphy and non-vanishing of Artin
L-functions} (2025), Theorem 5.  The displayed formula covers all complex
$|c|\le1$ and makes its hypotheses explicit.

\subsection*{MT-L02. One irrational parameter detects the zero-free half-plane}
Fix an irrational real number $0<c<1$ and a nontrivial irreducible Artin
representation $\rho$.  Then the following are equivalent:
\begin{enumerate}
 \item $L(s,\rho)$ is holomorphic and nonzero on $H$;
 \item the Euler product $D_c(s,\rho)$ admits a single-valued meromorphic
       continuation to $H$.
\end{enumerate}
In the second assertion meromorphic is a global, single-valued condition;
a branch on a slit domain is not enough.

MT-L01 proves the forward implication.  Conversely $L$ is meromorphic by
Brauer induction and $E_c$ is holomorphic on $H$ by the same normally
convergent series without any hypothesis about the zeros of $L$.
Logarithmic differentiation in $\operatorname{Re}s>1$ gives
\[
 \frac{D_c'}{D_c}=c\frac{L'}L+E_c'.
\]
Both sides are meromorphic on connected $H$ if the asserted continuation
of $D_c$ exists, so the identity holds on $H$.  Taking residues at any
$s_0\in H$ gives
\[
 \operatorname{ord}_{s_0}D_c
      =c\operatorname{ord}_{s_0}L.
\]
Both orders are integers, and $c$ is irrational.  Both must be zero.
Thus $L$ is holomorphic and nonvanishing on $H$, proving the equivalence.
If instead $c=a/b$ in lowest terms, the same equality proves
$b\mid\operatorname{ord}_{s_0}L$ everywhere in $H$.

GRH in the meromorphic-divisor sense for the nontrivial irreducible Artin
factors is therefore equivalent, by the functional equation, to these
irrational-parameter continuation tests.  To include the trivial factor one
must separately add RH for $\zeta_K$; that input is already part of $M_H$.
This is an analytic multiplicity test.  It neither constructs the
continuation from arithmetic lattices nor removes poles on the critical
line.  With a trivial summand of multiplicity $r$, the known pole at $1$
contributes local monodromy $e^{-2\pi i c r}$; this is why the no-invariants
hypothesis was retained rather than suppressed.

% END root
\part{Unchecked claims and unresolved gates}
The complete proofs above establish only their stated implications.
The rows below are retained as research gates or as carried claims whose
complete source proofs were unavailable. They are not additional theorems.
The manuscript faults MT-M01 and MT-M02 have proved counterexamples;
the associated repairs do not certify every downstream operator or trace.

\paragraph{NT-Q06} Which pair-correlation and low-zero statistics are actually established?

\textbf{State:} proved/\allowbreak{}classical NT-Z01–\allowbreak{}NT-Z05 with support/\allowbreak{}moment restrictions; full GUE/\allowbreak{}correlations remain gate NT-Z06

\textbf{Source:} Montgomery1973 §1; HR2003 Thm3.\allowbreak{}1/\allowbreak{}5.\allowbreak{}1.

\paragraph{NT-Q07} Does rank-one Hecke GRH discharge the GL2 family hypothesis?

\textbf{State:} no such implication established; exact additional GL2-GRH gate NT-Z07 retained

\textbf{Source:} ILS2000 p56 and Thm1.\allowbreak{}1.

\paragraph{NT-Q10} Can carried class-group generator theorem be freshly source-inspected?

\textbf{State:} current gate:\allowbreak{} original AMS PDF access; inherited CX-N07 statement retained honestly

\textbf{Source:} Bach1990 Thm4 pp375–\allowbreak{}376.

\paragraph{AG-Q01} Which finite-type Z-scheme zeta functions have the adopted zero-support conclusion?

\textbf{State:} proved inherited exact Artin–\allowbreak{}Tate transfers; M adds identified potentiallyCM factors; residual nonCM continuation/\allowbreak{}RH and exceptional fibres remain explicit gates

\textbf{Source:} BRIDGES Q01; ZETA\_\allowbreak{}TWO ZT01–\allowbreak{}ZT07.

\paragraph{AG-Q08} Can the positive squared-variable determinant extend beyond \ensuremath{\zeta}?

\textbf{State:} proved AG-DET/\allowbreak{}MOMENT/\allowbreak{}HANKEL and AG-SIGNED finite pole witness; prime-defined cohomological realization and moment positivity for arbitrary Artin quotients remain gates

\textbf{Source:} ALPOGE\_\allowbreak{}BRIDGES BR-TR-01/\allowbreak{}02; SA23 geo-H-lidskii/\allowbreak{}geo-H-genus; cote1/\allowbreak{}79–\allowbreak{}80,111–\allowbreak{}112.

\paragraph{AG-Q11} Does finite-field purity imply openness of closure[pi0,H] in pi-prime under the original hypotheses?

\textbf{State:} original notation recovered:\allowbreak{} literal openness in pi-prime refuted AG-C34-TYPE; intrinsic closure[P,H] openH proved AG-C34-REPAIR from archive Thm1.\allowbreak{}1; exact index AG-C34-INDEX; section-dependent alternative meanings separately gated

\textbf{Source:} BRIDGES Q03; cote34/\allowbreak{}64 Rem1.\allowbreak{}17; definitions53–\allowbreak{}63.

\paragraph{AG-Q13} Do zero-location inputs construct integral Lefschetz correspondences?

\textbf{State:} current gate:\allowbreak{} cycle/\allowbreak{}rank comparison on X×X plus denominator control; no such construction follows from displayed analytic reductions

\textbf{Source:} Hodge sketch§§3.\allowbreak{}2–\allowbreak{}3.\allowbreak{}3; cote16/\allowbreak{}14–\allowbreak{}18 no7,44–\allowbreak{}50.

\paragraph{AG-Q14} Does normalized L-sign prove arbitrary-cycle arithmetic Hodge index?

\textbf{State:} current gate:\allowbreak{} exact height identity for the given cycles and covered L-function; conditional sign deduction AG-HEIGHT proved

\textbf{Source:} Hodge sketch§3.\allowbreak{}5; cote44 Theorems4–\allowbreak{}5.

\paragraph{AG-Q16} Which finite/\allowbreak{}profinite GT gates move?

\textbf{State:} finite quotient detection/\allowbreak{}generation effective AG-FROB,CX-GEO-004/\allowbreak{}008; full profinite kernel/\allowbreak{}associator gate remains

\textbf{Source:} ALPOGE\_\allowbreak{}BRIDGES R-GTHAT; SA22 GT lines.

\paragraph{AG-Q17} Does global Hecke GRH give general weight-monodromy?

\textbf{State:} current gate:\allowbreak{} all-power trace bounds for finite exterior constructions and duals; prior reduction retained

\textbf{Source:} ALPOGE\_\allowbreak{}BRIDGES D-WM,BR-WM-02,CC-WM01–\allowbreak{}03.

\paragraph{AG-Q18} What K-group statement is actually exported by A?

\textbf{State:} endpoint A does not itself export an indexed K-group theorem; inherited2026K8 claim not independently certified; integral comparison/\allowbreak{}class-group components remain gates

\textbf{Source:} BRIDGES Q06; Alpoge step17/\allowbreak{}fn6; ORDER31.

\paragraph{AG-Q23} What exact extra positivity would eliminate poles after Hecke/\allowbreak{}Brauer divisor localization?

\textbf{State:} proved equivalence with all finite signed Hankel PSD and an explicit finite witness for any negative multiplicity; proving these inequalities from independent arithmetic data remains open

\textbf{Source:} AG-SIGNED; rootMT-ART02; cote1Fredholm framework.

\paragraph{CX-Q04} What is the strongest available general finite-field factorization transfer?

\textbf{State:} proved conditional on M:\allowbreak{} CX-ALG-009; A\_\allowbreak{}D-only extension gate CX-ALG-010 (high valuation coefficient unit or alternative descent)

\textbf{Source:} CX\_\allowbreak{}GRH CX-N15/\allowbreak{}CX-U02.

\paragraph{CX-Q05} What does GRH give for class groups and what computation remains?

\textbf{State:} proved conditional generator transfer CX-ALG-011/\allowbreak{}013; full polynomial relation-search bound remains gate CX-ALG-012

\textbf{Source:} Bach1990T4; CX\_\allowbreak{}GRH CX-N07/\allowbreak{}N08.

\paragraph{CX-Q15} Do the cited GRH transfers close general derandomization or GCT hardness gates?

\textbf{State:} current gate:\allowbreak{} CX-LIM-001, uniform hitting sets or independent border-hardness theorem; no logical nonimplication claimed

\textbf{Source:} GCT V T4.\allowbreak{}7/\allowbreak{}L4.\allowbreak{}8; PIT/\allowbreak{}BPP questions.

\paragraph{CX-Q16} Does positivity/\allowbreak{}purity remove efficient multiplicity-construction gates?

\textbf{State:} specified earlier counterexamples carried CX-LIM-002; actual quotient ranks and uniform data construction remain

\textbf{Source:} ALPOGE\_\allowbreak{}GCT P-F7/\allowbreak{}P-F10 and PLETH PLE-F5/\allowbreak{}F7.

\paragraph{Q-ARTIN-POS} Can positivity remove critical-line Artin poles?

\textbf{State:} current gate MT-ART02:\allowbreak{} positive semidefinite local Heilbronn matrix at multiple zeta zeros; no proof of that positivity

\textbf{Source:} ORDER42; Alpoge steps4,7.
\section*{Statement index}
\addcontentsline{toc}{section}{Statement index}
The complete statement, exact input, proof location and work placement for
every ID below is recorded in \path{CONSEQUENCES.tsv}. This index keeps
the paper navigable without duplicating the entire machine-readable table.
\paragraph{number theory} NT-P01, NT-P02, NT-P03, NT-P04, NT-P05, NT-P06, NT-P07, NT-P08, NT-P09, NT-P10, NT-C01, NT-C02, NT-C03, NT-C04, NT-C05, NT-H01, NT-H02, NT-H03, NT-H04, NT-A01, NT-A02, NT-A03, NT-F01, NT-F02, NT-F03, NT-F04, NT-L01, NT-L02, NT-L03, NT-L04, NT-L05, NT-L06, NT-L07, NT-L08, NT-Z01, NT-Z02, NT-Z03, NT-Z04, NT-Z05, NT-Z06, NT-Z07, NT-Z08, NT-Z09, NT-L09, NT-L10, NT-Z10, NT-P11, NT-P12, NT-Z11, NT-Z12.
\paragraph{geometry programme} AG-AT, AG-ARTIN, AG-DED, AG-LOCAL, AG-MODEL, AG-CM, AG-WEIGHT, AG-SCHEMES, AG-NONCM, AG-NONCM-W0, AG-FROB, AG-LAMBDA, AG-SIGN, AG-HEIGHT, AG-SIGN-NORM, AG-CRITICAL, AG-DET, AG-MOMENT, AG-HANKEL, AG-EXT, AG-EXT-K, AG-C1, AG-C16, AG-C34-74, AG-C34-64, AG-C44, AG-C149, AG-NODE, AG-INERTIA, AG-GT, AG-PURITY, AG-WM, AG-STANDARD, AG-COVER, AG-KUMMER, AG-KUMMER-2, AG-CM-REC, AG-CM-PERIOD, AG-SIGNED, AG-C34-TYPE, AG-C34-REPAIR, AG-C34-INDEX.
\paragraph{complexity} CX-ALG-001, CX-ALG-002, CX-ALG-003, CX-ALG-004, CX-ALG-005, CX-ALG-006, CX-ALG-007, CX-ALG-008, CX-ALG-009, CX-ALG-010, CX-ALG-011, CX-ALG-012, CX-ALG-013, CX-ALG-014, CX-HN-001, CX-HN-002, CX-HN-003, CX-HN-004, CX-GEO-001, CX-GEO-002, CX-GEO-003, CX-GEO-004, CX-GEO-005, CX-GEO-006, CX-NN-001, CX-NN-002, CX-NN-003, CX-NN-004, CX-LIM-001, CX-LIM-002, CX-LIM-003, CX-GEO-007, CX-GEO-008, CX-ALG-015, CX-ALG-016, CX-ALG-017, CX-ALG-018, CX-GEO-009, CX-ALG-019, CX-ALG-020, CX-ALG-021, CX-ALG-022.
\paragraph{root} MT-A01, MT-M01, MT-M02, MT-M03, MT-M04, MT-ART01, MT-ART02, MT-ART03, MT-L01, MT-L02, MT-ART04, MT-Q01, MT-Q02.
\section*{Source registry and exact locators}
\addcontentsline{toc}{section}{Source registry and exact locators}
The inspection descriptions below record their actual scope. A located
classical theorem is not labelled as a fresh whole-proof audit unless
its event explicitly records that reading.

\paragraph{S-NT-LLS} Lamzouri–\allowbreak{}Li–\allowbreak{}Soundararajan, Conditional bounds for the least quadratic non-residue and related problems, corrected arXiv v3 (8 Dec 2016).

\textbf{Locator:} Theorem1.\allowbreak{}1 and Corollary1.\allowbreak{}1 p2; Corollary1.\allowbreak{}2 p3; Theorem1.\allowbreak{}5 pp3–\allowbreak{}4; Corollary1.\allowbreak{}3 and (1.\allowbreak{}3) p4.\allowbreak{} Author-page numbering.\allowbreak{}.

\textbf{Inspection:} Primary PDF text read for full statements, thresholds, U/\allowbreak{}V normalization and correction-version choice.\allowbreak{} These results were already background IA02; no fresh reproof of the long literature argument.\allowbreak{}

\url{https://arxiv.org/pdf/1309.3595v3}.

\paragraph{S-NT-GM} Grenié–\allowbreak{}Molteni, An explicit Chebotarev density theorem under GRH, J.\allowbreak{}Number Theory200(2019),441–\allowbreak{}485.

\textbf{Locator:} Theorem1.\allowbreak{}1 author p2; Corollary1.\allowbreak{}2 author p3; AppendixA TheoremA.\allowbreak{}1 author p26.\allowbreak{}.

\textbf{Inspection:} Original author PDF pp2–\allowbreak{}3 rendered and visually inspected for both explicit error polynomials and unramified convention.\allowbreak{} AppendixA finite-order Hecke zero-count statement read; that appendix is unconditional.\allowbreak{} The local O(log T) deduction uses consecutive heights.\allowbreak{}

\url{https://sites.unimi.it/molteni/research/papers-pdf/44-molteni-An_explicit_Chebotarev_density_theorem_under_GRH.pdf}.

\paragraph{S-NT-BS} Bach–\allowbreak{}Sorenson, Explicit bounds for primes in residue classes, Math.\allowbreak{}Comp.\allowbreak{}65(1996),1717–\allowbreak{}1735.

\textbf{Locator:} §2 printed1720; Corollary3.\allowbreak{}3 printed1721; Theorem5.\allowbreak{}1 printed1729/\allowbreak{}PDF13.\allowbreak{}.

\textbf{Inspection:} Primary author-upload PDF opened in this agent.\allowbreak{} Geometry agent confirmed precise Thm5.\allowbreak{}1 statement; RH \ensuremath{\zeta}\_\allowbreak{}L and \ensuremath{\zeta}\_\allowbreak{}Q; degree-one prime; L\ensuremath{\ne}Q.\allowbreak{} Arbitrary extra-prime exclusion is not in this theorem.\allowbreak{}

\url{https://www.researchgate.net/profile/Jonathan-Sorenson-2/publication/220576951_Explicit_bounds_for_primes_in_residue_classes/links/55100ea70cf2ac2905afbceb/Explicit-bounds-for-primes-in-residue-classes.pdf}.

\paragraph{S-NT-SCHOENFELD} Schoenfeld, Sharper bounds for the Chebyshev functions \ensuremath{\theta}(x) and \ensuremath{\psi}(x).\allowbreak{}II, Math.\allowbreak{}Comp.\allowbreak{}30(1976),337–\allowbreak{}360.

\textbf{Locator:} Theorem10 equations(6.\allowbreak{}1)–\allowbreak{}(6.\allowbreak{}5), printed337; Corollary1 to Theorem10.\allowbreak{}.

\textbf{Inspection:} Primary indexed full statement of Theorem10 read, including 73.\allowbreak{}2/\allowbreak{}599 and 23·10\textasciicircum{}8.\allowbreak{} Direct AMS PDF403; mirror https:\allowbreak{}/\allowbreak{}/\allowbreak{}electronicsandbooks.\allowbreak{}com/\allowbreak{}edt/\allowbreak{}manual/\allowbreak{}Magazine/\allowbreak{}M/\allowbreak{}Mathematics\%20of\%20Computation/\allowbreak{}1960-2002/\allowbreak{}pdf/\allowbreak{}1976\_\allowbreak{}v030\_\allowbreak{}n134/\allowbreak{}2005976.\allowbreak{}pdf timed out via web.\allowbreak{} Corollary1 original page still to inspect; do not label full source proof checked.\allowbreak{}

\url{https://www.ams.org/mcom/1976-30-134/S0025-5718-1976-0457374-X/S0025-5718-1976-0457374-X.pdf}.

\paragraph{S-NT-DGM} Dudek–\allowbreak{}Grenié–\allowbreak{}Molteni, Primes in explicit short intervals on RH, Int.\allowbreak{}J.\allowbreak{}NumberTheory12(2016),1391–\allowbreak{}1407.

\textbf{Locator:} Theorem1.\allowbreak{}1 author p1.\allowbreak{}.

\textbf{Inspection:} Primary theorem statement, c=1/\allowbreak{}2+2/\allowbreak{}log x and x\ensuremath{\ge}2 read.\allowbreak{} Computational c=.\allowbreak{}55 refinement not used.\allowbreak{}

\url{https://sites.unimi.it/molteni/research/papers-pdf/36-molteni-Primes_in_explicit_short_intervals_on_RH.pdf}.

\paragraph{S-NT-GMP} Grenié–\allowbreak{}Molteni–\allowbreak{}Perelli, Primes and prime ideals in short intervals, arXiv:\allowbreak{}1602.\allowbreak{}02906v2.

\textbf{Locator:} Theorems1,2 and subsequent asymptotic statements, author p2; Proposition1 author p3.\allowbreak{}.

\textbf{Inspection:} Primary theorem statements and distinction Dirichlet-GRH versus RH \ensuremath{\zeta}\_\allowbreak{}K read.\allowbreak{} Absolute constants remain unspecified.\allowbreak{}

\url{https://arxiv.org/pdf/1602.02906}.

\paragraph{S-NT-SHOUP} Shoup, Searching for primitive roots in finite fields (1992).

\textbf{Locator:} Theorem3, author PDF p4.\allowbreak{}.

\textbf{Inspection:} Complexity agent downloaded original PDF with a standard browser user agent.\allowbreak{} This agent inspected the actual rendered p4 image at work/\allowbreak{}complexity/\allowbreak{}sources/\allowbreak{}Shoup\_\allowbreak{}Theorem3.\allowbreak{}png.\allowbreak{} Exact r\textasciicircum{}4(logr+1)\textasciicircum{}4log²p bound confirmed; no numerical implicit constant claimed.\allowbreak{}

\url{https://www.shoup.net/papers/primroots.pdf}.

\paragraph{S-NT-BACH} Bach, Explicit bounds for primality testing and related problems, Math.\allowbreak{}Comp.\allowbreak{}55(1990),355–\allowbreak{}380.

\textbf{Locator:} Theorem4 printed375 and corollary376, inherited exact locator.\allowbreak{}.

\textbf{Inspection:} Carried as background with no fault found; direct AMS requests inaccessible.\allowbreak{} Primary fresh inspection not claimed.\allowbreak{} Complexity agent owns the algorithmic application.\allowbreak{}

\url{https://www.ams.org/mcom/1990-55-191/S0025-5718-1990-1023756-8/S0025-5718-1990-1023756-8.pdf}.

\paragraph{S-NT-LENSTRA} Lenstra, On Artin’s conjecture and Euclid’s algorithm in global fields, Invent.\allowbreak{}Math.\allowbreak{}42(1977),201–\allowbreak{}224.

\textbf{Locator:} §2 pp202–\allowbreak{}206; (2.\allowbreak{}14),(2.\allowbreak{}15) pp205–\allowbreak{}206; Theorem(3.\allowbreak{}1)206; Theorem(4.\allowbreak{}1)207; Theorem(9.\allowbreak{}1),(9.\allowbreak{}2)218.\allowbreak{}.

\textbf{Inspection:} Primary definitions and theorem statements read; original p218 rendered and inspected, confirming \#S\ensuremath{\ge}2.\allowbreak{} Full elementary reduction in Theorem3.\allowbreak{}1/\allowbreak{}lemma3.\allowbreak{}2 read, but Cooke–\allowbreak{}Weinberger base theorem not independently reproved.\allowbreak{}

\url{https://pub.math.leidenuniv.nl/~lenstrahw/PUBLICATIONS/1977d/fulltext-12.pdf}.

\paragraph{S-NT-HOOLEY} Hooley, On Artin’s conjecture, J.\allowbreak{}reine angew.\allowbreak{}Math.\allowbreak{}225(1967),209–\allowbreak{}220.

\textbf{Locator:} Unnumbered final theorem pp219–\allowbreak{}220; original internal locator corroborated by later primary research.\allowbreak{}.

\textbf{Inspection:} Original publication metadata checked; original fulltext endpoint blocked.\allowbreak{} For the exact proof used here, cite inspected Lenstra Theorem3.\allowbreak{}1 instead.\allowbreak{} Original theorem internal-page check remains a source gate, not a mathematical gate on the classical corollary.\allowbreak{}

\url{https://doi.org/10.1515/crll.1967.225.209}.

\paragraph{S-NT-TITCHMARSH} Titchmarsh, The Theory of the Riemann Zeta-function, 2nd ed.\allowbreak{}, revised by Heath-Brown (1986).

\textbf{Locator:} Lemma3.\allowbreak{}12 and proof printed61–\allowbreak{}62; Theorem14.\allowbreak{}2 and (14.\allowbreak{}2.\allowbreak{}5),(14.\allowbreak{}2.\allowbreak{}6) printed336–\allowbreak{}337; Theorem14.\allowbreak{}14(A); Theorem14.\allowbreak{}25(A)–\allowbreak{}(C) printed369–\allowbreak{}370.\allowbreak{}.

\textbf{Inspection:} Downloaded local PDF and text.\allowbreak{} Relevant theorem statements and the whole printed 14.\allowbreak{}25 argument read; Perron kernel locator checked for the new transfer proof.\allowbreak{} Not a check of the entire textbook.\allowbreak{} Two printed pages per PDF page.\allowbreak{} Full reference text excluded from deliverables.\allowbreak{}

\url{https://sites.math.rutgers.edu/~zeilberg/EM18/TitchmarshZeta.pdf}.

\paragraph{S-NT-SIMONIC} Simonič, Estimates for L-functions in the Critical Strip Under GRH with Effective Applications, Mediterr.\allowbreak{}J.\allowbreak{}Math.\allowbreak{}20(2023),87.

\textbf{Locator:} Theorem1 equations(4)–\allowbreak{}(6); Corollaries2–\allowbreak{}3; §2.\allowbreak{}1 SP axioms; Theorem2 equations(7),(8); proof §3.\allowbreak{}.

\textbf{Inspection:} Primary full-text theorem and axioms read.\allowbreak{} Fixed-character consequences derived in full in section.\allowbreak{}tex; effective Mertens constants checked in original HTML.\allowbreak{} No uniform hybrid-conductor Lindelöf estimate claimed.\allowbreak{}

\url{https://link.springer.com/article/10.1007/s00009-023-02289-2}.

\paragraph{S-NT-HARPER} Harper, Sharp conditional bounds for moments of the Riemann zeta function, arXiv:\allowbreak{}1305.\allowbreak{}4618.

\textbf{Locator:} Theorem1 author p1.\allowbreak{}.

\textbf{Inspection:} Primary theorem statement read with k\ensuremath{\ge}0 fixed and implicit constant depending on k.\allowbreak{} Source introduction distinguishes order bounds from precise moments.\allowbreak{} No complete reproof of the paper.\allowbreak{}

\url{https://arxiv.org/pdf/1305.4618}.

\paragraph{S-NT-MONTGOMERY} Montgomery, The pair correlation of zeros of the zeta function, Proc.\allowbreak{}Sympos.\allowbreak{}PureMath.\allowbreak{}24(1973),181–\allowbreak{}193.

\textbf{Locator:} §1 definition and unnumbered Theorem p181; Corollary2 p182; conjecture following restricted theorem.\allowbreak{}.

\textbf{Inspection:} Primary web PDF text read for normalization, \ensuremath{\alpha} range, uniformity and simplicity proportion; direct local download403.\allowbreak{} Full pair-correlation conjecture carefully separated.\allowbreak{}

\url{https://public.websites.umich.edu/hlm/paircor1.pdf}.

\paragraph{S-NT-HR} Hughes–\allowbreak{}Rudnick, Linear statistics of low-lying zeros of L-functions, Q.\allowbreak{}J.\allowbreak{}Math.\allowbreak{}54(2003),309–\allowbreak{}333.

\textbf{Locator:} v2:\allowbreak{} ordering p1; admissibility §2; Theorem3.\allowbreak{}1p5; Theorem4.\allowbreak{}1p7; Theorem5.\allowbreak{}1p9; Corollary8.\allowbreak{}2p20; Theorem8.\allowbreak{}3pp20–\allowbreak{}21.\allowbreak{}.

\textbf{Inspection:} Primary text read for unconditional complex-zero formulation, Fourier support, m<2/\allowbreak{}\ensuremath{\alpha} condition, and first NONNEGATIVE zero convention.\allowbreak{} Source statements checked, not full moment proof recertified.\allowbreak{}

\url{https://arxiv.org/pdf/math/0208230}.

\paragraph{S-NT-ILS} Iwaniec–\allowbreak{}Luo–\allowbreak{}Sarnak, Low lying zeros of families of L-functions, Publ.\allowbreak{}Math.\allowbreak{}IHES91(2000),55–\allowbreak{}131.

\textbf{Locator:} Global RH assumption printed56; Theorem1.\allowbreak{}1 printed58–\allowbreak{}59; proof Theorem7.\allowbreak{}2.\allowbreak{}.

\textbf{Inspection:} Primary published PDF text inspected.\allowbreak{} The family GL2-RH hypothesis remains separate from auxiliary Dirichlet-GRH; not silently discharged by M.\allowbreak{}

\url{https://www.numdam.org/article/PMIHES_2000__91__55_0.pdf}.

\paragraph{S-NT-RT} Rodgers–\allowbreak{}Tao, The de Bruijn–\allowbreak{}Newman constant is non-negative, arXiv:\allowbreak{}1801.\allowbreak{}05914v5.

\textbf{Locator:} §1 definition/\allowbreak{}equivalence pp1–\allowbreak{}2; Theorem1p3.\allowbreak{}.

\textbf{Inspection:} Primary version5 dated3Jul2021; theorem and classical RH equivalence read.\allowbreak{} Combination with A gives Λ=0, with normalization retained.\allowbreak{}

\url{https://arxiv.org/pdf/1801.05914}.

\paragraph{S-NT-LAGARIAS} Lagarias, An elementary problem equivalent to the Riemann hypothesis, arXiv:\allowbreak{}math/\allowbreak{}0008177v2.

\textbf{Locator:} Theorem1.\allowbreak{}1 author p1.\allowbreak{}.

\textbf{Inspection:} Primary theorem statement read, including equality only at n=1; no reproduction of entire classical proof.\allowbreak{}

\url{https://arxiv.org/pdf/math/0008177}.

\paragraph{S-NT-MILNE} Milne, Class Field Theory v4.\allowbreak{}03 (2020).

\textbf{Locator:} ChapterVI Corollary2.\allowbreak{}12 and Example2.\allowbreak{}13, printed191/\allowbreak{}PDF200.\allowbreak{}.

\textbf{Inspection:} Primary formula read for residue, full regulator, real-quadratic normalization, roots-of-unity factor.\allowbreak{} Source is author-hosted course notes.\allowbreak{}

\url{https://www.jmilne.org/math/CourseNotes/CFT.pdf}.

\paragraph{S-NT-JOHNSTON} Johnston, Improving bounds on prime counting functions by partial verification of the Riemann hypothesis, RamanujanJ59(2022),1307–\allowbreak{}1321.

\textbf{Locator:} Theorem3.\allowbreak{}2 equations(3.\allowbreak{}3)–\allowbreak{}(3.\allowbreak{}6); §2 paragraph after(2.\allowbreak{}1), ordinary-count normalization.\allowbreak{}.

\textbf{Inspection:} Primary full-text theorem and normalization statement read.\allowbreak{} Under A the finite-height hypothesis can be met for every x, yielding \ensuremath{\psi} threshold59 and \ensuremath{\pi} threshold2657.\allowbreak{} This supplies an exact directly checked alternative proof locator to inaccessible Schoenfeld Corollary1.\allowbreak{}

\url{https://link.springer.com/article/10.1007/s11139-022-00616-x}.

\paragraph{S-NT-RV} Ramakrishnan–\allowbreak{}Valenza, Fourier Analysis on Number Fields, GTM186(1999).

\textbf{Locator:} §7.\allowbreak{}4 Theorem7-19 printed279; §7.\allowbreak{}7 Theorem7-28 printed289.\allowbreak{}.

\textbf{Inspection:} Original PDF downloaded and relevant full statements read directly in extracted text:\allowbreak{} completed unitary-Hecke meromorphy/\allowbreak{}functional equation and pure-norm pole exception; nonvanishing on Re s=1 for every unitary idele-class character.\allowbreak{} Web fetch exceeds size limit; local download succeeded.\allowbreak{} No page-image check or whole-book proof check claimed.\allowbreak{}

\url{https://webusers.imj-prg.fr/~julien.marche/Ref_GL_GT/Fourier_Analysis_on_Number_Fields.pdf}.

\paragraph{S-NT-FPS} Fiorilli–\allowbreak{}Parks–\allowbreak{}Södergren, Low-lying zeros of quadratic Dirichlet L-functions:\allowbreak{} Lower order terms for extended support, CompositioMath153(2017),1196–\allowbreak{}1216; arXiv1601.\allowbreak{}06833.

\textbf{Locator:} Definitions(1.\allowbreak{}1)–\allowbreak{}(1.\allowbreak{}2) and Theorem1.\allowbreak{}1 authorp2; Lemma2.\allowbreak{}5p7; Theorem3.\allowbreak{}5p12; final proof of Theorem1.\allowbreak{}1p14.\allowbreak{}.

\textbf{Inspection:} Primary PDF downloaded; exact family \ensuremath{\chi}\_\allowbreak{}\{8d\} for odd squarefree d of both signs, even nonnegative Schwartz weights, Fourier convention and support (\ensuremath{-}2,2) read.\allowbreak{} Checked Dirichlet-only GRH in the prime-sum step.\allowbreak{} Full literature proof not independently revalidated; the new nonvanishing deduction in section.\allowbreak{}tex is complete original prose.\allowbreak{}

\url{https://arxiv.org/pdf/1601.06833}.

\paragraph{S-AG-ZETA} Supplied ZETA\_\allowbreak{}TWO return.

\textbf{Locator:} RETURN\_\allowbreak{}ZETA\_\allowbreak{}TWO\_\allowbreak{}final\_\allowbreak{}IMPLICATIONS.\allowbreak{}md §§1–\allowbreak{}4, ZT01–\allowbreak{}ZT08, CM01, SV01–\allowbreak{}SV02; matching LEDGER and QUESTIONS.

\textbf{Inspection:} Full implication text read and relevant reductions rechecked algebraically; original earlier full proof file not included.\allowbreak{} Earlier findings retain their IDs and are not claimed new.\allowbreak{}


\paragraph{S-AG-BRIDGES} Supplied ALPOGE\_\allowbreak{}BRIDGES and original BRIDGES returns.

\textbf{Locator:} RETURN\_\allowbreak{}ALPOGE\_\allowbreak{}BRIDGES\_\allowbreak{}cp02\_\allowbreak{}IMPLICATIONS.\allowbreak{}md and QUESTIONS:\allowbreak{} BR-TR-01/\allowbreak{}02, BR-MIX-01, BR-WM-02, AQ-FINITE, R-GTHAT; OURS\_\allowbreak{}BRIDGES\_\allowbreak{}BRIDGES\_\allowbreak{}20261120.\allowbreak{}tex F1–\allowbreak{}F4,Q01–\allowbreak{}Q08.

\textbf{Inspection:} Implication/\allowbreak{}ledger/\allowbreak{}question records read; reused background is distinguished from the original proof package, which is not attached.\allowbreak{}


\paragraph{S-AG-BS} Bach–\allowbreak{}Sorenson, Explicit bounds for primes in residue classes, Math.\allowbreak{}Comp.\allowbreak{}65(1996),1717–\allowbreak{}1735.

\textbf{Locator:} Theorem5.\allowbreak{}1 printed1729/\allowbreak{}PDF13; §2 printed1720 (RH \ensuremath{\zeta}E,\ensuremath{\zeta}Q); Corollary3.\allowbreak{}3 printed1721.

\textbf{Inspection:} Original author-upload PDF theorem and hypothesis text read; no reproof of full classical explicit estimate.\allowbreak{}

\url{https://www.researchgate.net/profile/Jonathan-Sorenson-2/publication/220576951_Explicit_bounds_for_primes_in_residue_classes/links/55100ea70cf2ac2905afbceb/Explicit-bounds-for-primes-in-residue-classes.pdf}.

\paragraph{S-AG-DG} Deligne–\allowbreak{}Goncharov, Groupes fondamentaux motiviques de Tate mixte.

\textbf{Locator:} arXiv math/\allowbreak{}0302267v2,27Feb2003; (2.\allowbreak{}1.\allowbreak{}3), Propositions2.\allowbreak{}2–\allowbreak{}2.\allowbreak{}3, §2.\allowbreak{}4, printed14/\allowbreak{}PDF16.

\textbf{Inspection:} Primary formula and surrounding hypotheses read, distinguishing rational Ext and free Lie presentation from canonical generators and integral torsion.\allowbreak{}

\url{https://arxiv.org/pdf/math/0302267}.

\paragraph{S-AG-MCMULLEN} C.\allowbreak{}McMullen, Advanced Complex Analysis, Harvard Math213a course notes.

\textbf{Locator:} Version16Nov2026; Theorem3.\allowbreak{}16 printed92/\allowbreak{}PDF93; Theorems3.\allowbreak{}10–\allowbreak{}3.\allowbreak{}12 zero exponent.

\textbf{Inspection:} Primary author notes used for exact classical Hadamard locator; statement and applicability read.\allowbreak{} The theorem is classical, not presented as a2026 innovation.\allowbreak{}

\url{https://people.math.harvard.edu/~ctm/home/text/class/harvard/101/22/html/www/math213a/home/course/course.pdf}.

\paragraph{S-AG-KUMMER} Stacks Project, Kummer theory.

\textbf{Locator:} Lemma59.\allowbreak{}28.\allowbreak{}1 and immediately following long exact cohomology sequence; fppf variant Lemma59.\allowbreak{}28.\allowbreak{}3 kept separate.

\textbf{Inspection:} Entire proof of Lemma59.\allowbreak{}28.\allowbreak{}1 and displayed sequence read; ell-invertibility and inverse-limit transition conventions checked.\allowbreak{}

\url{https://stacks.math.columbia.edu/tag/03PK}.

\paragraph{S-AG-MILNE-ANT} J.\allowbreak{}S.\allowbreak{}Milne, Algebraic Number Theory, v3.\allowbreak{}08,19July2020.

\textbf{Locator:} Theorem4.\allowbreak{}3 PDF72 (Minkowski bound); Theorem4.\allowbreak{}4 class group finiteness.

\textbf{Inspection:} Exact bound and hypotheses read; explicit -35 ideal-class example proved in this run, not copied from source.\allowbreak{}

\url{https://www.jmilne.org/math/CourseNotes/ANT.pdf}.

\paragraph{S-AG-149} Alexander Grothendieck, cote149, Autour de la Longue Marche à travers la théorie de Galois.

\textbf{Locator:} 136PDFpages including cover.\allowbreak{} PDF80/\allowbreak{}footer79 eqs13–\allowbreak{}17; PDF81/\allowbreak{}footer80 eqs18–\allowbreak{}19; PDF82/\allowbreak{}footer81 eqs20–\allowbreak{}21; PDF85/\allowbreak{}footer84 eq22.\allowbreak{} Images PDF79–\allowbreak{}86 inspected.\allowbreak{}.

\textbf{Inspection:} Original official PDF downloaded and selected original page images read.\allowbreak{} The content is normalizer/\allowbreak{}inertia extensions, normal line bundles and Kummer cohomology.\allowbreak{} ORDER31 association to LM6 A44–\allowbreak{}A45 is separate inherited evidence; no K8 theorem attributed to sheets.\allowbreak{}

\url{https://grothendieck.umontpellier.fr/149.pdf}.

\paragraph{S-AG-MASULLO} Supplied Masullo The Music of the Primes draft.

\textbf{Locator:} Original PDF425 Lemma9.\allowbreak{}3.\allowbreak{}2.\allowbreak{}4;426 Theorem9.\allowbreak{}3.\allowbreak{}2.\allowbreak{}6 eq9.\allowbreak{}3.\allowbreak{}2.\allowbreak{}10;804 DefinitionD.\allowbreak{}3.\allowbreak{}3.\allowbreak{}1;11 norm/\allowbreak{}SE definition.

\textbf{Inspection:} Original images425,426,804 inspected, norm text11 read.\allowbreak{} Independent whole-proof checks MT-M01–\allowbreak{}03 recorded in events.\allowbreak{}tsv.\allowbreak{} No claim that M itself is disproved by auxiliary counterexamples.\allowbreak{}


\paragraph{S-AG-HODGE} Supplied Hodge bridge sketch.

\textbf{Locator:} OURS\_\allowbreak{}HODGE\_\allowbreak{}BRIDGE\_\allowbreak{}SKETCH\_\allowbreak{}20261121.\allowbreak{}md §§1–\allowbreak{}5, archive map§5; cote16 scans14–\allowbreak{}18,44–\allowbreak{}50; cote44 scans2,4,6,8–\allowbreak{}9; cote47 scans3,5–\allowbreak{}9.

\textbf{Inspection:} Relevant arguments and reported archive locators read.\allowbreak{} Offset in cote16 scan convention retained.\allowbreak{} Separate2026CM-Hodge claims not independently certified from this sketch.\allowbreak{}


\paragraph{S-AG-SA} Supplied SA22 and SA23 intakes.

\textbf{Locator:} OURS\_\allowbreak{}INTAKE\_\allowbreak{}RESEARCH\_\allowbreak{}SA22\_\allowbreak{}20261120.\allowbreak{}md §5.\allowbreak{}1 A6-node-character/\allowbreak{}A:\allowbreak{}inertia-maximality,§5.\allowbreak{}3Soulé/\allowbreak{}GT,§5.\allowbreak{}4–\allowbreak{}5.\allowbreak{}8; SA23§4.\allowbreak{}1 geo-H-lidskii/\allowbreak{}geo-H-genus.

\textbf{Inspection:} Relevant statements and intake check records read; original proofs/\allowbreak{}scans absent except149 separately retrieved.\allowbreak{} Cote1 scans79–\allowbreak{}80,111–\allowbreak{}112 are inherited locators.\allowbreak{}


\paragraph{S-AG-ANDERSON} G.\allowbreak{}W.\allowbreak{}Anderson, Cyclotomy and an extension of the Taniyama group, Compositio Mathematica57(1986),153–\allowbreak{}217.

\textbf{Locator:} §5.\allowbreak{}3 printed178/\allowbreak{}PDF27; Theorem6 §5.\allowbreak{}6 printed179–\allowbreak{}180/\allowbreak{}PDF28–\allowbreak{}29; Cor5.\allowbreak{}7.\allowbreak{}2/\allowbreak{}Rmk5.\allowbreak{}7.\allowbreak{}3/\allowbreak{}Cor5.\allowbreak{}7.\allowbreak{}4 printed181/\allowbreak{}PDF30; Prop5.\allowbreak{}7.\allowbreak{}5 printed181–\allowbreak{}182/\allowbreak{}PDF30–\allowbreak{}31; §5.\allowbreak{}8 Theorem7 printed182/\allowbreak{}PDF31.

\textbf{Inspection:} Original primary PDF downloaded; images27–\allowbreak{}31 and complete proofs Theorem6, Prop5.\allowbreak{}7.\allowbreak{}5 and Theorem7 read.\allowbreak{} Deeper cited Serre/\allowbreak{}Langlands/\allowbreak{}Deligne/\allowbreak{}Blasius results are classical inputs, not re-proved.\allowbreak{} Printed Rmk5.\allowbreak{}7.\allowbreak{}3 refers to Cor5.\allowbreak{}8.\allowbreak{}2, an internal numbering error noted explicitly.\allowbreak{}

\url{https://www.numdam.org/item/CM_1986__57_2_153_0.pdf}.

\paragraph{S-AG-FRAMEWORK} Classical geometric framework locators carried in ZETA\_\allowbreak{}TWO.

\textbf{Locator:} SGA4 XVII§§5.\allowbreak{}1–\allowbreak{}5.\allowbreak{}3; SGA5 XV§3 no2 unnumbered theorem eq1; EGA IV3(8.\allowbreak{}8.\allowbreak{}2); WeilI Thm1.\allowbreak{}6; WeilII Thm3.\allowbreak{}3.\allowbreak{}1; SGA6 V§§2,7(7.\allowbreak{}1.\allowbreak{}1); Deligne1979§§1.\allowbreak{}1–\allowbreak{}1.\allowbreak{}3,(3.\allowbreak{}1.\allowbreak{}2).

\textbf{Inspection:} Carried exact locators, not a claimed new full reading of the original SGA/\allowbreak{}EGA/\allowbreak{}Deligne works.\allowbreak{}


\paragraph{S-AG-34} Alexander Grothendieck, cote34, ExposéII.

\textbf{Locator:} 210PDFpages includingcover; Thm1.\allowbreak{}1 PDF53/\allowbreak{}footer52; struckRem1.\allowbreak{}17 PDF64/\allowbreak{}footer63; Thm1.\allowbreak{}1 proof §§2–\allowbreak{}3 PDF70–\allowbreak{}78/\allowbreak{}footer69–\allowbreak{}77; (2.\allowbreak{}5.\allowbreak{}5) PDF74/\allowbreak{}footer73.

\textbf{Inspection:} Official84,513,436bytePDF downloaded; pageimages53–\allowbreak{}78 read, with enlarged crops53/\allowbreak{}64.\allowbreak{} Exact normal/\allowbreak{}locallyfinite/\allowbreak{}geoconnected hypotheses and pro-ell solvable quotient wording recovered.\allowbreak{} Full theorem1.\allowbreak{}1 proof route §§2–\allowbreak{}3 read and geometric reductions checked; foundational SGA citations not separately re-proved.\allowbreak{} Webopen failed because filelarge; original local download succeeded.\allowbreak{}

\url{https://grothendieck.umontpellier.fr/34.pdf}.

\paragraph{S-AG-RV} D.\allowbreak{}Ramakrishnan and R.\allowbreak{}J.\allowbreak{}Valenza, Fourier Analysis on Number Fields(1999).

\textbf{Locator:} §7.\allowbreak{}4 Theorem7-19 printed279; §7.\allowbreak{}7 Theorem7-28 printed289, whole proof289–\allowbreak{}292.

\textbf{Inspection:} Primary PDF downloaded by number\_\allowbreak{}theory; theorem/\allowbreak{}hypothesis text and complete boundary proof289–\allowbreak{}292 directly read here.\allowbreak{} Boundary includes all unitary characters and norm twists; completed norm-character pole exceptions retained.\allowbreak{} Full proof of Theorem7-19 not separately re-refereed.\allowbreak{}

\url{https://webusers.imj-prg.fr/~julien.marche/Ref_GL_GT/Fourier_Analysis_on_Number_Fields.pdf}.

\paragraph{S-AG-MILNE-JAV} J.\allowbreak{}S.\allowbreak{}Milne, Jacobian Varieties and Abelian Varieties, revised notes of1986 chapters.

\textbf{Locator:} JacobianVarieties version12June2021 Prop9.\allowbreak{}1 pp30–\allowbreak{}32; AbelianVarieties version2Jan2022 Remarks15.\allowbreak{}3/\allowbreak{}15.\allowbreak{}5 pp28–\allowbreak{}29.

\textbf{Inspection:} Primary statements and comparison proof read; the Galois-compatible Tate-module/\allowbreak{}abelianized-fundamental-group identification is used, not an asserted cycle comparison.\allowbreak{}

\url{https://www.jmilne.org/math/xnotes/JVs.pdf}.

\paragraph{CX-S-KOIRAN} Pascal Koiran, Hilbert’s Nullstellensatz is in the Polynomial Hierarchy, DIMACS 96-27.

\textbf{Locator:} Section 2 Theorem 2 pp.\allowbreak{}5–\allowbreak{}6; Section 5 Theorem 9 and Corollary 1 printed p.\allowbreak{}12 (PDF13); 8 August 1997 erratum p.\allowbreak{}1.

\textbf{Inspection:} Downloaded both author PS files, converted to PDF, read theorem statements, protocol proof, number-field input, entire erratum.\allowbreak{}

\url{https://perso.ens-lyon.fr/pascal.koiran/Publis/DIMACS.96-27.ps}.

\paragraph{CX-S-BACH} Eric Bach, Explicit bounds for primality testing and related problems, Math.\allowbreak{} Comp.\allowbreak{} 55 (1990), 355–\allowbreak{}380.

\textbf{Locator:} Section 6 Theorems 2–\allowbreak{}3 p.\allowbreak{}372, Miller consequence p.\allowbreak{}373; Section 7 Theorem 4 p.\allowbreak{}375, corollary p.\allowbreak{}376.

\textbf{Inspection:} Original AMS PDF returned 403.\allowbreak{} Exact locators and statements carried from CX\_\allowbreak{}GRH/\allowbreak{}ALPOGE\_\allowbreak{}GCT; not represented as freshly checked original proof.\allowbreak{}

\url{https://doi.org/10.1090/S0025-5718-1990-1023756-8}.

\paragraph{CX-S-SCHOOF} René Schoof, Four primality testing algorithms, MSRI Publications 44 (2008).

\textbf{Locator:} Theorem 2.\allowbreak{}2 pp.\allowbreak{}105–\allowbreak{}106, including full proof specifying quadratic Dirichlet characters.

\textbf{Inspection:} Read statement and full proof via primary PDF text.\allowbreak{} PDF text loses minus signs; transfer uses standard Miller test and the theorem’s stated cutoff.\allowbreak{}

\url{https://pub.math.leidenuniv.nl/~stevenhagenp/ANTproc/05rene.pdf}.

\paragraph{CX-S-AKS} Agrawal, Kayal and Saxena, PRIMES is in P, Ann.\allowbreak{} Math.\allowbreak{} 160 (2004), 781–\allowbreak{}793.

\textbf{Locator:} Theorems 4.\allowbreak{}1 and 5.\allowbreak{}1; 2019 erratum Ann.\allowbreak{} Math.\allowbreak{} 189, 317–\allowbreak{}318.

\textbf{Inspection:} Primary article and erratum publication located; used only for unconditional P classification.\allowbreak{}

\url{https://annals.math.princeton.edu/wp-content/uploads/annals-v160-n2-p12.pdf}.

\paragraph{CX-S-SHOUP} Victor Shoup, Searching for Primitive Roots in Finite Fields, Math.\allowbreak{} Comp.\allowbreak{} 58 (1992), 369–\allowbreak{}380.

\textbf{Locator:} Author version Theorem 3 p.\allowbreak{}4, Section 6 p.\allowbreak{}13; Proposition 9 p.\allowbreak{}10.

\textbf{Inspection:} Downloaded author PDF and visually inspected p.\allowbreak{}4:\allowbreak{} Theorem3 is O(r\textasciicircum{}4(log r+1)\textasciicircum{}4(log p)\textasciicircum{}2), r=omega(p\ensuremath{-}1).\allowbreak{} Text encoding is corrupted; image is authoritative.\allowbreak{} No full original proof check claimed.\allowbreak{}

\url{https://shoup.net/papers/primroots.pdf}.

\paragraph{CX-S-HAC} Menezes, van Oorschot and Vanstone, Handbook of Applied Cryptography (1997).

\textbf{Locator:} Algorithm 3.\allowbreak{}34 and Fact 3.\allowbreak{}35 p.\allowbreak{}100; Fact 3.\allowbreak{}46 (composite square-root/\allowbreak{}factoring relation).

\textbf{Inspection:} Primary author-hosted chapter, exact square-root algorithm and its random step checked.\allowbreak{}

\url{https://cacr.uwaterloo.ca/hac/about/chap3.pdf}.

\paragraph{CX-S-GCTV} Ketan Mulmuley, Geometric Complexity Theory V:\allowbreak{} Efficient Algorithms for Noether Normalization, arXiv:\allowbreak{}1209.\allowbreak{}5993v5.

\textbf{Locator:} Theorem 2.\allowbreak{}10 p.\allowbreak{}25; Theorem 3.\allowbreak{}4 pp.\allowbreak{}27–\allowbreak{}28; Theorem 4.\allowbreak{}1 p.\allowbreak{}30; Theorem 4.\allowbreak{}7/\allowbreak{}Lemma 4.\allowbreak{}8 p.\allowbreak{}33; Theorem 4.\allowbreak{}10 p.\allowbreak{}34.

\textbf{Inspection:} Primary PDF downloaded and relevant statements/\allowbreak{}proof of 3.\allowbreak{}4 read; output and equation-access qualifications checked.\allowbreak{}

\url{https://arxiv.org/pdf/1209.5993}.

\paragraph{CX-S-FS} Michael Forbes and Amir Shpilka, A PSPACE Construction of a Hitting Set for the Closure of Small Algebraic Circuits, ECCC TR17-163.

\textbf{Locator:} Theorem 7.\allowbreak{}1 p.\allowbreak{}21; introduction Theorems 1.\allowbreak{}2–\allowbreak{}1.\allowbreak{}3.

\textbf{Inspection:} Primary theorem and unary-encoding condition checked.\allowbreak{}

\url{https://eccc.weizmann.ac.il/report/2017/163/download/}.

\paragraph{CX-S-GSS} Guo, Saxena and Sinhababu, Algebraic Dependencies and PSPACE Algorithms in Approximative Complexity over Any Field, Theory of Computing 15(16) (2019).

\textbf{Locator:} Theorem1.\allowbreak{}3 p5; Lemma4.\allowbreak{}1/\allowbreak{}Theorem4.\allowbreak{}2 pp15–\allowbreak{}16; Theorems1.\allowbreak{}5–\allowbreak{}1.\allowbreak{}6 p6.

\textbf{Inspection:} Primary journal PDF freshly retrieved via theoryofcomputing.\allowbreak{}org; checked APS statement/\allowbreak{}closure characterization and precise separating-e.\allowbreak{}s.\allowbreak{}o.\allowbreak{}p.\allowbreak{} statements in Theorems1.\allowbreak{}5–\allowbreak{}1.\allowbreak{}6.\allowbreak{}

\url{https://theoryofcomputing.org/articles/v015a016/v015a016.pdf}.

\paragraph{CX-S-IKRS} Ivanyos, Karpinski, Rónyai and Saxena, Trading GRH for algebra:\allowbreak{} algorithms for factoring polynomials and related structures.

\textbf{Locator:} 2008 author version Main Theorem p.\allowbreak{}3, Lemma4.\allowbreak{}2 p.\allowbreak{}14, Proposition4.\allowbreak{}6 pp.\allowbreak{}16–\allowbreak{}17; published version Theorems1.\allowbreak{}1,1.\allowbreak{}3,1.\allowbreak{}7.

\textbf{Inspection:} Fresh published primary source read for factoring-vs-automorphism and quasipolynomial scope; exact older-version component locators retained from CX\_\allowbreak{}GRH.\allowbreak{}

\url{https://www.cse.iitk.ac.in/users/nitin/papers/nogrh.pdf}.

\paragraph{CX-S-BS} Eric Bach and Jonathan Sorenson, Explicit Bounds for Primes in Residue Classes, Math.\allowbreak{} Comp.\allowbreak{}65 (1996),1717–\allowbreak{}1735.

\textbf{Locator:} Section2 p.\allowbreak{}1720; Theorem5.\allowbreak{}1 p.\allowbreak{}1729.

\textbf{Inspection:} Primary author-uploaded PDF retrieved; statement and precise analytic premise read; geometry agent independently verified same theorem.\allowbreak{}

\url{https://www.researchgate.net/profile/Jonathan-Sorenson-2/publication/220576951_Explicit_bounds_for_primes_in_residue_classes/links/55100ea70cf2ac2905afbceb/Explicit-bounds-for-primes-in-residue-classes.pdf}.

\paragraph{CX-S-MILLER} Gary L.\allowbreak{} Miller, Riemann’s Hypothesis and Tests for Primality, JCSS13(3) (1976),300–\allowbreak{}317.

\textbf{Locator:} Lemma5 and Theorem3 p310; Theorem4 pp311–\allowbreak{}312.

\textbf{Inspection:} Author-hosted primary PDF downloaded; p310 visually inspected and entire Lemma5 proof read.\allowbreak{} \ensuremath{\lambda}\ensuremath{{}^{\prime}} multiple and polynomial bit-growth conditions confirmed; explicit-input version proved here.\allowbreak{}

\url{https://www.cs.cmu.edu/~glmiller/Publications/Papers/Mi76.pdf}.

\paragraph{CX-S-BIMS} Bhargava, Ivanyos, Mittal and Saxena, Irreducibility and deterministic r-th root finding over finite fields, ISSAC2017.

\textbf{Locator:} Corollary1.\allowbreak{}3 authorPDF p2; Theorem3.\allowbreak{}1 and full proof p3; Corollary3.\allowbreak{}2 p4.

\textbf{Inspection:} Statement and complete Theorem3.\allowbreak{}1 proof read; original p3 rendered and inspected.\allowbreak{} Literal nonvanishing step fails for the explicit F7 cubic example CX-ALG-016.\allowbreak{} Corrected basis-search proof CX-ALG-017; even-degree quadratic case independently re-proved CX-ALG-018.\allowbreak{} Corollary conclusions not refuted.\allowbreak{}

\url{https://www.cse.iitk.ac.in/users/nitin/papers/irred-and-rth-root.pdf}.

\paragraph{CX-S-BE} Jean-François Biasse and Muhammed Rashad Erukulangara, A proof of the conjectured run time of the Hafner–\allowbreak{}McCurley class group algorithm, Advances in Mathematics of Communications17(6) (2023),1388–\allowbreak{}1405.

\textbf{Locator:} Author PDF MainTheorem and Conjecture1 p4; Algorithm4 p15; Proposition6 pp15–\allowbreak{}16; Corollary1 p16.

\textbf{Inspection:} Direct PDF request returned503 but web source recovered through indexed reference.\allowbreak{} Main theorem, explicit Hecke-GRH definition, complete Algorithm4, Proposition6 and Corollary1 with its proof read.\allowbreak{} Entire 18-page proof not claimed re-audited; prior CX-N12 detailed transfer retained.\allowbreak{}

\url{https://par.nsf.gov/servlets/purl/10324557}.

\paragraph{CX-S-JMV} David Jao, Stephen D.\allowbreak{} Miller and Ramarathnam Venkatesan, Expander graphs based on GRH with an application to elliptic curve cryptography, J.\allowbreak{} Number Theory129 (2009),1491–\allowbreak{}1504.

\textbf{Locator:} Author arXiv0811.\allowbreak{}0647v2:\allowbreak{} Theorem1.\allowbreak{}1 p3; Remark1.\allowbreak{}2(a)/\allowbreak{}Corollary1.\allowbreak{}3 p4; Theorem1.\allowbreak{}5 p5; Theorem1.\allowbreak{}6 pp6–\allowbreak{}7; Lemma2.\allowbreak{}1 pp7–\allowbreak{}8; proof of Theorem1.\allowbreak{}1 p9 and Cor1.\allowbreak{}3 p10; Proposition3.\allowbreak{}1 p11, Theorem3.\allowbreak{}3 p12.

\textbf{Inspection:} Primary author-hosted PDF read for all listed statements.\allowbreak{} Full proofs of Theorem1.\allowbreak{}1, Cor1.\allowbreak{}3 and Lemma2.\allowbreak{}1 read; verified estimates are character-by-character, enabling the genus quotient refinement.\allowbreak{} Isogeny statements and supplied-solver scope checked; no full underlying CM/\allowbreak{}Deuring proof audit claimed.\allowbreak{}

\url{https://web.williams.edu/Mathematics/sjmiller/public_html/crypto/handouts/JaoMillerVenkatesan_ExpanderGraphsGRHEllCurveCrypto.pdf}.

\paragraph{A-SOURCE} Levent Alpoge, Remark.

\textbf{Locator:} PDF pp3–\allowbreak{}4 steps2–\allowbreak{}20; ALPOGE\_\allowbreak{}Main.\allowbreak{}lean five declarations; ALPOGE\_\allowbreak{}Statements.\allowbreak{}lean definitions.

\textbf{Inspection:} User-supplied PDF text pp1–\allowbreak{}4 read; declarations and definitions read; no independent Lean kernel run.\allowbreak{} Cited theorem input A\_\allowbreak{}D as instructed.\allowbreak{}

\url{https://alpo.ge/remark.pdf}.

\paragraph{M-SOURCE} A.\allowbreak{}M.\allowbreak{} Masullo, The Music of the Primes, supplied 940-page draft.

\textbf{Locator:} §1.\allowbreak{}2 pp12–\allowbreak{}19; Lemma9.\allowbreak{}3.\allowbreak{}2.\allowbreak{}4 p425; Theorem9.\allowbreak{}3.\allowbreak{}2.\allowbreak{}6 p426; DefD.\allowbreak{}3.\allowbreak{}3.\allowbreak{}1 p804; Theorem12.\allowbreak{}5.\allowbreak{}1.\allowbreak{}1 pp519–\allowbreak{}520;13.\allowbreak{}3.\allowbreak{}1.\allowbreak{}1 and placeholder13.\allowbreak{}3.\allowbreak{}1.\allowbreak{}2 p522.

\textbf{Inspection:} Original PDF425,426 images inspected; source definitions11,804 read.\allowbreak{} Claims M remain conditional, independent intermediate assertions not inferred from endpoint.\allowbreak{}


\paragraph{KW09} Chandrashekhar Khare and Jean-Pierre Wintenberger, Serre's modularity conjecture (I).

\textbf{Locator:} Theorem10.\allowbreak{}1(ii), author PDF20; Cor10.\allowbreak{}2(ii), author PDF21.

\textbf{Inspection:} Primary source theorem and Corollary read, including A5 scope and Artin-L-function discussion.\allowbreak{}

\url{https://www.math.ucla.edu/~shekhar/papers/results.pdf}.

\paragraph{STARK74} Harold M.\allowbreak{} Stark, Some effective cases of the Brauer-Siegel theorem, Invent.\allowbreak{}Math.\allowbreak{}23(1974),135–\allowbreak{}152.

\textbf{Locator:} Theorem3,p144.

\textbf{Inspection:} Exact source locator corroborated by primary later paper Gun-Hazra-Sahu2025 introduction; original volume web fetch timed out.\allowbreak{} Classical theorem statement used; no claim to have read original proof.\allowbreak{}

\url{https://gdz.sub.uni-goettingen.de/download/pdf/PPN356556735_0023/PPN356556735_0023.pdf}.

\paragraph{GHS25} Sanoli Gun, Suhita Hazra and Dhananjaya Sahu, On holomorphy and non-vanishing of Artin L-functions, Monatshefte fur Mathematik206(2025),853–\allowbreak{}883.

\textbf{Locator:} Theorem5 and Corollary6; §6; introduction for Stark statement.

\textbf{Inspection:} Primary statements and proof mechanism of integer-root Euler deformation checked.\allowbreak{} MT-L01/\allowbreak{}L02 gives own proof for general |c|\ensuremath{\le}1 and irrational c.\allowbreak{}

\url{https://link.springer.com/article/10.1007/s00605-025-02060-7}.

\paragraph{GM19-A1} Loic Grenie and Giuseppe Molteni, An explicit Chebotarev density theorem under GRH.

\textbf{Locator:} AppendixA TheoremA.\allowbreak{}1, author PDF26.

\textbf{Inspection:} Number\_\allowbreak{}theory worker verified primary finite-order Hecke zero count; root uses the result as a cited classical input and subtracts consecutive-height counts.\allowbreak{}

\url{https://sites.unimi.it/molteni/research/papers-pdf/44-molteni-An_explicit_Chebotarev_density_theorem_under_GRH.pdf}.

\end{document}
